\pdfinfoomitdate=1
\pdftrailerid{}
\pdfsuppressptexinfo=15
\documentclass[11pt,letterpaper]{article}
\usepackage[T1]{fontenc}
\usepackage{lmodern,microtype,amsmath,amssymb,amsthm,mathtools,booktabs,array,longtable}
\usepackage[margin=1in]{geometry}
\usepackage{xcolor}
\usepackage[colorlinks=true,linkcolor=blue!45!black,urlcolor=blue!45!black,citecolor=blue!45!black]{hyperref}
\usepackage{fancyhdr}
\usepackage{enumitem}
\hypersetup{pdftitle={Extreme-part conditions in partitions (OEIS A350879, A237753, A237751, A118096, A117086)},pdfauthor={Research report},pdfsubject={Merged report a350879-extreme-part-conditions (batch 111): fixed-slope extremes (Report 216), ratio two (Report 217), divisible extremes (Report 218)}}
\pagestyle{fancy}\fancyhf{}\fancyhead[L]{\small Extreme-part conditions}\fancyhead[R]{\small Reports 216, 217, 218}\fancyfoot[C]{\thepage}
\setlength{\headheight}{14pt}
\newtheorem{theorem}{Theorem}[section]
\newtheorem{lemma}[theorem]{Lemma}\newtheorem{proposition}[theorem]{Proposition}\newtheorem{corollary}[theorem]{Corollary}
\theoremstyle{remark}\newtheorem{remark}[theorem]{Remark}
% Report 216 (Part I)
\newcommand{\ind}{\mathbf 1}\newcommand{\ee}{\mathrm e}\newcommand{\ceil}[1]{\left\lceil#1\right\rceil}
\newcommand{\Ekb}{E_{k,b}}\newcommand{\Tkb}{T_{k,b}}
% Report 217 (Part II); its \ceil is Part I's
\newcommand{\Li}{\operatorname{Li}_2}
\newcommand{\E}{\mathbb E}
\newcommand{\Q}{\mathbb Q}
\newcommand{\R}{\mathbb R}
\newcommand{\coeff}[2]{[#1^{#2}]}
\newcommand{\ii}{\mathrm i}
% Report 218 (Part III); its \ceil and \coeff are the above
\newcommand{\A}{\mathsf A}
\newcommand{\B}{\mathcal B}
\newcommand{\dd}{\mathrm d}
\newcommand{\e}{\mathrm e}
\newcommand{\ord}{\operatorname{ord}}
\newcommand{\floor}[1]{\left\lfloor #1\right\rfloor}
% Added in the merge (batch 111, 7 October 2026).
\newcommand{\file}[1]{\nolinkurl{#1}}
\expandafter\def\expandafter\UrlBreaks\expandafter{\UrlBreaks\do\-\do\_\do\.\do\:}
\newcommand{\vol}{\mathrm{vol}}
\newenvironment{writenote}[1][]{\par\smallskip\begingroup\small
 \noindent\textbf{[write]}\if\relax\detokenize{#1}\relax\else~\textbf{#1.}\fi\ }%
 {\par\endgroup\smallskip}
\newenvironment{partabstract}{\begin{quote}\small\noindent\textbf{Abstract of this Part (as delivered).}\ }{\end{quote}}
\newcommand{\partsource}[1]{\par\noindent\begin{minipage}{\textwidth}\small #1\end{minipage}\par\medskip}
\newlength{\xpcdefparindent}\setlength{\xpcdefparindent}{\parindent}
\newcounter{xpcsavesec}
\title{\textbf{Extreme-part conditions in integer partitions}\\[4pt]\large (OEIS A350879, A237753, A237751, A118096 and A117086)\\[8pt]
\normalsize Part I: fixed slope partition extremes and inverse asymptotics\\[2pt]
Part II: exact ratio two partitions at fixed asymptotic order\\[2pt]
Part III: deficit asymptotics for partitions with divisible extreme parts}
\author{Three research reports of one external research session\\(bundle Reports 216, 217 and 218), bound into one ProveIt report}
\date{All Parts: 4 October 2026\quad merged: 7 October 2026}
\begin{document}
\hypersetup{pageanchor=false}
\maketitle\thispagestyle{fancy}
\begin{abstract}
Three manuscripts on partitions $\lambda$ of $n$ constrained by a relation between their extreme parts, bound into one report. Part~I fixes integers $k\ge2$, $b\ge0$ and counts partitions with largest part $M$ and length $L$ satisfying $M=kL+b$ or $M\ge kL+b$ (by conjugation, the columns of A350879; A237753 and A237751 at $k=2$): exact finite-product sieves, expansions to every fixed order with leading coefficient $k!$, explicit corrections, and a Lambert $W_{-1}$ inverse with a shrinking two-ceiling bracket. Part~II counts partitions whose largest part is exactly twice the smallest (A118096): a complete-circle estimate, a discrete two-variable saddle and an exact finite coefficient rule give every fixed order on the scale $e^{\pi\sqrt{2n/15}}$, four explicit corrections, rounding-aware inverses, and eventual strict increase and log-concavity. Part~III counts partitions whose largest part is divisible by the smallest (A117086) and proves that the deficit $D(n)=p(n)-a(n)$ satisfies $D(n)\sim\frac{\pi}{2\sqrt{6n}}p(n)$ with an expansion to every fixed order, fixed-minimum residue equidistribution, a stretched-exponential comparison with a reciprocal-minimum weighted count, a rational certificate at $n=1000$, and inverse brackets. Parts~I and~III share the finite-shift Rademacher transfer; Part~II's generating function is the $L=2$ summand of Part~III's. The write records that A350879's alternative elementary form, with $\exp(\sqrt{2\pi n/3})$, is false (the rate is $\exp(\pi\sqrt{2n/3})$), proves that A117086 increases strictly for every $n$, quotes the five OEIS entries, and compares the inverses with the repository's transseries volume. The report is unrefereed, and nothing in it is formalized.
\end{abstract}
\hypersetup{pageanchor=true}
{\small\setcounter{tocdepth}{1}\setlength{\parskip}{0pt}\tableofcontents}
\newpage

\section*{Guide to this report}\phantomsection\label{xpc:sec:guide}
\addcontentsline{toc}{section}{Guide to this report}
The three Parts are three manuscripts of the session bundle of Reports
1--243 (arrival commit \file{60f54ea06}), placed in this directory by commit
\file{d451ef3d8} (batch~111).
\begin{center}
\small
\begin{tabular}{@{}l l >{\raggedright\arraybackslash}p{0.5\textwidth} l@{}}
\toprule
Part & Bundle report & Delivered title & Labels\\
\midrule
I & 216 (base) & Fixed slope partition extremes and inverse asymptotics (4~October 2026) & \file{xpc:sl:}\\
II & 217 & Exact ratio two partitions at fixed asymptotic order (4~October 2026; source edition ``v2'') & \file{xpc:r2:}\\
III & 218 & Deficit asymptotics for partitions with divisible extreme parts (4~October 2026) & \file{xpc:dv:}\\
\bottomrule
\end{tabular}
\end{center}

\emph{What binds the Parts.} All three count partitions through a condition
coupling the extreme parts: $M=kL+b$ or $M\ge kL+b$ with $M$ the largest
part and $L$ the length (Part~I; by conjugation, $M$ and $L$ trade places, so
these are also conditions on the smallest-part multiplicity), $\max=2\min$
(Part~II) and $\min\mid\max$ (Part~III). Part~II's generating function
$\sum_kq^{3k}/\prod_{j=k}^{2k}(1-q^j)$ is the $L=2$ summand
of~\eqref{xpc:dv:eq:exact-a}: the partitions of Part~II are exactly those of
A117086 whose ratio $\max/\min$ equals~$2$. Parts~I and~III share their
analytic machinery: a finite sieve or polynomial reduction whose discarded
part is bounded on the entire Cauchy circle, and a finite-shift transfer from
the Rademacher series with the same terminating polynomials $Q_\ell$. Part~II
uses a different method (a discrete two-variable saddle at the golden-ratio
point and Bessel identification) and a different scale.
\begin{center}
\small
\begin{tabular}{@{}>{\raggedright\arraybackslash}p{0.14\textwidth}>{\raggedright\arraybackslash}p{0.26\textwidth}>{\raggedright\arraybackslash}p{0.26\textwidth}>{\raggedright\arraybackslash}p{0.26\textwidth}@{}}
\toprule
& Part I (216) & Part II (217) & Part III (218)\\
\midrule
Condition & $M=kL+b$, $M\ge kL+b$ ($k\ge2$, $b\ge0$ fixed) & $\max=2\min$ & $\min\mid\max$ (the deficit $D=p-a$)\\
OEIS & A350879 columns; A237753, A237751 & A118096 & A117086\\
Scale & $k!\,B(n)t^{k}$ and $k!\,B(n)t^{k-1}$, $t\sim\pi/\sqrt{6n}$ & $Cn^{-1/2}e^{\pi\sqrt{2n/15}}$ & $B(n)t/2$\\
Orders & every fixed order; four coefficients for A237753, A237751 & every fixed order; $c_1,\dots,c_4$ & every fixed order; $c_0,\dots,c_5$\\
Method & finite sieve, circle bound, Rademacher shift transfer & complete-circle bound, discrete saddle, Bessel identification & polynomial reduction, first-harmonic gap, shift transfer\\
Inverse & $W_{-1}$ core, triangular reversion, two ceilings & $\log$-$\log$ polynomials $P_j$, two ceilings & $W_{-1}$ core, triangular reversion, two ceilings; exact two-index shortcut\\
Also & fixed-offset calculus & strict log-concavity & residue equidistribution; weighted count; rational certificate\\
\bottomrule
\end{tabular}
\end{center}

\emph{Arrangement and duplication.} The base of the merge is Report~216: its
\file{Report216.tex} was staged as this \file{article.tex}; it is the most
general (a two-parameter family), it arrived first of the three, and it
carries the batch's one recorded OEIS correction. Report~217 is Part~II
and Report~218 Part~III. No Part cites another. Part~III re-derives the
partition input that Part~I proves: the one-term Rademacher gap, the
terminating polynomials $Q_\ell$ and the finite-shift transfer
(Section~\ref{xpc:dv:sec:transfer} against Lemmas~\ref{xpc:sl:lem:gap}
and~\ref{xpc:sl:lem:Q} and Proposition~\ref{xpc:sl:prop:transfer}), and the
same triangular inverse reversion (Section~\ref{xpc:dv:sec:inverse} against
Proposition~\ref{xpc:sl:prop:inverse-coefs}). These are kept, with dated
notes naming Part~I's counterparts, because Part~III's later statements cite
them by their own labels; the proofs are the same argument.

\emph{Numbering.} Sections are numbered continuously, statements within
sections and equations globally, as delivered. Part~I keeps all its numbers,
including its Tables~1--2 and its Appendix~A. Report~217's Section~$k$ is
Section~$k+10$ here and its equation~($k$) is equation~($k+65$); Report~218's
Section~$k$ is Section~$k+21$, its equation~($k$) is equation~($k+103$), and
its appendix keeps the letter~A. Statements move with their sections. Part
labels carry \file{xpc:sl:}, \file{xpc:r2:} and \file{xpc:dv:}; the labels
added by the write carry the Part prefix, or \file{xpc:} for this Guide. The
bibliography is merged at the end: the OEIS entries of Parts~II and~III
(both delivered as \file{oeis}) are \file{oeisA118096} and
\file{oeisA117086}; Part~II's \file{dlmf} (Bessel sections) is
\file{dlmfBessel}; Parts~I and~III's \file{dlmf} (Section~23.18) and
\file{johansson} are printed once; Part~III's \file{nr} is Part~I's
\file{ngorhoades} (the same paper).

\begin{writenote}[Provenance, added 7 October 2026, batch~111]\raggedright
This report was built from three manuscripts, all dated 4~October 2026,
arrived in \file{60f54ea06} and placed in \file{d451ef3d8}. \emph{Part~I}:
Report~216, archive \file{Report216-reproducibility.zip} ($567{,}001$
bytes, SHA-256 \file{ca858cb05b9b...ea13edb9ec180c7}, 20 files in a wrapper
directory \file{Report216/}); text \file{Report216.tex} (738 lines, 20
pages as delivered). \emph{Part~II}: Report~217, archive
\file{Report217-source-v2.zip} ($409{,}022$ bytes, SHA-256
\file{96812623e011...abcc67943772177}, 21 files in \file{Report217/});
text \file{article.tex} (594 lines, 14 pages); the archive name carries the
edition mark ``v2'', the only such mark in the bundle; no other edition was
delivered. \emph{Part~III}: Report~218, archive
\file{Report218-reproducibility.zip} ($544{,}706$ bytes, SHA-256
\file{f688cd814dde...7e4377cbbb2b09f}, 17 files in \file{Report218/}); text
\file{report218.tex} (976 lines, 22 pages). Report~216's title page has an
empty author line and an empty PDF author field; Reports~217 and~218 read
``Report 217'' and ``Report 218''. None names a person, tool or addressee;
none carries a ``prepared for private review'' line, an e-mail address or
personal data, and none records a ProveIt commit. Each Part's programs are
shipped under \file{code/} and its data, generated tables and recorded
outputs under \file{data/}, with the prefixes \file{216-slope-},
\file{217-ratio2-} and \file{218-divisible-}; Report~216's \file{SOURCES.txt},
\file{SOURCE_FILES.txt} and code README, Report~217's
\file{SOURCE_NOTES.md} and code README, and Report~218's code README are
shipped at the top level under the same prefixes; Report~217's
\file{MANIFEST.json}, which also records its claim and source-comparison
fields, is shipped as \file{data/217-ratio2-MANIFEST.json}. Report~216's
delivered \file{README.txt} was staged as the report README and is replaced
by it. Not shipped, and retrievable from the arrival commit: the three
delivered PDFs, the texts of Reports~217 and~218 (they are Parts~II and~III
of this file), their delivered READMEs, and the pure checksum manifests
(Report~216 \file{MANIFEST.json}, 19 entries; Report~217
\file{generated/SHA256SUMS.json}, 5; Report~218 \file{MANIFEST.json} and
\file{SHA256SUMS}, 15 and 16), all verified at the write, as was
Report~217's \file{MANIFEST.json} (20 entries). Report~216's two generated
tables and Report~218's four generated tables, which the delivered texts
read with \verb|\input|, are printed inline here (the shipped files under
\file{data/} are byte-identical to the delivered ones). Where the merge had to
choose: the base and the order; the citation keys above; keeping Part~III's
re-derivations; the per-Part paragraph and list spacing of the deliveries
(Part~I sets no paragraph indent, Part~II tight lists, Part~III looser lists
and allowed display breaks), restored at each Part header. The write
prefixed the 92, 54 and 90 delivered labels (Report~216's 92 include the two
table labels of its generated tables) and updated the 83, 57 and 86
references to them (66, 36 and 60 \verb|\eqref|; 17, 21 and 26 \verb|\ref|);
it labelled three unlabelled sections of Report~217 and two of Report~218,
and added this Guide, the Part headers, the remarks marked [write]
(Remarks~\ref{xpc:sl:rem:oeis}, \ref{xpc:sl:rem:a350879},
\ref{xpc:sl:rem:transseries}, \ref{xpc:r2:rem:oeis},
\ref{xpc:r2:rem:transseries}, \ref{xpc:dv:rem:oeis},
\ref{xpc:dv:rem:transseries}) and the dated notes. The added remarks are the
last statements of their sections and the added displays are unnumbered, so
every delivered number is unchanged up to the stated shifts (label tables of
builds of the three delivered texts and of this one compared).

\emph{The sources, as the write read them} (7~October 2026). OEIS A350879,
A237753, A237751, A118096 and A117086 with their b-files, in the internal
format (Remarks~\ref{xpc:sl:rem:oeis}, \ref{xpc:r2:rem:oeis}
and~\ref{xpc:dv:rem:oeis}); the transseries volume
\path{Analysis/Transseries/docs/series-and-transseries/Transseries_And_Inversion/transseries_and_inversion.tex}
(\file{p0:def:model}, \file{p0:def:core}, \file{p0:thm:lambert-core},
\file{p0:thm:lambert-centered}, \file{p0:thm:flattening},
\file{p0:thm:backward-error}, \file{p0:def:three-inverses},
\file{p0:thm:staircase}). The write tried to open Pittel's 2007 paper at its
publisher (refused, HTTP~403); like the source, it did not read it. Not read
by the write: the rest of the cited literature; the Parts' accounts of what
they inspected are reported, not confirmed.

\emph{What was checked at intake.} The batch-111 dossier (7~October 2026)
read the three manuscripts, found no error in them, confirmed the A350879
discrepancy against the live entry, and reran the three delivered suites on
copies: their outputs reproduce, the only failures being Windows artifacts
(line-ending bytes in hashed outputs, a recorded Python version string, a
console code page), none mathematical. \emph{What the write checked.} Every
proof line by line, with no error found. With its own code (Python~3.14.4,
SymPy~1.14.0, mpmath~1.3.0) it recomputed, for Part~I, $E_{2,0}$ and
$T_{2,1}$ by the sieve to $n=3000$ (all A237753 and A237751 b-file terms in
range agree, as do brute-force counts for $n\le40$ and all columns $k\ge2$ of
A350879's 50-row b-file), the shipped selected counts, the radial
series~\eqref{xpc:sl:eq:radE}--\eqref{xpc:sl:eq:radT}, the coefficients
\eqref{xpc:sl:eq:Ecoefs} and~\eqref{xpc:sl:eq:Tcoefs}, the general first
coefficients~\eqref{xpc:sl:eq:c1E}--\eqref{xpc:sl:eq:c1T} for $2\le k\le6$,
$0\le b\le3$, the inverse coefficients~\eqref{xpc:sl:eq:u1}--\eqref{xpc:sl:eq:u3}
by solving~\eqref{xpc:sl:eq:reversion}, and~\eqref{xpc:sl:eq:indexcorr};
for Part~II, A118096 to $n=3000$ from~\eqref{xpc:r2:eq:F} (b-file and brute
force agree), the saddle data $H_2,H_3,H_4$, $\ell'(x_0)$, $\ell''(x_0)$,
$e_1$, $e_3$, $h(x_0)=A$, $\lambda$, the four coefficients~\eqref{xpc:r2:eq:cs}
from the finite rule~\eqref{xpc:r2:eq:crule} (see the note at the end of
Section~\ref{xpc:r2:sec:rule}), $d_1,\dots,d_4$ both from~\eqref{xpc:r2:eq:ds}
and from~\eqref{xpc:r2:eq:transfer}, $P_1$ and $P_2$, and the diagnostics of
Section~\ref{xpc:r2:sec:repro}; for Part~III, the counts against A117086's
b-file and brute force, $D(17)$, $D(18)$ and $D(1000)$, $g_1,\dots,g_8$ from
both~\eqref{xpc:dv:eq:G} and~\eqref{xpc:dv:eq:Gtail}, $c_0,\dots,c_5$,
$\rho_1,\rho_2,\rho_3$, $d^{(D)}_1,d^{(D)}_2,d^{(D)}_3$, the general
$v_1,v_2,v_3$, $v^{(D)}_{1,2}$, $v^{(a)}_{1,2}$, the
centres~\eqref{xpc:dv:eq:inverse-D}, \eqref{xpc:dv:eq:inverse-a}
and~\eqref{xpc:dv:eq:half-index}, $s_3(1000)$, $u_3(1000)$ and the two bounds
of~\eqref{xpc:dv:eq:worked-bound} from the printed $E$ (the certificate's
$E$ itself was not recomputed), and the rational inequalities of
Section~\ref{xpc:dv:sec:quantitative}. Finite and floating computations are
observations, not certificates.

\emph{Relation to the repository.} Before batch~111 no file of the
repository named A350879, A237753, A237751, A118096 or A117086. No statement
of this report is formalized in Lean or Rocq, and its place in the
collection confers no formal status. The nearest reports are in the same
directory
\file{SetTheory/Cardinals/docs/reports/generating-functions-and-asymptotics/oeis-sequence-asymptotics/}:
\file{a239950-maximal-schreier-supports} (batch~111: least part equal to the
number of distinct sizes, another extreme-part condition, on a different
exponential scale) and \file{a239964-sizes-equal-max-multiplicity}
(batch~111). They share no result with this report, so no reciprocal note
is proposed. The inverses are compared with the transseries volume in
Remarks~\ref{xpc:sl:rem:transseries}, \ref{xpc:r2:rem:transseries}
and~\ref{xpc:dv:rem:transseries}.
\end{writenote}
\begin{writenote}[What is not claimed, collected from the three Parts]
\emph{Part~I.} Fixed slope, fixed offset and fixed order only: $k$, $b$ and
$R$ may not grow with $n$, no convergence or numerical error bound for the
formal series, no effective onset, no $k=1$ (rank) conclusion, no
growing-offset sum; the inverse constants are not effective and the bracket
is not a certified finite-$y$ enclosure; no historical priority: the leading
equality law is Kot\v{e}\v{s}ovec's, partition-product transfer is not new,
the comparison with Szekeres, Jiang--Wang, Melczer--Panova--Pemantle,
Richmond, Hwang, Ngo--Rhoades, Dyson and Berkovich--Garvan is bounded; the
algebraic checks are not clean-room and do not form a separate proof;
diagnostics are not certified; same-toolchain byte stability only.
\emph{Part~II.} No novelty claim and no solution of a known open problem:
the leading equivalent is Kot\v{e}\v{s}ovec's (2025), the radial all-order
machinery is McIntosh's and Zhou's, and the comparison with Pittel (2007) is
unresolved (abstract only; ``not evidence of absence of overlapping
results''); no uniformity in the order, no explicit threshold, no effective
last exception for increase or log-concavity; the inversion and ceiling
mechanisms are standard; diagnostics are not certified; $H(q)$ is not
identified with a mock theta function. \emph{Part~III.} No historical
priority, no uniform growing-order expansion, no convergence; the twisted
product, all-order transfer and arbitrary-weight tail frameworks are prior
(Bridges--Franke--Garnowski, Ngo--Rhoades, Cesana--Craig--Males,
Andrews--Freitas, the last not read in the original); the stretched-exponential
bound is an upper bound, not an asymptotic, and proves no
$\exp(-c\sqrt n)$ rate, and its optimality concerns only the two-term
majorant; no global monotonicity of $D$; the inverse constants are
ineffective and no blind rounding is licensed; no quantum-modularity claim;
the A117086 revision history was not inspected. \emph{The write adds:} its
checks are floating, finite or symbolic computations; the three transseries
remarks claim no novelty for any inversion; the merge claims no result
connecting the Parts beyond the shared summand and machinery.
\end{writenote}
\medskip
\begingroup\small
\noindent\emph{Reading conventions} (letters the Parts use differently, with
the tempting false readings; no symbol of any source was renamed; symbols of
the transseries volume carry the subscript ``vol'' in the three transseries
remarks).\par\smallskip
\renewcommand{\arraystretch}{1.1}
\renewcommand{\theHtable}{xpcguide.\arabic{table}}% no anchor clash with Part I's Tables 1-2
\begin{longtable}{@{}>{\raggedright\arraybackslash}p{0.6in}>{\raggedright\arraybackslash}p{1.75in}>{\raggedright\arraybackslash}p{1.75in}>{\raggedright\arraybackslash}p{1.75in}@{}}
\toprule
Symbol & Part I & Part II & Part III\\
\midrule
\endhead
$A$ & $\pi^2/6$ & $\pi^2/30$ (and $a$ the curvature $3\phi-4$) & $\mathsf A=\pi^2/6$\\
$B$ & $B(n)$, the partition scale & $2\sqrt A$ & $\mathcal B(n)$, the partition scale\\
$C$ & $C_M$, an inverse constant & $1/\sqrt{5+\sqrt5}$ & $C_D$, $C_a$; $C_0$, $C_1$ explicit constants\\
$D$ & $D_j$, sieve exponents; the inverse constant $D$ & $D=2A-ax_0^2$ & $D(n)=p(n)-a(n)$, the deficit\\
$a(n)$, $X$ & $X\in\{E,T\}$ & $a(n)$ = A118096; $X_J(Y)$ the inverse & $a(n)$ = A117086; $X\in\{D,a\}$\\
$t$, $s$ & $t=\sqrt{A/(n-\frac1{24})}$; $s=s_X$ the valuation & $t=\sqrt{A/n}$; $s=\sqrt t$, but $s=B\sqrt n$ in Section~\ref{xpc:r2:sec:inverse} & $t=\sqrt{\mathsf A/N}$, $N=n-\frac1{24}$\\
$c_m$, $c_j$ & relative coefficients in $t$ & radial coefficients in $t$ (and $d_r$ in $n^{-1/2}$) & relative coefficients in $t$\\
$d$, $d_j$ & $d=s+2$ & $d_r$ & $d_j$, coefficients in $w^{-1}$\\
$G$, $g_\ell$ & $G^X_{k,b}$, $g_\ell$ & $g(x)$ the amplitude & $G(z)$, $g_\ell$; $g(x)$ in Section~\ref{xpc:dv:sec:weighted}\\
$H$ & $H(q)$ a sieve polynomial; $H(v)$ the forward series & $H_m=h^{(m)}(x_0)$; $H(q)$ the weak-class series; $H_J$ & $H_k(z,q)$; $H(q)$ the weighted series; $H$ an exponential degree\\
$Q$ & $Q_\ell(t)$ & $Q(y,v)$ the quadratic form & $Q_\ell(t)$ (the same as Part~I's)\\
$T$ & $T_{k,b}$, the tail count & $T_k(q)$ summands; $T=\log(Y/(CB))$ & $T_K(q)$, the minimum tail\\
$L$ & $L(\lambda)$, the length & $L$ a block length; $L_t$; $L_J(x)$ & $L$ in~\eqref{xpc:dv:eq:exact-a}; $L=\log(1/t)$ in Section~\ref{xpc:dv:sec:quantitative}\\
$w$, $w_0$ & $w=2\sqrt{A(n-\frac1{24})}$, $w_0$ the Lambert root & $w=\log T$ & $w=2\sqrt{\mathsf A(n-\frac1{24})}$, $w_0$\\
$\nu$, $N$ & $\nu_X(y)$ thresholds; $N=n-\frac1{24}$ & $N(Y)$ the threshold & $\nu_X(y)$; $N=n-\frac1{24}$\\
$u_j$, $v_j$, $\Delta$ & $u_j$, the inverse coefficients & --- & $v_j$, the same coefficients for data $d_j=(2\mathsf A)^jc_j$\\
$k$, $b$, $M$ & slope, offset; $M$ the largest part, also an inverse order & $k$ the smallest part; $M$ an order & $k$ the minimum; $b$ a residue; $M$ the maximum, an order\\
$X_{\vol}$, $a_{\vol}$, $b_{\vol}$, $L_{\vol}$ & \multicolumn{3}{>{\raggedright\arraybackslash}p{5.25in}}{the transseries volume's core variable and data, never the Parts' $X$, $a$, $b$, $L$}\\
\bottomrule
\end{longtable}
\addtocounter{table}{-1}% the uncaptioned longtable must not consume a table number
\endgroup
\clearpage

\part{Fixed slope partition extremes and inverse asymptotics}\label{xpc:sl:part}
\setlength{\parindent}{0pt}\setlength{\parskip}{4pt}
\partsource{Bundle Report~216, archive \file{Report216-reproducibility.zip}, delivered as \file{Report216.tex} (20 pages), the base of the merge. Delivered title block ``Fixed slope partition extremes and inverse asymptotics / Report216 / 4 October 2026'' (empty author line). Printed as delivered, with \textbf{[write]} remarks and notes; section, statement, equation and table numbers are those of the delivery, its labels carry \file{xpc:sl:}, and its two generated tables are printed inline. This Part is self-contained.}
\begin{partabstract}
Fix integers $k\ge2$ and $b\ge0$. We count partitions of $n$ whose largest part $M$ and number of parts $L$ satisfy $M=kL+b$ or $M\ge kL+b$. Exact finite-product sieves give, for every fixed order, explicit expansions in powers of $(n-1/24)^{-1/2}$ relative to the principal partition scale. The leading coefficients are $k!$, with powers $k$ and $k-1$, respectively. The proof controls the discarded sieve terms on an entire Cauchy circle, including its $t^{-3/2}$ loss, then transfers each retained polynomial by a finite-shift form of the Rademacher formula. We give rational coefficient generators, the first four coefficients for A237753 and A237751, fixed-offset companions, and a triangular Lambert $W_{-1}$ inverse with shrinking two-ceiling uncertainty for the first index whose count is at least a target. Kot\v{e}\v{s}ovec's posted leading equality equivalent is credited; classical joint-extreme and general transfer results are compared in their actual parameter regimes. The claims concern fixed slope, fixed offset, and fixed order. Numerical diagnostics are not certified error bounds, and no historical priority, effective onset, or growing-order assertion is made.
\end{partabstract}

\section{The statistics and the result}\label{xpc:sl:sec:result}
For a nonempty integer partition $\lambda$, let $M(\lambda)$ be its largest part and $L(\lambda)$ its number of positive parts. Throughout this article $k\ge2$ and $b\ge0$ are fixed integers. Define
\begin{align}
 E_{k,b}(n)&=\#\{\lambda\vdash n:M(\lambda)=kL(\lambda)+b\},\label{xpc:sl:eq:Edef}\\
 T_{k,b}(n)&=\#\{\lambda\vdash n:M(\lambda)\ge kL(\lambda)+b\}.\label{xpc:sl:eq:Tdef}
\end{align}
The empty partition is excluded, so both counts are zero at $n=0$. Let $p(n)$ denote the unrestricted partition number, with $p(0)=1$ and $p(n)=0$ for negative integers $n$.

Conjugation interchanges $M$ and $L$. It therefore identifies $E_{k,0}$ with column $k$ of OEIS A350879, whose displayed convention is $kM=L$. In particular,
\begin{equation}\label{xpc:sl:eq:OEIS}
 E_{2,0}(n)=\mathrm{A237753}(n),\qquad
 T_{2,1}(n)=\mathrm{A237751}(n).
\end{equation}
For the second identity, the OEIS condition $2M<L$ conjugates to $M>2L$, which is $M\ge2L+1$ on integers. The offset $1$ is essential for the lower-order coefficients.

The leading equality equivalent
\begin{equation}\label{xpc:sl:eq:prior}
 E_{k,0}(n)\sim k!\left(\frac{\pi}{\sqrt{6n}}\right)^k p(n)
\end{equation}
is already attributed to V\'{a}clav Kot\v{e}\v{s}ovec, 17 October 2024, in A350879 and its $k=2$ specialization~\cite{oeis350879,oeis237753}. The equality generating function at $b=0$ is also recorded there. Our purpose is a fully justified, explicit fixed-order expansion and inversion for both statistics with fixed offsets. We do not claim that the leading scale or the general idea of partition-product transfer is new. Section~\ref{xpc:sl:sec:prior} explains the relationship with the relevant primary literature.

Set
\begin{equation}\label{xpc:sl:eq:scale}
 A=\frac{\pi^2}{6},\qquad N=n-\frac1{24},\qquad
 t=\sqrt{\frac{A}{N}},\qquad
 B(n)=\frac{\exp(2\sqrt{AN})}{4\sqrt3\,N}.
\end{equation}
The parameter $t$ is positive for integers $n\ge1$ and tends to zero. Write $X=E$ or $X=T$ for the kind of statistic and define
\begin{equation}\label{xpc:sl:eq:valuation}
 \epsilon_E=1,\quad\epsilon_T=0,\qquad
 s_E=k,\quad s_T=k-1.
\end{equation}
We sometimes write simply $s=s_X$ and $X_{k,b}$.

For $j\ge1$ let
\begin{equation}\label{xpc:sl:eq:Dj}
 D_j=D_j(k,b)=\frac{(2k+1)j^2+(2b-1)j}{2}.
\end{equation}
This is an integer: the numerator is even because $j^2-j$ is even. Define formal power series
\begin{align}
 G^E_{k,b}(z)&=\sum_{j\ge1}(-1)^{j-1}\ee^{-D_jz}
                   \prod_{r=j}^{kj}(1-\ee^{-rz}),\label{xpc:sl:eq:GE}\\
 G^T_{k,b}(z)&=\sum_{j\ge1}(-1)^{j-1}\ee^{-D_jz}
                   \prod_{r=j+1}^{kj}(1-\ee^{-rz}).\label{xpc:sl:eq:GT}
\end{align}
The $j$th summand has valuation $(k-1)j+\epsilon_X$. Thus every coefficient uses finitely many terms; no convergence of the infinite series near $z=0$ is asserted or needed. Put $g_\ell=[z^\ell]G^X_{k,b}(z)$.

Define terminating polynomials
\begin{equation}\label{xpc:sl:eq:Q}
 Q_\ell(t)=\sum_{h=0}^{\ell+1}
  \frac{(-1)^h(\ell+h+1)!}{(\ell+1-h)!\,h!}
  \left(\frac{t}{4A}\right)^h\qquad(\ell\ge0),
\end{equation}
and define $c_m=c_m(k,b,X)$ by the formal identity
\begin{equation}\label{xpc:sl:eq:cgenerator}
 \sum_{m\ge0}c_mt^m
 =\frac1{k!t^s}\sum_{\ell\ge s}g_\ell t^\ell Q_\ell(t).
\end{equation}
Each $c_m$ is a polynomial in $A^{-1}$ with rational coefficients. This is also a finite algorithm: to compute $c_0,\ldots,c_R$, only $g_s,\ldots,g_{s+R}$ are required.

\begin{theorem}[Every fixed order]\label{xpc:sl:thm:forward}
For every fixed nonnegative integer $R$,
\begin{equation}\label{xpc:sl:eq:main}
 X_{k,b}(n)=k!\,B(n)t^s
 \left(\sum_{m=0}^{R}c_m(k,b,X)t^m
              +O_{k,b,R}(t^{R+1})\right).
\end{equation}
Here $c_0=1$, and the coefficients are the explicit rational polynomials defined by \eqref{xpc:sl:eq:GE}--\eqref{xpc:sl:eq:cgenerator}. In particular,
\begin{align}
 E_{k,b}(n)&\sim k!\left(\frac{\pi}{\sqrt{6n}}\right)^k p(n),\label{xpc:sl:eq:Eleading}\\
 T_{k,b}(n)&\sim k!\left(\frac{\pi}{\sqrt{6n}}\right)^{k-1}p(n).\label{xpc:sl:eq:Tleading}
\end{align}
\end{theorem}
These are fixed-depth asymptotic statements. The constants and the onset may depend on $k,b,R$. They do not permit $k$, $b$, or $R$ to grow with $n$, and do not imply convergence or a numerical error bound for the full formal series.

The first relative coefficients are
\begin{align}
 c_1(k,b,E)&=-\frac{k(k+5)}4-b-12\ind_{\{k=2\}}
                    -\frac{(k+1)(k+2)}{4A},\label{xpc:sl:eq:c1E}\\
 c_1(k,b,T)&=-\frac{k(k+5)-2}4-b-6\ind_{\{k=2\}}
                    -\frac{k(k+1)}{4A}.\label{xpc:sl:eq:c1T}
\end{align}
In particular the leading tail companion for A237751 is
\begin{equation}\label{xpc:sl:eq:A237751leading}
 \mathrm{A237751}(n)\sim\frac{2\pi}{\sqrt{6n}}p(n).
\end{equation}
This is a consequence of the theorem, not a claim of historical firstness.

\paragraph{A literal discrepancy in the posted equivalent.}
At the source check on 4 October 2026, A350879 printed an alternative elementary exponential expression with $\exp(\sqrt{2\pi n/3})$, while its partition-normalized expression is \eqref{xpc:sl:eq:prior}. The correct rate equivalent to \eqref{xpc:sl:eq:prior} is $\exp(\pi\sqrt{2n/3})$. Indeed the usual equivalent for $p(n)$ gives
\begin{equation}\label{xpc:sl:eq:correctliteral}
 E_{k,0}(n)\sim
 \frac{k!\pi^k}{4\sqrt3\,6^{k/2}n^{1+k/2}}
       \exp\!\left(\pi\sqrt{\frac{2n}{3}}\right).
\end{equation}
The discrepancy is a missing factor of $\pi$ under the square root in that alternative expression. It does not invalidate the already posted $p(n)$-normalized leading law. No OEIS edit is part of this work.

\begin{remark}[{[write]} The OEIS entries, quoted, 7~October 2026]\label{xpc:sl:rem:oeis}
Read on 7~October 2026 in the internal format.
\begin{itemize}[leftmargin=*]
\item \emph{A350879}, revision \#51 (18~October 2024), author Seiichi Manyama
(21~January 2022): ``Triangle T(n,k), n >= 1, 1 <= k <= n, read by rows,
where T(n,k) is the number of partitions of n such that k*(greatest part) =
(number of parts).''; a comment gives the conjugate reading ``(greatest part)
= k*(number of parts)''; the column generating function of Manyama; columns
$k=1,\dots,7$ are A047993, A237753, A237756, A348163, A348164, A377107,
A377108; b-file by Andrew Howroyd (rows 1--50). Kot\v{e}\v{s}ovec's comment of
17~October 2024 reads: ``Column k > 1 is asymptotic to k! * Pi\^{}k *
exp(sqrt(2*Pi*n/3)) / (2\^{}((k+4)/2) * 3\^{}((k+1)/2) * n\^{}((k+2)/2)).
Equivalently, for fixed k > 1, T(n,k) \textasciitilde{} k! * Pi\^{}k *
A000041(n) / (6\^{}(k/2) * n\^{}(k/2)).'' The first form is false; see
Remark~\ref{xpc:sl:rem:a350879}.
\item \emph{A237753}, revision \#32 (11~July 2026), author Clark Kimberling
(13~February 2014): ``Number of partitions of n such that 2*(greatest part) =
(number of parts).''; Manyama's generating function; Kot\v{e}\v{s}ovec's
``a(n) \textasciitilde{} Pi\^{}2 * exp(Pi*sqrt(2*n/3)) / (4 * 3\^{}(3/2) *
n\^{}2)'' (17~October 2024), which is~\eqref{xpc:sl:eq:correctliteral} at
$k=2$ and is correct; b-file to $n=10000$ (Kot\v{e}\v{s}ovec, terms to 1000 from
Manyama).
\item \emph{A237751}, revision \#11 (24~January 2022), author Clark
Kimberling (13~February 2014): ``Number of partitions of n such that
2*(greatest part) < (number of parts).''; it adds the conjugate reading
``(greatest part) > 2*(number of parts)'' and ``a(n) = A000041(n) -
A237755(n).''; b-file by Manyama to $n=1000$. It states no asymptotic
formula; \eqref{xpc:sl:eq:A237751leading} is not recorded there.
\end{itemize}
The write computed $E_{2,0}$ and $T_{2,1}$ from~\eqref{xpc:sl:eq:Esieve}
and~\eqref{xpc:sl:eq:Tsieve} with its own code for $n\le3000$: every term of
the A237753 and A237751 b-files in that range agrees, as do brute-force counts
for $n\le40$; and $E_{k,0}(n)$ agrees with every entry of columns $k\ge2$ of
A350879's b-file. Nothing was submitted to the OEIS.
\end{remark}

\begin{remark}[{[write]} A350879's alternative form is false, 7~October 2026]\label{xpc:sl:rem:a350879}
\emph{Claim.} For every fixed $k\ge2$ the first form of
Kot\v{e}\v{s}ovec's comment in A350879 (Remark~\ref{xpc:sl:rem:oeis}),
\[
T(n,k)\sim\frac{k!\,\pi^k\exp\bigl(\sqrt{2\pi n/3}\bigr)}{2^{(k+4)/2}3^{(k+1)/2}n^{(k+2)/2}},
\]
is false, while its second form, the $p(n)$-normalized law~\eqref{xpc:sl:eq:prior},
is true. \emph{Proof.} By conjugation $T(n,k)=E_{k,0}(n)$, and
\eqref{xpc:sl:eq:correctliteral} (Theorem~\ref{xpc:sl:thm:forward} with the
Hardy--Ramanujan equivalent of $p(n)$) gives $T(n,k)\sim
k!\pi^k\exp(\pi\sqrt{2n/3})/(4\sqrt3\,6^{k/2}n^{1+k/2})$. Since
$2^{(k+4)/2}3^{(k+1)/2}=4\sqrt3\,6^{k/2}$, the quotient of the posted first
form by this equivalent is $\exp\bigl(-(\pi-\sqrt\pi)\sqrt{2n/3}\bigr)\to0$,
so the first form is not asymptotic to $T(n,k)$. Only the exponential rate is
wrong: $\sqrt{2\pi n/3}$ lacks a factor $\pi$ under the root, as the source
says in Section~\ref{xpc:sl:sec:result}; the constant and the power of $n$ are
right. \emph{Numerically} (not part of the proof), $A237753(3000)$ is
$0.7274\ldots$ times the correct equivalent and $2.84\ldots\cdot10^{26}$
times the posted first form. No OEIS edit was made.
\end{remark}

\section{Exact finite product sieves}\label{xpc:sl:sec:sieves}
For $|q|<1$, write
\[
 (a;q)_m=\prod_{r=0}^{m-1}(1-aq^r),\qquad
 (q;q)_m=\prod_{r=1}^m(1-q^r),\qquad
 P(q)=\frac1{(q;q)_\infty}=\sum_{n\ge0}p(n)q^n.
\]
Empty finite products equal one.

\begin{proposition}[Exact identities]\label{xpc:sl:prop:sieve}
As formal power series and as analytic functions for $|q|<1$,
\begin{align}
 F^E_{k,b}(q):=\sum_{n\ge0}E_{k,b}(n)q^n
  &=P(q)\sum_{j\ge1}(-1)^{j-1}q^{D_j}
                         \prod_{r=j}^{kj}(1-q^r),\label{xpc:sl:eq:Esieve}\\
 F^T_{k,b}(q):=\sum_{n\ge0}T_{k,b}(n)q^n
  &=P(q)\sum_{j\ge1}(-1)^{j-1}q^{D_j}
                         \prod_{r=j+1}^{kj}(1-q^r).\label{xpc:sl:eq:Tsieve}
\end{align}
For every fixed coefficient $q^n$, only indices with $D_j\le n$ contribute.
\end{proposition}
\begin{proof}
For a fixed length $l\ge1$, subtracting one from every part gives the unrestricted generating function $q^l/(q;q)_l$. If $M\le kl+b-1$, the remaining diagram fits a rectangle of height $l$ and width $kl+b-2$. Its generating function is
\[
 q^l\frac{(q^{kl+b-1};q)_l}{(q;q)_l}.
\]
The fixed-length tail is therefore
\begin{equation}\label{xpc:sl:eq:fixedlength}
 F^T_{k,b}(q)=\sum_{l\ge1}\frac{q^l}{(q;q)_l}
                      \bigl(1-(q^{kl+b-1};q)_l\bigr).
\end{equation}
Using the finite $q$-binomial identity in its explicit form
\begin{equation}\label{xpc:sl:eq:qbinomial}
 (a;q)_l=\sum_{j=0}^l
  \frac{(q;q)_l}{(q;q)_j(q;q)_{l-j}}
                   (-a)^j q^{j(j-1)/2},
\end{equation}
the $j$th selected term in \eqref{xpc:sl:eq:fixedlength} is
\[
 \frac{(-1)^{j-1}q^{l+(kl+b-1)j+j(j-1)/2}}
      {(q;q)_j(q;q)_{l-j}}\qquad(1\le j\le l).
\]
Set $l=j+h$. The exponent becomes $D_j+(kj+1)h$, and Euler's identity gives
\begin{equation}\label{xpc:sl:eq:euler}
 \sum_{h\ge0}\frac{q^{(kj+1)h}}{(q;q)_h}
 =\frac1{(q^{kj+1};q)_\infty}=P(q)(q;q)_{kj}.
\end{equation}
After division by $(q;q)_j$ the remaining factors are $j+1,\ldots,kj$, proving \eqref{xpc:sl:eq:Tsieve}. Finally
\begin{equation}\label{xpc:sl:eq:taildiff}
 E_{k,b}(n)=T_{k,b}(n)-T_{k,b+1}(n),
 \qquad D_j(k,b+1)=D_j(k,b)+j,
\end{equation}
so subtraction inserts the missing factor $1-q^j$ and gives \eqref{xpc:sl:eq:Esieve}.

All manipulations are coefficientwise legitimate: the exponents are nonnegative, and only finitely many length and sieve indices can affect a fixed degree; in particular $D_j\to\infty$ quadratically. The absolute bounds in Section~\ref{xpc:sl:sec:circle} justify the analytic interchanges near the unit circle. Any smaller closed disk is dominated by a larger such radius, so they also establish the analytic assertions throughout the open disk.
\end{proof}

For the equality remainder it is advantageous to start from a different exact expression, instead of subtracting two separately bounded tails.
\begin{lemma}[Equality rectangle]\label{xpc:sl:lem:rectangle}
One has
\begin{equation}\label{xpc:sl:eq:rectangle}
 F^E_{k,b}(q)=\sum_{m\ge0}q^{(k+1)m+k+b}
                   \frac{(q^{k(m+1)+b};q)_m}{(q;q)_m}.
\end{equation}
Expanding its numerator through elementary-symmetric degrees $i=0,\ldots,J-1$ and summing over $m$ gives precisely the terms $j=1,\ldots,J$ in \eqref{xpc:sl:eq:Esieve}.
\end{lemma}
\begin{proof}
Take $L=m+1$ and $M=k(m+1)+b$. Removing the first row and first column removes $M+L-1=(k+1)m+k+b$ cells. What remains is an arbitrary diagram in a rectangle of height $m$ and width $M-1$. Its Gaussian-binomial generating function is
\[
 \frac{(q^M;q)_m}{(q;q)_m},
\]
which proves \eqref{xpc:sl:eq:rectangle}. Selecting degree $i$ in \eqref{xpc:sl:eq:qbinomial} and then putting $m=i+h$, the $h$ coefficient in the exponent is $k(i+1)+1$. Applying \eqref{xpc:sl:eq:euler} with $j=i+1$ leaves $(q;q)_{kj}/(q;q)_{j-1}$, namely the factors $j,\ldots,kj$, with exponent $D_j$. Thus the truncation indices agree exactly, including $i=0\leftrightarrow j=1$.
\end{proof}

\section{A remainder bound on the entire circle}\label{xpc:sl:sec:circle}
A radial expansion of $G^X_{k,b}(t)$ alone does not establish a coefficient asymptotic. We instead truncate before extracting coefficients and control the discarded terms at every angle on a full circle.

Let $r=\ee^{-t}$, $\delta=1-r$, and $|q|=r$, with $t>0$ tending to zero. Then $\delta\sim t$. For complex numbers $x_1,\ldots,x_m$, put $S=\sum_i|x_i|$. If a product $\prod_i(1-x_i)$ is truncated before elementary-symmetric degree $J$, its remainder is bounded by
\begin{equation}\label{xpc:sl:eq:elementary}
 \sum_{i\ge J}e_i(|x_1|,\ldots,|x_m|)
 \le\sum_{i\ge J}\frac{S^i}{i!}
 \le\frac{S^J\ee^S}{J!}.
\end{equation}
This follows from $i!e_i\le S^i$ and $(J+h)!\ge J!h!$. Also
\begin{align}
 \frac1{|(q;q)_m|}&\le\frac1{(r;r)_m}
   =P(r)(r^{m+1};r)_\infty\notag\\
 &\le P(r)\exp\!\left(-\frac{r^{m+1}}{\delta}\right),\label{xpc:sl:eq:denombound}
\end{align}
using $\log(1-x)\le-x$ for $0<x<1$.

\begin{lemma}[Geometric mesh]\label{xpc:sl:lem:mesh}
For fixed real numbers $a,c>0$, and $u_m=r^{m+1}$,
\begin{equation}\label{xpc:sl:eq:mesh}
 \sum_{m\ge0}u_m^a\exp(-cu_m/\delta)
      =O_{a,c}(\delta^a/t).
\end{equation}
\end{lemma}
\begin{proof}
The function $h(x)=\exp(-ax-c\ee^{-x}/\delta)$ is unimodal on $[0,\infty)$. Comparing each side of its maximum with the corresponding integral bounds its mesh sum by
\[
 \sum_{m\ge0}h(t(m+1))\le t^{-1}\int_0^\infty h(x)\,dx+2\sup_{x\ge0}h(x).
\]
With $u=\ee^{-x}$,
\[
 \int_0^\infty h(x)\,dx
 =\int_0^1u^{a-1}\ee^{-cu/\delta}\,du
 \le\Gamma(a)(\delta/c)^a.
\]
The maximum of $u^a\ee^{-cu/\delta}$ on $u>0$ is $(a\delta/(c\ee))^a$. Since $t\to0$, these bounds give \eqref{xpc:sl:eq:mesh}.
\end{proof}

Write $F^X_{k,b,J}$ for \eqref{xpc:sl:eq:Esieve} or \eqref{xpc:sl:eq:Tsieve} with only $j=1,\ldots,J$ retained. For $J=0$ the truncated expression is zero.
\begin{proposition}[Global finite-depth bound]\label{xpc:sl:prop:global}
For every fixed nonnegative integer $J$,
\begin{align}
 \sup_{|q|=r}|F^E_{k,b}(q)-F^E_{k,b,J}(q)|
  &=O_{k,b,J}\!\left(P(r)t^{k+(k-1)J}\right),\label{xpc:sl:eq:BE}\\
 \sup_{|q|=r}|F^T_{k,b}(q)-F^T_{k,b,J}(q)|
  &=O_{k,b,J}\!\left(P(r)t^{\,(k-1)(J+1)}\right).\label{xpc:sl:eq:BT}
\end{align}
\end{proposition}
\begin{proof}
For equality, use \eqref{xpc:sl:eq:rectangle} and set $u=r^{m+1}$. The sum of the numerator-factor moduli is
\[
 S_m=r^b u^k\frac{1-r^m}{\delta}
       \le\frac{u^k(1-u)}{\delta}.
\]
Here $b\ge0$ and $r^m=u/r\ge u$. Since $k\ge2$,
\begin{equation}\label{xpc:sl:eq:quarter}
 u^{k-1}(1-u)\le u(1-u)\le\frac14\qquad(0<u<1).
\end{equation}
Consequently $-u/\delta+S_m\le-3u/(4\delta)$. The exterior power in \eqref{xpc:sl:eq:rectangle} has modulus $r^{b-1}u^{k+1}\le r^{-1}u^{k+1}$. Combining \eqref{xpc:sl:eq:elementary} and \eqref{xpc:sl:eq:denombound}, the error after degrees $0,\ldots,J-1$ is at most
\[
 \frac{P(r)r^{-1}}{J!\delta^J}
      \sum_{m\ge0}u^{k+1+kJ}\exp\!\left(-\frac{3u}{4\delta}\right).
\]
The mesh lemma makes this
\[
 O\!\left(P(r)\frac{\delta^{k+1+(k-1)J}}t\right)
   =O\!\left(P(r)t^{k+(k-1)J}\right).
\]
Lemma~\ref{xpc:sl:lem:rectangle} identifies the retained terms, proving \eqref{xpc:sl:eq:BE}.

For the tail, use \eqref{xpc:sl:eq:fixedlength} and put $u=r^{l+1}$. Now
\[
 S_l=r^{b-k-1}u^k\frac{1-u/r}{\delta}
    \le r^{-k-1}\frac{u^k(1-u)}{\delta}.
\]
For $r$ sufficiently close to $1$, depending only on the fixed $k$, we have $r^{-k-1}\le2$. Thus \eqref{xpc:sl:eq:quarter} gives
\[
 -u/\delta+S_l\le-u/(2\delta).
\]
The prefactor $r^l$ equals $r^{-1}u$. Retaining degrees $1,\ldots,J$ in $1-\prod(1-x_i)$ leaves degree at least $J+1$. Its error is bounded by a constant times
\[
 P(r)\delta^{-J-1}
    \sum_{l\ge1}u^{1+k(J+1)}\exp(-u/(2\delta)).
\]
The mesh estimate gives the power
\[
 \delta^{-J-1}\frac{\delta^{1+k(J+1)}}t
       =O\!\left(t^{(k-1)(J+1)}\right),
\]
proving \eqref{xpc:sl:eq:BT}.

The same bounds, with $J=0$ and the numerator absolute sum bounded by $\ee^{S_m}$ or $\ee^{S_l}$, prove absolute summability of the expanded products. For the tail one may also use $\ee^{S_l}-1\le S_l\ee^{S_l}$. This supplies the absolute interchange used in Proposition~\ref{xpc:sl:prop:sieve}. Every inequality has been expressed in $r$, so the estimates hold on the entire circle, not only on its positive radius.
\end{proof}

\subsection{The Cauchy cost and the required depth}\label{xpc:sl:sec:cauchy}
The eta inversion formula~\cite[23.18.5--23.18.7]{dlmf} yields the exact identity
\begin{equation}\label{xpc:sl:eq:eta}
 P(\ee^{-t})=\sqrt{\frac{t}{2\pi}}
   \exp\!\left(\frac At-\frac t{24}\right)
   P(\ee^{-4\pi^2/t}).
\end{equation}
For clarity, substituting $\tau=it/(2\pi)$ into
$\eta(-1/\tau)=\sqrt{-i\tau}\,\eta(\tau)$ and using
$\eta(\tau)=\ee^{\pi i\tau/12}\prod_{r\ge1}(1-\ee^{2\pi i r\tau})$
gives precisely \eqref{xpc:sl:eq:eta}. The last factor equals
$1+O(\ee^{-4\pi^2/t})$ by its convergent power series at zero.

Choose $t$ as in \eqref{xpc:sl:eq:scale}. Because $Nt=A/t$,
\begin{equation}\label{xpc:sl:eq:cauchycost}
 \frac{r^{-n}P(r)}{B(n)}
   =\frac{4\sqrt3\,A}{\sqrt{2\pi}}t^{-3/2}
       \left(1+O(\ee^{-4\pi^2/t})\right).
\end{equation}
Cauchy's coefficient inequality and Proposition~\ref{xpc:sl:prop:global} therefore give remainders
\begin{align*}
 O\!\left(B(n)t^{k+(k-1)J-3/2}\right)&\quad\text{for equality},\\
 O\!\left(B(n)t^{(k-1)(J+1)-3/2}\right)&\quad\text{for the tail}.
\end{align*}
Relative to $B(n)t^s$, both powers are $(k-1)J-3/2$. Thus the convenient integer choice
\begin{equation}\label{xpc:sl:eq:Jdepth}
 J\ge\ceil{\frac{R+3}{k-1}}
\end{equation}
ensures a remainder $O(B(n)t^{s+R+1})$. In fact the left side of
$(k-1)J\ge R+5/2$ is integral, so \eqref{xpc:sl:eq:Jdepth} is the least depth guaranteed by this particular estimate. The loss $t^{-3/2}$ has been included; a merely radial valuation count would miss it.

For this $J$, the retained expression is $P(q)$ times a fixed polynomial in $q$. Its asymptotic analysis requires only finitely many shifted partition numbers, which we now control separately.

\section{Finite shift transfer from the partition formula}\label{xpc:sl:sec:transfer}
We use the Rademacher formula in Johansson's exact $U$-series normalization~\cite[(1.4)--(1.8)]{johansson}, rather than using a leading Hardy--Ramanujan equivalent as if it were an all-order error estimate.

\begin{lemma}[One-term exponential gap]\label{xpc:sl:lem:gap}
Let
\begin{equation}\label{xpc:sl:eq:f}
 f(N)=\frac{\exp(2\sqrt{AN})}{4\sqrt3\,N}
                \left(1-\frac1{2\sqrt{AN}}\right).
\end{equation}
For positive integers $n$ tending to infinity,
\begin{equation}\label{xpc:sl:eq:p-gap}
 p(n)=f(n-1/24)+O\!\left(\exp(\sqrt{A(n-1/24)})\right).
\end{equation}
The same estimate holds at each of finitely many fixed integer shifts of $n$.
\end{lemma}
\begin{proof}
Put $w=2\sqrt{AN}$. The cited convergent series is
\begin{equation}\label{xpc:sl:eq:U}
 p(n)=\frac1{2\sqrt3\,N}\sum_{a\ge1}
       \frac{A_a(n)}{\sqrt a}\,U(w/a),\qquad
 U(x)=\cosh x-\frac{\sinh x}{x}.
\end{equation}
Here $A_1(n)=1$, and $A_a(n)$ is a sum of at most $a$ unit-modulus terms, so $|A_a(n)|\le a$. For $x\ge0$,
\[
 U(x)=\frac1x\int_0^x u\sinh u\,du,
 \qquad 0\le U(x)\le\frac{x^2\ee^x}{3},
\]
where at $x=0$ the formulas are understood by continuity. The bound follows from $\sinh u\le u\ee^u\le u\ee^x$. Hence the terms with $a\ge2$ have total absolute value at most a constant times
\[
 N^{-1}w^2\ee^{w/2}\sum_{a\ge2}a^{-3/2}=O(\ee^{w/2}),
\]
since $w^2/N=4A$. The $a=1$ term equals
\[
 \frac1{4\sqrt3\,N}
       \left(\ee^w(1-w^{-1})+\ee^{-w}(1+w^{-1})\right).
\]
Its growing term is exactly \eqref{xpc:sl:eq:f}, and its other term is exponentially small. This proves \eqref{xpc:sl:eq:p-gap}. Finitely many fixed shifts change $w$ by $O(N^{-1/2})$, preserving the gap.
\end{proof}

\begin{lemma}[Exact derivative polynomials]\label{xpc:sl:lem:Q}
For every fixed $\ell\ge0$,
\begin{equation}\label{xpc:sl:eq:fderiv}
 f^{(\ell)}(N)=B(N+1/24)t^\ell Q_\ell(t),\qquad t=\sqrt{A/N},
\end{equation}
where derivatives are with respect to the real variable $N$, and $Q_\ell$ is \eqref{xpc:sl:eq:Q}.
\end{lemma}
\begin{proof}
The case $\ell=0$ follows from $Q_0(t)=1-t/(2A)$. Direct differentiation gives
\[
 \frac{B'(N+1/24)}{B(N+1/24)}=t-\frac{t^2}A,
 \qquad \frac{dt}{dN}=-\frac{t^3}{2A}.
\]
Thus \eqref{xpc:sl:eq:fderiv} is equivalent to the recurrence
\begin{equation}\label{xpc:sl:eq:Qrec}
 Q_{\ell+1}(t)=\left(1-\frac{(\ell+2)t}{2A}\right)Q_\ell(t)
                     -\frac{t^2}{2A}Q_\ell'(t).
\end{equation}
To verify the terminating formula directly, write $Q_\ell(t)=\sum_hq_{\ell,h}t^h$ with coefficients outside $0\le h\le\ell+1$ set to zero. The recurrence is
\[
 q_{\ell+1,h}=q_{\ell,h}-\frac{\ell+h+1}{2A}q_{\ell,h-1}.
\]
Substituting the factorial expression in \eqref{xpc:sl:eq:Q} satisfies this identity, including both endpoints. Induction proves the result.
\end{proof}

\begin{proposition}[Polynomial transfer]\label{xpc:sl:prop:transfer}
Let $H(q)=\sum_dh_dq^d$ be a fixed polynomial. For every fixed nonnegative integer $L$,
\begin{equation}\label{xpc:sl:eq:transfer}
 [q^n]P(q)H(q)=B(n)\sum_{\ell=0}^{L}
     [z^\ell]H(\ee^{-z})\,t^\ell Q_\ell(t)
          +O_{H,L}(B(n)t^{L+1}).
\end{equation}
\end{proposition}
\begin{proof}
Coefficient extraction is the exact finite sum $\sum_dh_dp(n-d)$. For sufficiently large $n$, every $n-d$ in that sum is positive. Lemma~\ref{xpc:sl:lem:gap} replaces $p(n-d)$ by $f(N-d)$ with an error exponentially smaller than $B(n)$ times any fixed power of $t$. Taylor's theorem at the finitely many bounded shifts $d$ gives
\[
 f(N-d)=\sum_{\ell=0}^{L}\frac{(-d)^\ell}{\ell!}f^{(\ell)}(N)
            +O_{d,L}(B(n)t^{L+1}).
\]
Indeed Lemma~\ref{xpc:sl:lem:Q} bounds $f^{(L+1)}(N-u)$ by $O(B(n)t^{L+1})$ uniformly for $u$ in the fixed interval containing these shifts. Summing and using
\[
 [z^\ell]H(\ee^{-z})=\sum_dh_d\frac{(-d)^\ell}{\ell!}
\]
proves \eqref{xpc:sl:eq:transfer}. Only the explicit smooth function $f$ is differentiated; no asymptotic error for integer partition numbers is differentiated.
\end{proof}

\begin{proof}[Proof of Theorem~\ref{xpc:sl:thm:forward}]
Choose $J$ by \eqref{xpc:sl:eq:Jdepth}, and let $H$ be the finite sieve polynomial consisting of the first $J$ terms in \eqref{xpc:sl:eq:Esieve} or \eqref{xpc:sl:eq:Tsieve}. Section~\ref{xpc:sl:sec:cauchy} replaces $F^X_{k,b}$ by $P(q)H(q)$ with coefficient error $O(B(n)t^{s+R+1})$. Apply Proposition~\ref{xpc:sl:prop:transfer} with $L=s+R$.

The coefficients below degree $s$ in $H(\ee^{-z})$ vanish exactly. Also the omitted $j>J$ terms have valuation at least
\[
 (k-1)(J+1)+\epsilon_X=s+(k-1)J>s+R,
\]
so the coefficients through degree $s+R$ agree with those of $G^X_{k,b}$. Discarding total $t$-degrees beyond $s+R$ in \eqref{xpc:sl:eq:transfer} gives precisely \eqref{xpc:sl:eq:main} and \eqref{xpc:sl:eq:cgenerator}. The $j=1$ product has leading coefficient $\prod_{r=1}^kr=k!$ for equality, and $\prod_{r=2}^kr=k!$ for the tail. All $j\ge2$ terms have strictly larger valuation. Therefore $g_s=k!$ and $c_0=1$.

Finally \eqref{xpc:sl:eq:p-gap} gives $p(n)=B(n)(1-t/(2A))$ plus an exponentially smaller term; $t\sim\pi/\sqrt{6n}$. This proves \eqref{xpc:sl:eq:Eleading}--\eqref{xpc:sl:eq:Tleading}.
\end{proof}

\section{Coefficients and fixed offset companions}\label{xpc:sl:sec:coefficients}
Formula \eqref{xpc:sl:eq:cgenerator} can be written without formal-series notation as
\begin{equation}\label{xpc:sl:eq:cfinite}
 c_m=\frac1{k!}\sum_{\substack{s\le\ell\le s+m\\0\le h=s+m-\ell\le\ell+1}}
       g_\ell\frac{(-1)^h(\ell+h+1)!}
                         {(\ell+1-h)!\,h!\,(4A)^h}.
\end{equation}
The constrained sum avoids any negative-factorial convention. The finite upper bound for a purely algebraic coefficient calculation is
\begin{equation}\label{xpc:sl:eq:Jformal}
 j\le\left\lfloor\frac{s+R-\epsilon_X}{k-1}\right\rfloor
\end{equation}
when computing $g_0,\ldots,g_{s+R}$. This algebraic bound differs from the larger analytical depth \eqref{xpc:sl:eq:Jdepth}; the latter proves the error bound and cannot be replaced by a radial valuation argument.

\subsection{The first correction and the exceptional second term}
For the $j=1$ term, $D_1=k+b$ and
\[
 \prod_{r\in I}(1-\ee^{-rz})
   =\left(\prod_{r\in I}r\right)z^{|I|}
             \left(1-\frac z2\sum_{r\in I}r+O(z^2)\right).
\]
For equality $I=\{1,\ldots,k\}$, giving a relative radial coefficient
$-k(k+5)/4-b$. For the tail $I=\{2,\ldots,k\}$, it is
$-[k(k+5)-2]/4-b$.

The $j=2$ summand begins at relative degree $k-1$. It enters the first correction only for $k=2$. Its signed leading coefficients are $-(2\cdot3\cdot4)=-24$ for equality and $-(3\cdot4)=-12$ for the tail. Dividing by $k!=2$ gives the extra $-12$ and $-6$. Finally the coefficient of $t$ in $Q_s$ is $-(s+1)(s+2)/(4A)$. This proves \eqref{xpc:sl:eq:c1E}--\eqref{xpc:sl:eq:c1T}. In particular, retaining only one selected large part captures the leading equivalent but loses a first-order contribution at slope $2$.

Normalizing instead by the exact partition number gives
\begin{align}
 \frac{E_{k,b}(n)}{k!t^kp(n)}
  &=1-\left(\frac{k(k+5)}4+b+12\ind_{\{k=2\}}
                  +\frac{k(k+3)}{4A}\right)t+O(t^2),\label{xpc:sl:eq:Epfirst}\\
 \frac{T_{k,b}(n)}{k!t^{k-1}p(n)}
  &=1-\left(\frac{k(k+5)-2}4+b+6\ind_{\{k=2\}}
                  +\frac{(k-1)(k+2)}{4A}\right)t+O(t^2).\label{xpc:sl:eq:Tpfirst}
\end{align}
At all fixed orders the conversion is multiplication by the formal reciprocal of $1-t/(2A)$. The shift $N=n-1/24$ is kept throughout, so these formulas should not be mixed termwise with coefficients in powers of $n^{-1/2}$ without re-expanding that shift.

The leading terms are independent of the fixed offset $b$, whereas the first relative coefficient changes by exactly $-b$. Identity \eqref{xpc:sl:eq:taildiff} holds exactly. At higher orders the factors $\ee^{-bjz}$ in the sieve determine the offset dependence algorithmically. This is a fixed-offset statement; summing these expansions over a growing range of offsets would require a new uniform remainder analysis.

\subsection{The two OEIS specializations}
For A237753, corresponding to $(k,b,X)=(2,0,E)$, the radial coefficients begin
\begin{equation}\label{xpc:sl:eq:radE}
 G^E_{2,0}(z)=2z^2-31z^3+\frac{2090}{3}z^4
                 -\frac{79007}{4}z^5+\frac{30640414}{45}z^6+\cdots.
\end{equation}
Its first four normalized coefficients are
\begin{align}
 c_0&=1,\notag\\
 c_1&=-\frac{31}{2}-\frac3A,\notag\\
 c_2&=\frac{1045}{3}+\frac{155}{2A}+\frac{15}{4A^2},\label{xpc:sl:eq:Ecoefs}\\
 c_3&=-\frac{79007}{8}-\frac{5225}{2A}
                    -\frac{1395}{8A^2}-\frac{15}{8A^3}.\notag
\end{align}
They multiply $2B(n)t^2$ in \eqref{xpc:sl:eq:main}.

For A237751, corresponding to $(k,b,X)=(2,1,T)$,
\begin{equation}\label{xpc:sl:eq:radT}
 G^T_{2,1}(z)=2z-20z^2+\frac{931}{3}z^3
                    -\frac{20270}{3}z^4+\frac{11331751}{60}z^5+\cdots,
\end{equation}
and the coefficients multiplying $2B(n)t$ are
\begin{align}
 c_0&=1,\notag\\
 c_1&=-10-\frac3{2A},\notag\\
 c_2&=\frac{931}{6}+\frac{30}{A}+\frac3{4A^2},\label{xpc:sl:eq:Tcoefs}\\
 c_3&=-\frac{10135}{3}-\frac{4655}{6A}-\frac{75}{2A^2}.\notag
\end{align}
The absence of an $A^{-3}$ term in the last line is consistent with the terminating degree of $Q_1$. These expansions are not obtained by merely substituting $k=2$ into a formula that discards the $j=2$ sieve term.

\section{All fixed order inverse asymptotics}\label{xpc:sl:sec:inverse}
Fix one of the sequences in Theorem~\ref{xpc:sl:thm:forward}. For real $n>1/24$ put
\begin{equation}\label{xpc:sl:eq:inverse-scale}
 w=2\sqrt{A(n-1/24)},\quad d=s+2,\quad
 D=\frac{k!\,2^s A^{s+1}}{\sqrt3}.
\end{equation}
Then $t=2A/w$ and the forward result is
\begin{equation}\label{xpc:sl:eq:w-forward}
 X_{k,b}(n)=D\ee^ww^{-d}
    \left(1+\sum_{j=1}^Rc_j\left(\frac{2A}{w}\right)^j
                        +O(w^{-R-1})\right)
\end{equation}
on integer $n$. The expression outside the error is an explicit smooth function; the integer count is not being extended or differentiated.

\subsection{Lambert core and triangular reversion}
For sufficiently large real $y$, let
\begin{equation}\label{xpc:sl:eq:w0}
 w_0=-d\,W_{-1}\!\left(-\frac1d\left(\frac Dy\right)^{1/d}\right).
\end{equation}
This is the large positive solution of
\begin{equation}\label{xpc:sl:eq:leading-inverse}
 D\ee^{w_0}w_0^{-d}=y.
\end{equation}
The branch $W_{-1}$, not $W_0$, is required: if $y>D(\ee/d)^d$, then $w_0>d$ and $w_0\to\infty$ as $y\to\infty$.

Let $H(v)=1+\sum_{j\ge1}c_jv^j$ be a formal series. Set $z=w_0^{-1}$, and seek
\begin{equation}\label{xpc:sl:eq:delta}
 \Delta(z)=\sum_{j\ge1}u_jz^j
\end{equation}
so that substituting $w=w_0+\Delta(z)$ in the logarithm of the forward expression yields zero residual. Using \eqref{xpc:sl:eq:leading-inverse}, the formal equation is
\begin{equation}\label{xpc:sl:eq:reversion}
 \Delta-d\log(1+z\Delta)
       +\log H\!\left(\frac{2Az}{1+z\Delta}\right)=0.
\end{equation}
All logarithms here are formal logarithms of series with constant term one.

\begin{proposition}[Triangular inverse coefficients]\label{xpc:sl:prop:inverse-coefs}
Equation \eqref{xpc:sl:eq:reversion} has a unique formal solution of the form \eqref{xpc:sl:eq:delta}. For every $m\ge1$, its coefficient of $z^m$ equals $u_m$ plus a polynomial in $u_1,\ldots,u_{m-1}$ and $c_1,\ldots,c_m$, with coefficients rational in $A$ and $d$. Thus $u_m$ is determined by a finite triangular calculation. The first coefficients are
\begin{align}
 u_1&=-2Ac_1,\label{xpc:sl:eq:u1}\\
 u_2&=-2Adc_1-4A^2c_2+2A^2c_1^2,\label{xpc:sl:eq:u2}\\
 u_3&=-\frac{8A^3}{3}c_1^3+8A^3c_1c_2-8A^3c_3
          +2A^2dc_1^2-4A^2c_1^2\notag\\
     &\hspace{13mm}-4A^2dc_2-2Ad^2c_1.\label{xpc:sl:eq:u3}
\end{align}
\end{proposition}
\begin{proof}
The direct term $\Delta$ contributes $u_m$ to degree $m$. In $z\Delta$, the same unknown first appears at degree $m+1$, and in the substitution into $H$ it can only appear at still higher degree. Therefore all other degree-$m$ terms depend only on earlier $u$'s. The constant term vanishes; solving degree by degree proves existence and uniqueness. Expanding the two formal logarithms through degrees one, two, and three gives \eqref{xpc:sl:eq:u1}--\eqref{xpc:sl:eq:u3}.
\end{proof}

For each fixed $M\ge1$ define the real center
\begin{equation}\label{xpc:sl:eq:xM}
 \Delta_M=\sum_{j=1}^M\frac{u_j}{w_0^j},\qquad
 x_M(y)=\frac1{24}+\frac{(w_0+\Delta_M)^2}{4A}.
\end{equation}
The square in this definition is retained in full, even though some terms lie beyond its promised order. Define also the Lambert-only center $x_0=1/24+w_0^2/(4A)$ for numerical comparison.

\begin{lemma}[Error of the smooth inverse center]\label{xpc:sl:lem:inverse-error}
Let
\begin{equation}\label{xpc:sl:eq:smoothM}
 \mathcal F_M(x)=D\ee^{w(x)}w(x)^{-d}
           \left(1+\sum_{j=1}^Mc_j(2A/w(x))^j\right),
 \qquad w(x)=2\sqrt{A(x-1/24)}.
\end{equation}
On its large positive branch this function is positive and strictly increasing. If $\widehat x_M$ is its large real solution of $\mathcal F_M(x)=y$, then
\begin{equation}\label{xpc:sl:eq:smooth-error}
 \widehat x_M=x_M(y)+O(w_0^{-M}).
\end{equation}
\end{lemma}
\begin{proof}
The polynomial in parentheses tends to one and is positive for large $x$. Its logarithmic derivative with respect to $w$ is $O(w^{-2})$, so
\begin{equation}\label{xpc:sl:eq:log-deriv}
 \frac{d}{dx}\log\mathcal F_M(x)
    =\frac{2A}{w}\left(1-\frac dw+O(w^{-2})\right)
    =t+O(t^2)>0
\end{equation}
for large $x$. This also gives a unique large root when $y$ is large. Substitution of $w_0+\Delta_M$ in the logarithmic equation leaves residual $O(w_0^{-M-1})$, by Proposition~\ref{xpc:sl:prop:inverse-coefs}. The derivative of that equation with respect to $w$ is $1+O(w_0^{-1})$ in a bounded neighborhood. The mean value theorem therefore gives an error $O(w_0^{-M-1})$ in $w$. Since $dx/dw=w/(2A)=O(w_0)$, the error in $x$ is $O(w_0^{-M})$.
\end{proof}

In particular the first two index corrections, viewed as an expansion of the smooth inverse, are
\begin{equation}\label{xpc:sl:eq:indexcorr}
 \widehat x_M=\frac1{24}+\frac{w_0^2}{4A}-c_1
       +\frac{-dc_1-2Ac_2+Ac_1^2}{w_0}
       +O(w_0^{-2})\qquad(M\ge2).
\end{equation}
For $M=1$ the constant correction $-c_1$ is valid with error $O(w_0^{-1})$. Thus this constant is $31/2+3/A$ for A237753 and $10+3/(2A)$ for A237751. These constants improve a real center, not a universal rounding prescription.

\subsection{Integer thresholds and the two ceilings}\label{xpc:sl:sec:threshold}
Use the exact threshold convention
\begin{equation}\label{xpc:sl:eq:nu}
 \nu_X(y)=\min\{n\ge0:X_{k,b}(n)\ge y\}.
\end{equation}
The inequality is weak; this matters at targets equal to a sequence value.

\begin{lemma}[Eventual strict increase]\label{xpc:sl:lem:monotone}
Each sequence in Theorem~\ref{xpc:sl:thm:forward} is positive, unbounded, and strictly increasing for all sufficiently large integers $n$.
\end{lemma}
\begin{proof}
Write the expansion through relative order one with remainder $O(t^2)$. The ratio of successive smooth leading factors $B(n)t(n)^s$ is $1+t+O(t^2)$, by explicit differentiation or direct expansion of \eqref{xpc:sl:eq:scale}. The factors $1+c_1t+O(t^2)$ at $n$ and $n+1$ have ratio $1+O(t^2)$ because $t(n+1)-t(n)=O(t^3)$. Consequently
\begin{equation}\label{xpc:sl:eq:ratio}
 \frac{X_{k,b}(n+1)}{X_{k,b}(n)}=1+t+O(t^2)>1
\end{equation}
for all sufficiently large $n$. The leading equivalent proves eventual positivity and unboundedness. This argument compares two error bounds; it never differentiates an unknown integer-index error.
\end{proof}

\begin{theorem}[Shrinking threshold uncertainty]\label{xpc:sl:thm:inverse}
For every fixed $M\ge1$ there are constants $C_M>0$ and $y_M$ such that, for every real $y\ge y_M$,
\begin{equation}\label{xpc:sl:eq:ceilings}
 \ceil{x_M(y)-C_Mw_0^{-M}}
       \le\nu_X(y)\le
 \ceil{x_M(y)+C_Mw_0^{-M}}.
\end{equation}
The constants may depend on the fixed $k,b,X,M$. They are not computed effectively here.
\end{theorem}
\begin{proof}
Let $\widehat x_M$ be the root from Lemma~\ref{xpc:sl:lem:inverse-error}, and put $h=Cw_0^{-M}$ with $C>0$ to be chosen. Set
\[
 n_- =\ceil{\widehat x_M-h}-1,
 \qquad n_+=\ceil{\widehat x_M+h}.
\]
The first is strictly below $\widehat x_M-h$, and the second is at least $\widehat x_M+h$. Both are within a bounded distance of $\widehat x_M$, since $h\to0$. Their corresponding $w$ values are $w_0+O(w_0^{-1})$.

By \eqref{xpc:sl:eq:log-deriv}, the logarithmic change of $\mathcal F_M$ between its root and either test integer has the required sign and magnitude at least a fixed positive constant times $h/w_0=Cw_0^{-M-1}$. The forward theorem at order $M$ gives
\[
 X_{k,b}(n)=\mathcal F_M(n)\bigl(1+O(w_0^{-M-1})\bigr)
\]
uniformly at these two nearby integers; the factor in \eqref{xpc:sl:eq:smoothM} is bounded away from zero there. Choose $C$ large enough to dominate the two relative errors. Then
\[
 X_{k,b}(n_-)<y\le X_{k,b}(n_+)
\]
for large $y$. Eventual monotonicity excludes every integer at most $n_-$ in the eventual increasing segment and includes $n_+$. The finite initial segment is below $y$ once $y$ exceeds its maximum. Thus
\[
 \ceil{\widehat x_M-h}\le\nu_X(y)\le\ceil{\widehat x_M+h}.
\]
Finally \eqref{xpc:sl:eq:smooth-error} lets us enlarge the constant and replace $\widehat x_M$ by $x_M$. This proves \eqref{xpc:sl:eq:ceilings} with the exact weak threshold convention \eqref{xpc:sl:eq:nu}.
\end{proof}
For sufficiently large $y$ the interval in \eqref{xpc:sl:eq:ceilings} has width less than one, so the two ceilings are equal or adjacent. When they agree they identify the threshold exactly. When $x_M$ lies too close to an integer, an unconditional assertion $\nu_X(y)=\lceil x_M(y)\rceil$ is not justified. The theorem states an existential asymptotic bracket, not a certified finite-$y$ enclosure. In particular numerical errors in a few examples cannot be used to select a valid universal $C_M$.

\begin{remark}[{[write]} The inverse and the transseries volume, 7~October 2026]\label{xpc:sl:rem:transseries}
Statement by statement, against the volume named in the Guide (its symbols
carry the subscript ``vol'').
\begin{enumerate}[label=(\alph*),leftmargin=*]
\item \emph{Instance.} By Lemma~\ref{xpc:sl:lem:monotone}, for $y$ beyond the
finite initial segment $\nu_X(y)$ of~\eqref{xpc:sl:eq:nu} is the integer
staircase $N_*(y)$ of \file{p0:def:three-inverses}, so
\file{p0:thm:staircase}(1) applies. In the counts computed by the write,
$E_{2,0}$ increases strictly from $n=15$ and $T_{2,1}$ from $n=5$ through
$n=3000$; that is a finite observation, not a proof of the onset.
\item \emph{Instance.} The Lambert root $w_0$ of~\eqref{xpc:sl:eq:w0} is
\file{p0:thm:lambert-core}(3), case $b_\vol<0$, with $X_\vol=w_0$,
$a_\vol=1$, $b_\vol=-d$, $L_\vol=\log(y/D)$: the equation
\eqref{xpc:sl:eq:leading-inverse} is $w_0-d\log w_0=L_\vol$, and the large
solution $(|b_\vol|/a_\vol)(-W_{-1}(-(a_\vol/|b_\vol|)e^{-L_\vol/|b_\vol|}))$
is~\eqref{xpc:sl:eq:w0}; the condition $y>D(e/d)^d$ is the fold condition
$L_\vol>L_c$.
\item \emph{Formal instance.} The reversion equation~\eqref{xpc:sl:eq:reversion}
is the formal-exactness identity of \file{p0:thm:lambert-centered}(4) with
$a_\vol=1$, $b_\vol=-d$, $Q_\vol(\tau)=\log H(2A\tau)$, $\tau=1/w_0$ and
$U=\Delta$: with $\mathsf x=w$, $a(\mathsf x-\tau^{-1})+b\log(\tau\mathsf x)+
Q(1/\mathsf x)=\Delta-d\log(1+z\Delta)+\log H(2Az/(1+z\Delta))$. So
Proposition~\ref{xpc:sl:prop:inverse-coefs} is an instance, $u_m$ is the
volume's $c_m$, and the volume's Bell recurrence~\file{p0:eq:c-bell}
reproduces~\eqref{xpc:sl:eq:u1}--\eqref{xpc:sl:eq:u3} (checked by solving the
reversion through degree three).
\item \emph{Instance.} Put $Q_M(\tau)=\log(1+\sum_{j\le M}c_j(2A\tau)^j)$,
whose first $M$ coefficients are those of the full series. In the variable
$w$, the function $\mathcal F_M$ of~\eqref{xpc:sl:eq:smoothM} satisfies
\[
\log\mathcal F_M=\log D+w-d\log w+Q_M(1/w)
\]
exactly, so it is of exponential--power type (\file{p0:def:model}) with data
$(D,1,1,-d,1,Q_M)$. Hence the root
estimate behind~\eqref{xpc:sl:eq:smooth-error} is
\file{p0:thm:lambert-centered}(5) in $w$, followed by $x=\tfrac1{24}+w^2/(4A)$;
its residual-over-slope step is \file{p0:thm:backward-error}.
\item \emph{Analogue.} Theorem~\ref{xpc:sl:thm:inverse} is, as stated and
proved, an analogue of \file{p0:thm:staircase}(2): a two-ceiling bracket
proved directly at the integers, not deduced from an admissible
interpolation.
\end{enumerate}
No novelty is claimed for any of these inversions.
\end{remark}

\section{Exact checks and numerical diagnostics}\label{xpc:sl:sec:checks}
The companion package separates exact arithmetic checks from floating-point illustrations. The former verify finite identities and coefficient implementations; the analytic remainder and threshold theorems rest on the proofs above.

\subsection{Independent finite paths}
The public exact core uses Python's standard library, including arbitrary-precision integers and rational fractions. Partition numbers can be generated by Euler's pentagonal recurrence. For exact statistic values, each sieve summand is formed as a finite product and convolved with that partition sequence. It retains all $j$ for which $D_j\le n_{\max}$, so no asymptotic approximation enters these counts.

A separate bounded-part, exact-length dynamic program supplies an independent counting path. Before adding part size $M$, its entries count partitions with all parts less than $M$. The ordinary unbounded coin-change update adds partitions whose maximum is exactly $M$, keeping the length coordinate. Each newly added contribution is assigned to every selected relation $M=kL+b$ or $M\ge kL+b$. Summing over all maxima and lengths recovers an independently computed unrestricted partition sequence. This checks the sieve without reusing its inclusion--exclusion identity.

For coefficient checks, one path multiplies truncated exponential series with rational coefficients. A different path first constructs the finite polynomial $H(q)=\sum_dh_dq^d$, then computes moments $\sum_dh_d(-d)^\ell/\ell!$. The transfer is also checked by comparing the terminating factorial formula with the differential recurrence \eqref{xpc:sl:eq:Qrec}. The inverse generator verifies its formal residual to the requested order. All guards are explicit exceptions and remain active with Python's \texttt{-O} flag.

The saved exact checks compare all 98 parameter and kind combinations with $2\le k\le8$, $0\le b\le6$, and $0\le n\le160$, totaling 15,778 integer comparisons. They check 72 moment, transfer, and inverse cases through relative order six, the $Q_\ell$ recurrence through $\ell=20$, and 690 sufficient-depth combinations. There are 23 invalid-input and poisoned-receipt rejection controls, with an additional 26 package-integrity rejection controls. These figures describe finite checks, not the scope of the theorem. The independent algorithms were adapted from separate research-stage calculations; no claim of clean-room authorship or independent proof is made.

\subsection{How to read the diagnostics}
For a forward truncation order $R$, define
\[
 \mathcal A_R(n)=k!B(n)t^s\sum_{m=0}^Rc_mt^m,
 \qquad \operatorname{relerr}_R(n)=\frac{\mathcal A_R(n)}{X_{k,b}(n)}-1.
\]
For inverse order $M$, the diagnostic is $x_M(X_{k,b}(n))-n$, using the real center \eqref{xpc:sl:eq:xM}. Order zero uses the Lambert-only center. These are signed errors, not rigorous enclosures. The integer target is generated exactly; the subsequent transcendental evaluations use optional high-precision arithmetic. The diagnostic code records the precision and library version. At each of the 16 displayed targets $y=X_{k,b}(n)$, a scan of every exact count from zero through $n$, together with adjacent comparisons, verifies $\nu_X(y)=n$. This is a finite verification of these particular thresholds, independent of the asymptotic inverse errors.

% [write] inlined from the delivered tables/forward_table.tex (shipped as data/)
% Generated by code/reproduce.py --diagnostics; do not edit.
\begin{table}[htbp]
\centering
\small
\setlength{\tabcolsep}{4pt}
\begin{tabular}{ccr|rrrrr}
\hline
$X$ & $(k,b)$ & $n$ & $0$ & $1$ & $2$ & $3$ & $6$ \\
\hline
$E$ & $(2,0)$ & 1000 & $6.262\!\times\!10^{-1}$ & $-5.164\!\times\!10^{-1}$ & $5.452\!\times\!10^{-1}$ & $-7.057\!\times\!10^{-1}$ & $3.930\!\times\!10^{0}$ \\
$E$ & $(2,0)$ & 10000 & $2.105\!\times\!10^{-1}$ & $-5.842\!\times\!10^{-2}$ & $2.060\!\times\!10^{-2}$ & $-8.845\!\times\!10^{-3}$ & $1.741\!\times\!10^{-3}$ \\
$T$ & $(2,0)$ & 1000 & $3.487\!\times\!10^{-1}$ & $-1.935\!\times\!10^{-1}$ & $1.534\!\times\!10^{-1}$ & $-1.586\!\times\!10^{-1}$ & $5.469\!\times\!10^{-1}$ \\
$T$ & $(2,0)$ & 10000 & $1.196\!\times\!10^{-1}$ & $-2.271\!\times\!10^{-2}$ & $6.085\!\times\!10^{-3}$ & $-2.105\!\times\!10^{-3}$ & $2.612\!\times\!10^{-4}$ \\
$E$ & $(2,1)$ & 1000 & $6.749\!\times\!10^{-1}$ & $-5.699\!\times\!10^{-1}$ & $6.091\!\times\!10^{-1}$ & $-7.942\!\times\!10^{-1}$ & $4.471\!\times\!10^{0}$ \\
$E$ & $(2,1)$ & 10000 & $2.242\!\times\!10^{-1}$ & $-6.352\!\times\!10^{-2}$ & $2.264\!\times\!10^{-2}$ & $-9.787\!\times\!10^{-3}$ & $1.946\!\times\!10^{-3}$ \\
$T$ & $(2,1)$ & 1000 & $3.956\!\times\!10^{-1}$ & $-2.220\!\times\!10^{-1}$ & $1.767\!\times\!10^{-1}$ & $-1.831\!\times\!10^{-1}$ & $6.333\!\times\!10^{-1}$ \\
$T$ & $(2,1)$ & 10000 & $1.331\!\times\!10^{-1}$ & $-2.551\!\times\!10^{-2}$ & $6.859\!\times\!10^{-3}$ & $-2.377\!\times\!10^{-3}$ & $2.958\!\times\!10^{-4}$ \\
$E$ & $(3,0)$ & 1000 & $5.491\!\times\!10^{-1}$ & $-1.885\!\times\!10^{-2}$ & $-1.970\!\times\!10^{-1}$ & $1.448\!\times\!10^{-1}$ & $1.235\!\times\!10^{-1}$ \\
$E$ & $(3,0)$ & 10000 & $1.382\!\times\!10^{-1}$ & $6.249\!\times\!10^{-3}$ & $-6.839\!\times\!10^{-3}$ & $1.102\!\times\!10^{-3}$ & $2.349\!\times\!10^{-5}$ \\
$T$ & $(3,0)$ & 1000 & $3.925\!\times\!10^{-1}$ & $-2.110\!\times\!10^{-2}$ & $-8.118\!\times\!10^{-2}$ & $5.522\!\times\!10^{-2}$ & $3.361\!\times\!10^{-2}$ \\
$T$ & $(3,0)$ & 10000 & $1.057\!\times\!10^{-1}$ & $1.812\!\times\!10^{-3}$ & $-2.958\!\times\!10^{-3}$ & $4.666\!\times\!10^{-4}$ & $7.834\!\times\!10^{-6}$ \\
$E$ & $(2,3)$ & 1000 & $7.783\!\times\!10^{-1}$ & $-6.876\!\times\!10^{-1}$ & $7.544\!\times\!10^{-1}$ & $-9.998\!\times\!10^{-1}$ & $5.767\!\times\!10^{0}$ \\
$E$ & $(2,3)$ & 10000 & $2.520\!\times\!10^{-1}$ & $-7.437\!\times\!10^{-2}$ & $2.715\!\times\!10^{-2}$ & $-1.191\!\times\!10^{-2}$ & $2.421\!\times\!10^{-3}$ \\
$T$ & $(2,3)$ & 1000 & $4.952\!\times\!10^{-1}$ & $-2.878\!\times\!10^{-1}$ & $2.320\!\times\!10^{-1}$ & $-2.420\!\times\!10^{-1}$ & $8.454\!\times\!10^{-1}$ \\
$T$ & $(2,3)$ & 10000 & $1.605\!\times\!10^{-1}$ & $-3.172\!\times\!10^{-2}$ & $8.625\!\times\!10^{-3}$ & $-3.009\!\times\!10^{-3}$ & $3.777\!\times\!10^{-4}$ \\
\hline
\end{tabular}
\caption{Forward relative error $\mathcal A_R(n)/X(n)-1$ at fixed truncation orders. These 80-digit mpmath evaluations are non-certified diagnostics, not error bounds.}
\label{xpc:sl:tab:forward-diagnostics}
\end{table}
% [write] inlined from the delivered tables/inverse_table.tex (shipped as data/)
% Generated by code/reproduce.py --diagnostics; do not edit.
\begin{table}[htbp]
\centering
\small
\setlength{\tabcolsep}{4pt}
\begin{tabular}{ccr|rrrrr}
\hline
$X$ & $(k,b)$ & $n$ & $0$ & $1$ & $2$ & $3$ & $6$ \\
\hline
$E$ & $(2,0)$ & 1000 & $-1.257\!\times\!10^{1}$ & $4.826\!\times\!10^{0}$ & $-4.446\!\times\!10^{0}$ & $5.653\!\times\!10^{0}$ & $-3.644\!\times\!10^{1}$ \\
$E$ & $(2,0)$ & 10000 & $-1.513\!\times\!10^{1}$ & $2.203\!\times\!10^{0}$ & $-6.962\!\times\!10^{-1}$ & $2.966\!\times\!10^{-1}$ & $-6.788\!\times\!10^{-2}$ \\
$T$ & $(2,0)$ & 1000 & $-7.645\!\times\!10^{0}$ & $2.292\!\times\!10^{0}$ & $-1.722\!\times\!10^{0}$ & $1.810\!\times\!10^{0}$ & $-7.266\!\times\!10^{0}$ \\
$T$ & $(2,0)$ & 10000 & $-8.912\!\times\!10^{0}$ & $1.002\!\times\!10^{0}$ & $-2.582\!\times\!10^{-1}$ & $9.136\!\times\!10^{-2}$ & $-1.312\!\times\!10^{-2}$ \\
$E$ & $(2,1)$ & 1000 & $-1.333\!\times\!10^{1}$ & $5.075\!\times\!10^{0}$ & $-4.698\!\times\!10^{0}$ & $6.008\!\times\!10^{0}$ & $-3.932\!\times\!10^{1}$ \\
$E$ & $(2,1)$ & 10000 & $-1.601\!\times\!10^{1}$ & $2.318\!\times\!10^{0}$ & $-7.361\!\times\!10^{-1}$ & $3.153\!\times\!10^{-1}$ & $-7.324\!\times\!10^{-2}$ \\
$T$ & $(2,1)$ & 1000 & $-8.517\!\times\!10^{0}$ & $2.425\!\times\!10^{0}$ & $-1.838\!\times\!10^{0}$ & $1.945\!\times\!10^{0}$ & $-7.931\!\times\!10^{0}$ \\
$T$ & $(2,1)$ & 10000 & $-9.853\!\times\!10^{0}$ & $1.062\!\times\!10^{0}$ & $-2.758\!\times\!10^{-1}$ & $9.819\!\times\!10^{-2}$ & $-1.431\!\times\!10^{-2}$ \\
$E$ & $(3,0)$ & 1000 & $-1.147\!\times\!10^{1}$ & $-2.409\!\times\!10^{0}$ & $2.700\!\times\!10^{0}$ & $-1.092\!\times\!10^{0}$ & $-9.135\!\times\!10^{-1}$ \\
$E$ & $(3,0)$ & 10000 & $-1.029\!\times\!10^{1}$ & $-1.250\!\times\!10^{0}$ & $3.482\!\times\!10^{-1}$ & $-2.479\!\times\!10^{-2}$ & $-2.171\!\times\!10^{-4}$ \\
$T$ & $(3,0)$ & 1000 & $-8.570\!\times\!10^{0}$ & $-1.233\!\times\!10^{0}$ & $1.301\!\times\!10^{0}$ & $-5.087\!\times\!10^{-1}$ & $-3.261\!\times\!10^{-1}$ \\
$T$ & $(3,0)$ & 10000 & $-7.955\!\times\!10^{0}$ & $-6.295\!\times\!10^{-1}$ & $1.657\!\times\!10^{-1}$ & $-1.317\!\times\!10^{-2}$ & $-1.767\!\times\!10^{-4}$ \\
$E$ & $(2,3)$ & 1000 & $-1.488\!\times\!10^{1}$ & $5.553\!\times\!10^{0}$ & $-5.223\!\times\!10^{0}$ & $6.762\!\times\!10^{0}$ & $-4.564\!\times\!10^{1}$ \\
$E$ & $(2,3)$ & 10000 & $-1.779\!\times\!10^{1}$ & $2.543\!\times\!10^{0}$ & $-8.195\!\times\!10^{-1}$ & $3.550\!\times\!10^{-1}$ & $-8.499\!\times\!10^{-2}$ \\
$T$ & $(2,3)$ & 1000 & $-1.027\!\times\!10^{1}$ & $2.680\!\times\!10^{0}$ & $-2.082\!\times\!10^{0}$ & $2.235\!\times\!10^{0}$ & $-9.410\!\times\!10^{0}$ \\
$T$ & $(2,3)$ & 10000 & $-1.174\!\times\!10^{1}$ & $1.179\!\times\!10^{0}$ & $-3.129\!\times\!10^{-1}$ & $1.129\!\times\!10^{-1}$ & $-1.693\!\times\!10^{-2}$ \\
\hline
\end{tabular}
\caption{Inverse real-index error $x_M(X(n))-n$; order zero is the Lambert core. These 80-digit mpmath evaluations are non-certified diagnostics, not error bounds.}
\label{xpc:sl:tab:inverse-diagnostics}
\end{table}

Large alternating coefficients make these fixed-order expansions numerically demanding. At a moderate index, increasing the truncation order may make an approximation worse. Such behavior does not conflict with Theorem~\ref{xpc:sl:thm:forward}, which fixes the order before $n$ tends to infinity, and does not establish an optimal truncation or resummation rule. The package includes fixed-offset and slope-$3$ companions to check that the implementation is not specialized only to the two named sequences.

\section{Relation to prior partition asymptotics}\label{xpc:sl:sec:prior}
This section records a bounded comparison of the cited source statements, their parameters, and their errors. It is not an exhaustive historical survey or a priority certificate. The principal analytic inputs in the proof are classical eta inversion and the exact Rademacher series. The model-specific argument is the finite-product sieve with its global finite-depth remainder and its explicit fixed-slope coefficients and inverse.

\subsection{Joint extreme counts and rectangles}
Szekeres's \emph{Asymptotic distribution of partitions by number and size of parts}~\cite{szekeres} uses $p(n;r,k)$ for the \emph{exact} joint count with length $r$ and largest part $k$. Its Theorem~1 sets $c=\pi/\sqrt6$, $\lambda=ck/\sqrt n$, and $\mu=cr/\sqrt n$, and assumes
\[
 \frac14\log n<\lambda,\mu<\frac{n^{1/4}}{\log n}.
\]
The resulting leading equivalent contains correlation terms in the exponent. A corollary on printed page 530 gives an independent-extreme simplification under narrower assumptions. The proof on printed pages 532--537 keeps an $o(1)$ exponential remainder. Thus both the broad correlated theorem and the narrower independent-extreme law are relevant prior work. The original exact-count definition matters; a later paper's cumulative notation should not be substituted for it. These leading joint asymptotics do not themselves provide the arbitrary fixed-order diagonal sum and offset calculus proved here.

Jiang and Wang~\cite{jiangwang} give a generalized Hardy--Ramanujan formula for cumulative rectangular restrictions. Their Theorem~1 and Remarks~1.2--1.3 allow both cutoffs to be at least $4\sqrt n$, with unbounded normalized cutoffs, and have relative error $O(n^{-1/4+\varepsilon})$. This range overlaps the $\sqrt n\log n$ cutoff sizes relevant to fixed-slope extreme events. It would be incorrect to dismiss their theorem as confined to bounded multiples of $\sqrt n$. Its stated error, however, does not justify taking the rare diagonal differences and sums to every algebraic order.

Melczer, Panova, and Pemantle~\cite{mpp} treat partitions inside a rectangle. Their Theorems~1--2 require the normalized parameters $(\ell/m,n/m^2)$ to remain in a compact subset of $x\ge2y>0$. For $m$ of order $\sqrt n\log n$, the second coordinate tends to zero, outside that compact regime. Their refined consecutive-$n$ differences also hold the rectangle fixed. These are important rectangle results, but their particular compact-uniform theorem does not directly cover the full dominant fixed-slope sum.

Richmond~\cite{richmond} treats central extreme coordinates and extended one-sided cumulative ranges. The central window $|x_i|\le(\log n)/5$ has a displayed error $O(n^{-1/10}\log n)$. Equation~(10) also treats exact joint counts, so the work should not be described as cumulative-only. It includes successive-rank questions. Its stated leading-error framework does not supply the all-fixed-order slope-coupled expansion or inverse threshold result here.

\subsection{Sieve and transfer antecedents}
Hwang~\cite{hwang} develops an analytic sieve for partitions with many summands. Theorem~1 compares $p(n,m)$ with $p(n-m)$ in a large-length regime, with a stated relative power error; Proposition~1 gives a general finite sieve with a controlled remainder. Later theorems extend to restricted part sets and length laws. This is a substantive methodological antecedent. It does not state the summed $M=kL+b$ or $M\ge kL+b$ expansion and inverse treated here.

Ngo and Rhoades~\cite{ngorhoades} establish arbitrary-order expansions for partition statistics in a quantum-modular framework, with general amplitude-transfer results and polynomial and logarithmic variants. Therefore a claim that partition-product amplitudes admit arbitrary-order expansions would be too broad as a novelty claim. Our elementary finite-shift transfer makes the particular coefficient normalization transparent, while the substantive analytic obligation is the global remainder for this fixed-slope amplitude.

Johansson~\cite{johansson} provides the exact Rademacher normalization used in Lemma~\ref{xpc:sl:lem:gap}, including a truncation estimate that tends to zero. We proved the one-term exponential gap from that source's convergent formula, rather than merely quoting a leading equivalent. The classical original is Rademacher~\cite{rademacher}; the directly inspected and used formula is Johansson's. Eta inversion is taken in the normalization recorded by the DLMF~\cite{dlmf}.

\subsection{Rank terminology and the excluded slope}
At $k=1$, $M-L$ is Dyson's ordinary rank. Dyson's exact alternating identities and subsequent rank literature are prior work~\cite{dyson}. The finite-depth gain in Proposition~\ref{xpc:sl:prop:global} is $k-1$, which vanishes at that slope. No $k=1$ asymptotic conclusion is claimed here.

The named $k$-rank discussed by Berkovich and Garvan~\cite[equations (6.1) and (6.2)]{berkovichgarvan} is defined using successive Durfee squares, selected columns, and rows. It is not the statistic $M-kL$. We use the descriptive phrase ``fixed-slope extreme statistic'' to avoid identifying these different notions. A finite search of rank terminology cannot exclude an unindexed treatment of the same diagonal event elsewhere.

\section{Reproducibility and scope of the package}\label{xpc:sl:sec:package}
The compact archive contains the editable \LaTeX{} manuscript, its PDF, standard-library exact code, selected generated rational coefficients and exact counts, optional numerical diagnostics, the generated tables, source references, and a checksum manifest. It does not redistribute source papers. No third-party source PDF is required to rerun the exact finite checks, and no network access is required for regeneration.

The exact core and its tests remain usable without the optional numerical library. The full PDF-and-archive build additionally needs the documented \TeX{} toolchain and the optional library for the diagnostic tables. The entry-point README gives commands for a fresh output directory, normal and optimized execution, byte-level verification of regenerated artifacts, and package guard tests. Same-toolchain byte stability is a reproducibility property of these files; it is not a claim that different \TeX{}, font, Python, or numerical-library versions necessarily emit identical bytes.

Finite tests cover only the ranges explicitly recorded in their receipts. Agreement of independently generated finite counts supports the implementation and the exact identities, but cannot prove the asymptotic remainder. Numerical diagnostics support neither effective constants in \eqref{xpc:sl:eq:main} nor a certified finite-target enclosure in \eqref{xpc:sl:eq:ceilings}. The latter remain theorem statements with existential constants.

\section{Further directions}\label{xpc:sl:sec:directions}
Several extensions would require additional arguments rather than a change of notation.
\begin{enumerate}
\item An effective version of the global bounds, the Rademacher tail estimate, and the finite-shift remainder could in principle produce explicit constants and an onset for a chosen $k,b,R$. That work would turn the inverse bracket into a computable enclosure. It is not supplied by the present diagnostic tables.
\item Allowing $b$ or $k$ to vary with $n$ changes the uniformity of the numerator bounds and the finite polynomial retained in the transfer. In particular a growing shift is no longer a bounded Taylor shift. The fixed-parameter theorem cannot simply be reused.
\item The coefficients grow rapidly in the slope-$2$ examples. Understanding optimal truncation, large-order coefficient behavior, and exponentially small corrections would be a distinct problem. No convergent resummation follows from the all-fixed-order expansion.
\item The boundary $k=1$ belongs to the ordinary rank problem, with no positive sieve-valuation gain. A uniform transition towards that regime would need a different analytic organization and a careful comparison with existing rank asymptotics.
\item Similar finite-product identities may arise for other coupled diagram boundaries. Each model would need its own entire-circle majorant; a formal radial series is not a substitute for that bound.
\end{enumerate}
\begin{writenote}[7 October 2026]
Under Vladimir's standing rule of 4~October 2026 (unproved claims are moved
to questions, wrong ones are refuted on the record), the write found no claim
of this Part that is unproved as stated and none that is wrong. The one wrong
claim met in its sources is the OEIS alternative form discussed in
Section~\ref{xpc:sl:sec:result}; it is refuted, with proof, in
Remark~\ref{xpc:sl:rem:a350879}. On Item~1, note only that the write's exact
counts increase strictly from $n=15$ ($E_{2,0}$) and $n=5$ ($T_{2,1}$) through
$n=3000$, which proves no onset.
\end{writenote}

% [write] Part I's appendix keeps its letter; the section counter is restored at Part II.
\setcounter{xpcsavesec}{\value{section}}\setcounter{section}{0}\renewcommand{\thesection}{\Alph{section}}\renewcommand{\theHsection}{xpcslapp.\Alph{section}}
\section{Exact algorithms in finite arithmetic}\label{xpc:sl:app:algorithms}
This appendix records the finite calculations independently of any particular programming language.

\paragraph{Exact sequence values.}
To compute $X_{k,b}(0),\ldots,X_{k,b}(n_{\max})$, first compute $p(0),\ldots,p(n_{\max})$. For each $j$ with $D_j\le n_{\max}$, start with the truncated coefficient list of $P(q)$. Multiply it successively by $1-q^r$ for $r=j,\ldots,kj$ in the equality case, or $r=j+1,\ldots,kj$ for the tail. For a single factor, update coefficients in descending index order by $a_m\leftarrow a_m-a_{m-r}$. Shift by $D_j$, multiply by $(-1)^{j-1}$, and add to the output. All arithmetic is integral. Descending order is essential: ascending order would implement a different operation.

\paragraph{Rational forward coefficients.}
For the requested relative order $R$, set $L=s+R$. For each $j$ allowed by \eqref{xpc:sl:eq:Jformal}, multiply the truncated series
\[
 \ee^{-D_jz}=\sum_{a=0}^{L}\frac{(-D_j)^a}{a!}z^a+O(z^{L+1}),
 \qquad
 1-\ee^{-rz}=-\sum_{a=1}^{L}\frac{(-r)^a}{a!}z^a+O(z^{L+1}).
\]
Add the signed product to obtain $g_0,\ldots,g_L$. Then use \eqref{xpc:sl:eq:cfinite}. Every operation is rational, and $A^{-1}$ is an indeterminate until numerical evaluation is explicitly requested. This exact construction also verifies $g_0=\cdots=g_{s-1}=0$ and $g_s=k!$.

\paragraph{Independent coefficient moments.}
Build the finite sieve polynomial $H(q)$ in integer arithmetic instead. If $H(q)=\sum_dh_dq^d$, then
\[
 g_\ell=\frac1{\ell!}\sum_dh_d(-d)^\ell\qquad(0\le\ell\le L).
\]
Compute $Q_\ell$ by \eqref{xpc:sl:eq:Qrec}, rather than by its factorial formula. This path compares two different finite constructions of both the radial and transfer coefficients.

\paragraph{Triangular inverse.}
Suppose $u_1,\ldots,u_{m-1}$ have been found. Set $u_m=0$, substitute the truncated $\Delta$ into the left side of \eqref{xpc:sl:eq:reversion}, and extract its degree-$m$ coefficient $r_m$. Then set $u_m=-r_m$. Formal multiplication, reciprocal, and logarithm are all truncated at degree $m$. The inverse coefficients can be stored as rational Laurent polynomials in $A$, since the forward coefficients are rational polynomials in $A^{-1}$. The residual through the requested degree must vanish exactly. A negative control that changes one coefficient must fail that residual test.


\part{Exact ratio two partitions at fixed asymptotic order}\label{xpc:r2:part}
\setcounter{section}{\value{xpcsavesec}}\renewcommand{\thesection}{\arabic{section}}\renewcommand{\theHsection}{\arabic{section}}
\setlength{\parindent}{\xpcdefparindent}\setlength{\parskip}{0pt plus 1pt}\setlist{nosep}
\partsource{Bundle Report~217, archive \file{Report217-source-v2.zip}, delivered as \file{article.tex} (14 pages). Delivered title block ``Exact ratio two partitions at fixed asymptotic order / Report 217 / 4 October 2026''. Printed as delivered, with \textbf{[write]} remarks and notes; Report~217's Section~$k$ is Section~$k+10$ here and its equation~($k$) is equation~($k+65$); statements move with their sections. Its labels carry \file{xpc:r2:}; its three unlabelled sections were labelled; its citation keys \file{oeis} and \file{dlmf} are \file{oeisA118096} and \file{dlmfBessel}. This Part is self-contained.}
\begin{partabstract}
Let $a(n)$ count partitions of $n$ whose largest part is exactly twice their smallest part, the sequence OEIS A118096. We give a self-contained fixed-order coefficient expansion, an exact finite rule for every correction, the first four corrections, and rounding-aware inverse bounds. A complete-circle estimate suppresses every noncentral angle, and a discrete two-variable saddle provides a uniform finite-order remainder. This is an independently derived refinement, not a claim of novelty or a solution of a known open problem. The leading equivalent was already recorded by Kot\v e\v sovec in 2025; general radial all-order methods precede this account. Historical comparison is incomplete: Pittel's directly adjacent 2007 bounded-ratio paper was identified, but its full theorem text was not accessible for this comparison. That limitation is not evidence of absence of overlapping results.
\end{partabstract}

\section{Scope and main results}\label{xpc:r2:sec:scope}
A partition has positive integer parts. For $n\geq0$, let $a(n)$ count those with largest part $2k$ and smallest part $k$ for some $k\geq1$; in particular $a(0)=0$. Requiring one part of each endpoint gives, for $|q|<1$,
\begin{equation}\label{xpc:r2:eq:F}
 F(q)=\sum_{n\geq0}a(n)q^n
 =\sum_{k\geq1}T_k(q),\qquad
 T_k(q)=\frac{q^{3k}}{\prod_{j=k}^{2k}(1-q^j)}.
\end{equation}
The sum converges locally absolutely in the disk. The endpoint copies account for $3k$; all additional parts lie in $[k,2k]$. Set
\begin{equation}\label{xpc:r2:eq:constants}
 \phi=\frac{1+\sqrt5}{2},\quad x_0=\log\phi,\quad
 A=\frac{\pi^2}{30},\quad B=2\sqrt A,\quad
 C=\frac1{\sqrt{5+\sqrt5}}.
\end{equation}

\begin{theorem}[Coefficient expansion]\label{xpc:r2:thm:coeff}
For every fixed integer $M\geq0$, as integer $n\to\infty$,
\begin{equation}\label{xpc:r2:eq:main}
 a(n)=C n^{-1/2}e^{B\sqrt n}
 \left(\sum_{r=0}^{M}d_r n^{-r/2}
       +O_M(n^{-(M+1)/2})\right),\qquad d_0=1.
\end{equation}
The constants are determined by the exact finite rule in Section~\ref{xpc:r2:sec:rule} and the transfer formula
\begin{equation}\label{xpc:r2:eq:transfer}
 d_r=\sum_{j=\lceil r/2\rceil}^{r}
 c_j A^{j/2}\,
 \frac{(-1)^{r-j}r!}
 {(r-j)!(2j-r)!(4\sqrt A)^{r-j}}.
\end{equation}
In particular, $c_0=1$ and
\begin{align}
 c_1&=\frac{239}{120}+\sqrt5,\notag\\
 c_2&=\frac{232801}{28800}+\frac{431\sqrt5}{120},\notag\\
 c_3&=\frac{640181519}{10368000}+\frac{795361\sqrt5}{28800},\label{xpc:r2:eq:cs}\\
 c_4&=\frac{3268837156801}{4976640000}
       +\frac{3045555791\sqrt5}{10368000}.\notag
\end{align}
Consequently
\begin{align}
 d_1&=\sqrt A\,c_1,&
 d_2&=Ac_2-\tfrac12c_1,\notag\\
 d_3&=A^{3/2}c_3-\tfrac32\sqrt A\,c_2,&
 d_4&=A^2c_4-3Ac_3+\tfrac34c_2.\label{xpc:r2:eq:ds}
\end{align}
Numerically these are $2.424917491063$, $3.187616630712$,
$9.439637377883$, and $32.379308690155$, respectively.
\end{theorem}

The $M=0$ equivalent is the one already recorded in A118096~\cite{oeisA118096}:
\[
 \frac{\exp(\pi\sqrt{2n/15})}{5^{1/4}\sqrt{2\phi n}}.
\]
All $O$-constants in this article may depend on the fixed requested order. Neither uniformity as that order grows nor an explicit numerical threshold is asserted. The coefficients are computable exactly; the asymptotic remainder constants are not certified by the accompanying numerical experiments.

The proof separates two issues that should not be conflated. The real-axis expansion belongs to the existing Eulerian-series framework of McIntosh~\cite{mcintosh} and Zhou~\cite{zhou}. It cannot by itself be transferred to all coefficient orders by a leading-term Tauberian argument. Sections~\ref{xpc:r2:sec:angle}--\ref{xpc:r2:sec:identify} give the angular estimate, the uniform two-variable remainder, and a finite Gaussian-contour identification needed here. Section~\ref{xpc:r2:sec:inverse} then gives a standard asymptotic inversion with two ceilings, and Section~\ref{xpc:r2:sec:shape} proves eventual strict increase and log-concavity without differentiating an uncontrolled remainder.

\begin{remark}[{[write]} The OEIS entry, quoted, 7~October 2026]\label{xpc:r2:rem:oeis}
OEIS A118096, read on 7~October 2026 in the internal format, is at revision
\#57 (18~March 2026), author Emeric Deutsch (12~April 2006), offset~1. Its
name is ``Number of partitions of n such that the largest part is twice the
smallest part.''; its formula lines are ``G.f.: Sum\_\{k>=1\} x\^{}(3*k)/Product\_\{j=k..2*k\}
(1-x\^{}j).'' (this Part's~\eqref{xpc:r2:eq:F}) and Kot\v{e}\v{s}ovec's ``a(n)
\textasciitilde{} exp(Pi*sqrt(2*n/15)) / (5\^{}(1/4)*sqrt(2*phi*n))'' (13~June
2025), which is the case $M=0$ of Theorem~\ref{xpc:r2:thm:coeff}
($5^{1/4}\sqrt{2\phi}=\sqrt{5+\sqrt5}$); a comment gives the conjugate reading
(the number of parts is twice the multiplicity of the largest part). It also
links a 2025 preprint of Stephen McKean, which this Part does not cite and the
write did not read. Its b-file runs to $n=10000$ (Kot\v{e}\v{s}ovec; terms to
1000 from Alois P.~Heinz). The bibliography entry~\cite{oeisA118096} prints the
name as a title in its own words. The write computed $a(n)$
from~\eqref{xpc:r2:eq:F} for $n\le3000$: every b-file term in that range
agrees, as do brute-force counts for $n\le40$. A118096 is cross-referenced
from A117086, whose partitions with $\max/\min=2$ are exactly these
(Part~III). Nothing was submitted to the OEIS.
\end{remark}

\section{The radial saddle}\label{xpc:r2:sec:radial}
Define, on the positive real axis,
\begin{equation}\label{xpc:r2:eq:hg}
\begin{gathered}
 f(x)=-\log(1-e^{-x}),\qquad
 h(x)=\int_x^{2x}f(u)\,du=\Li(e^{-x})-\Li(e^{-2x}),\\
 g(x)=\frac{e^{-3x}}{\sqrt{(1-e^{-x})(1-e^{-2x})}}.
\end{gathered}
\end{equation}
Their analytic continuations are taken only in a fixed small complex neighborhood of $x_0$, with the branches determined by their positive real values. Put
\begin{equation}\label{xpc:r2:eq:alambda}
 a=-h''(x_0)=3\phi-4>0,\qquad
 \lambda=g(x_0)\sqrt{\frac{2\pi}{a}}
 =\sqrt{\frac{2\pi}{\sqrt5\phi}}.
\end{equation}
The symbol $a$ without an argument denotes this positive saddle curvature; $a(n)$ always denotes the partition count.

\begin{lemma}\label{xpc:r2:lem:saddle}
The function $h$ has a unique maximum on $(0,\infty)$, at $x_0$. Its value is $A=\pi^2/30$, and
\begin{equation}\label{xpc:r2:eq:D}
 D:=2A-a x_0^2=\int_{x_0}^{2x_0}u^2 f''(u)\,du>0.
\end{equation}
\end{lemma}
\begin{proof}
With $y=e^{-x}$,
\[
 h'(x)=2f(2x)-f(x)=-\log((1-y)(1+y)^2).
\]
The equation $h'(x)=0$ reduces to $y^2+y=1$. The sign is positive before $x_0$ and negative after it. Also $h(x)\to0$ at both endpoints of $(0,\infty)$, and direct differentiation gives $h''(x_0)=4-3\phi$.

For completeness, put $r=\phi^{-1}$ and $L=\log\phi$, so $r+r^2=1$. The dilogarithm identities needed here are
\begin{align*}
 \Li(r)+\Li(r^2)&=\pi^2/6-2L^2,\\
 \Li(r)+\Li(-r)&=\tfrac12\Li(r^2),\\
 \Li(r^2)&=-\Li(-r)-\tfrac12L^2.
\end{align*}
They are respectively reflection, duplication, and the identity
$\Li(x/(1+x))=-\Li(-x)-\tfrac12\log^2(1+x)$ at $x=r$.
Each follows by differentiation from
$\Li'(z)=-\log(1-z)/z$, with the constants obtained at zero or from
$\Li(1)=\sum_{m\geq1}m^{-2}=\pi^2/6$.
Solving gives $\Li(r)=\pi^2/10-L^2$ and
$\Li(r^2)=\pi^2/15-L^2$, proving $h(x_0)=A$.
Finally, twice integrating by parts gives
\[
 \int_x^{2x}u^2f''(u)\,du
 =x^2h''(x)-2x h'(x)+2h(x).
\]
At the saddle this is $D$; positivity follows from $f''(u)=e^{-u}/(1-e^{-u})^2>0$.
\end{proof}

\begin{lemma}[Radial tails]\label{xpc:r2:lem:tails}
Fix $0<\varepsilon<x_0<R$, with $\varepsilon\leq1$, and put
$\eta=A-\max\{h(\varepsilon),h(R)\}>0$. Uniformly for $0<t\leq1$ and $-\pi\leq\theta\leq\pi$,
\begin{equation}\label{xpc:r2:eq:tails}
 \sum_{kt\notin[\varepsilon,R]}|T_k(e^{-t+\ii\theta})|
 =O_{\varepsilon,R}\bigl(t^{-2}e^{(A-\eta)/t}\bigr).
\end{equation}
Moreover $F(e^{-t})=O(t^{-2}e^{A/t})$.
\end{lemma}
\begin{proof}
Since $f$ is decreasing,
\[
 \sum_{j=k}^{2k}f(jt)\leq t^{-1}h(kt)+f(kt),\qquad
 T_k(e^{-t})\leq\frac{e^{h(kt)/t-3kt}}{1-e^{-kt}}.
\]
For $k<\varepsilon/t$ use $1-e^{-kt}\geq1-e^{-t}\geq t/2$ and at most $\varepsilon/t$ terms. For $k>R/t$ use $1-e^{-kt}\geq1-e^{-R}$ and sum the geometric factor $e^{-3kt}$, whose sum is $O(t^{-1})$. Monotonicity of $h$ on the two tails supplies the exponent gap. On the compact interval each summand is at most a constant times $e^{A/t}$ and there are $O(t^{-1})$ terms. Finally $|1-q^j|\geq1-|q|^j$ implies $|T_k(q)|\leq T_k(|q|)$.
\end{proof}

\section{A complete circle estimate}\label{xpc:r2:sec:angle}
The following elementary bound controls rational and irrational angles at once.
\begin{lemma}\label{xpc:r2:lem:cos}
For an integer $L\geq2$, an integer $k$, and $|\theta|\leq\pi$,
\begin{equation}\label{xpc:r2:eq:cos}
 \sum_{j=k}^{k+L-1}(1-\cos(j\theta))
 \geq \frac{L}{81\pi^2}\min\{L^2\theta^2,1\}.
\end{equation}
\end{lemma}
\begin{proof}
Let $U=\sum_{j=0}^{L-1}e^{\ii j\theta}$. The left side is at least $L-|U|$, independently of $k$. For $|\theta|\leq2\pi/L$ set $m=\lfloor L/2\rfloor$. The identity
\[
 L^2-|U|^2=2\sum_{d=1}^{L-1}(L-d)(1-\cos(d\theta))
\]
and $1-\cos u\geq2u^2/\pi^2$ for $|u|\leq\pi$ imply
\[
 L-|U|\geq\frac{2\theta^2}{\pi^2 L}
 \sum_{d=1}^{m}(L-d)d^2
 \geq\frac{L^3\theta^2}{81\pi^2}.
\]
Here $L-d\geq L/2$, $m\geq L/3$, and
$\sum_{d=1}^m d^2\geq m^3/3$ suffice. If $|\theta|\geq2\pi/L$, the geometric sum gives
$|U|\leq1/|\sin(\theta/2)|\leq\pi/|\theta|\leq L/2$.
These two cases imply the stated, slightly weaker, bound.
\end{proof}

For $\varepsilon\leq kt\leq R$, write $r_j=e^{-jt}$. Exactly,
\[
 \frac{|T_k(e^{-t+\ii\theta})|}{T_k(e^{-t})}
 =\prod_{j=k}^{2k}
 \left(1+\frac{2r_j(1-\cos(j\theta))}{(1-r_j)^2}\right)^{-1/2}.
\]
For $0\leq u\leq2$ and $r\geq e^{-2R}$, concavity of the logarithm on this interval gives
\[
 \tfrac12\log\left(1+\frac{2ru}{(1-r)^2}\right)
 \geq u\,\operatorname{arctanh}(r)
 \geq\gamma_R u,\qquad \gamma_R=\operatorname{arctanh}(e^{-2R})>0.
\]
Indeed the chord slope from $u=0$ to $u=2$ is
$\frac14\log(1+4r/(1-r)^2)=\operatorname{arctanh}(r)$.
Apply Lemma~\ref{xpc:r2:lem:cos} with $L=k+1\geq\varepsilon/t$. Since $\varepsilon\leq1$,
\begin{equation}\label{xpc:r2:eq:angularT}
 |T_k(e^{-t+\ii\theta})|\leq T_k(e^{-t})
 \exp\left[-c\min\left\{\frac{\theta^2}{t^3},\frac1t\right\}\right],
 \qquad c=\frac{\gamma_R\varepsilon^3}{81\pi^2}>0.
\end{equation}
Summation and Lemma~\ref{xpc:r2:lem:tails} prove
\begin{equation}\label{xpc:r2:eq:angularF}
 |F(e^{-t+\ii\theta})|
 \leq F(e^{-t})e^{-c\min\{\theta^2/t^3,1/t\}}
       +O(t^{-2}e^{(A-\eta)/t}).
\end{equation}
Thus no other root of unity can contribute at algebraic order to the answer. This conclusion uses the entire circle, not merely a neighborhood of $q=1$.

\section{The discrete two variable saddle}\label{xpc:r2:sec:two}
Cauchy's formula gives
\begin{equation}\label{xpc:r2:eq:cauchy}
 a(n)=\frac1{2\pi}\int_{-\pi}^{\pi}
 \sum_{k\geq1}T_k(e^{-t+\ii\theta})e^{nt-\ii n\theta}\,d\theta.
\end{equation}
Choose $t=\sqrt{A/n}$ and $s=\sqrt t$. In the central region put
\begin{equation}\label{xpc:r2:eq:scaling}
 k=\frac{x_0}{t}+\frac y s,\qquad
 \theta=t^{3/2}v,\qquad
 z=t-\ii\theta=s^2(1-\ii s v),\qquad
 x=kz=(x_0+s y)(1-\ii s v).
\end{equation}
The possible $y$ form a lattice of mesh $s$, with a shift depending on $t$.

\subsection{Finite Euler Maclaurin expansion}
With $B_{2r}$ denoting Bernoulli numbers, define
\begin{equation}\label{xpc:r2:eq:e}
 e_{2r-1}(x)=\frac{B_{2r}}{(2r)!}
       \bigl(f^{(2r-1)}(2x)-f^{(2r-1)}(x)\bigr).
\end{equation}
Euler--Maclaurin applied to the finite sum $\sum_{j=k}^{2k}f(jz)$ gives, for every fixed $p\geq1$,
\begin{equation}\label{xpc:r2:eq:EM}
 T_k(e^{-z})=e^{h(x)/z}g(x)
 \exp\left(\sum_{r=1}^{p}e_{2r-1}(x)z^{2r-1}
                   +O(t^{2p-1})\right).
\end{equation}
The error estimate as written need not improve on the last displayed term; to obtain any specified finite power, apply this formula at a higher $p$ and discard unnecessary terms. Uniformity holds when $x$ and the line segment from $x$ to $2x$ remain in fixed compact analytic neighborhoods of their positive real counterparts and $|z|=O(t)$.

Here are the remainder details. Apply the real-variable Euler--Maclaurin identity to the complex-valued function $u\mapsto f(uz)$ on $[k,2k]$. Its $2p$th derivative is $z^{2p}f^{(2p)}(uz)$. The periodic Bernoulli remainder is bounded by a constant times
$k|z|^{2p}=O(t^{2p-1})$. The integral term is $h(kz)/z$ and the endpoint term is $(f(x)+f(2x))/2$. Combining the latter with $e^{-3x}$ gives $g(x)$. Bounded derivatives on the compact neighborhoods establish the asserted uniformity.

\subsection{Localization with explicit power bookkeeping}
For the fixed $\varepsilon,R$ above, set
\[
 c_h=\min_{x\in[\varepsilon,R]}
       \frac{A-h(x)}{(x-x_0)^2}>0,
\]
where the quotient at $x_0$ is defined as $a/2$. Let
$L_t=\sqrt{K\log(1/t)}$, with $K$ a sufficiently large fixed constant depending on the target order. The compact-$k$ part with $|y|>L_t$ is bounded radially by
$O(t^{-1}e^{A/t}t^{c_hK})$. In the full Cauchy integral, divided by its eventual leading scale $t e^{2A/t}$, this is $O(t^{c_hK-2})$.

For $|v|>L_t$, equation~\eqref{xpc:r2:eq:angularF} and
$F(e^{-t})=O(t^{-2}e^{A/t})$ give a relative bound
$O(t^{cK-3})$ plus an exponentially small term. For small enough $t$,
$K\log(1/t)<1/t$, as needed at the boundary of this region.
The noncompact $k$ tails are exponentially small by Lemma~\ref{xpc:r2:lem:tails}. Thus a relative remainder $O(t^N)$ permits restriction to
\begin{equation}\label{xpc:r2:eq:box}
 |y|,|v|\leq L_t
\end{equation}
by taking $c_hK\geq N+2$ and $cK\geq N+3$. We may increase $K$ further to cover the finitely many polynomial-Gaussian tails below. No choice here depends on $n$ apart from the stated logarithmic window.

\subsection{A uniform Gaussian remainder}
After multiplying the summand by $e^{nz}$ and extracting $g(x_0)e^{2A/t}$, its quadratic exponent is
\begin{equation}\label{xpc:r2:eq:Q}
 Q(y,v)=-\frac a2y^2+\ii a x_0yv-\frac D2v^2.
\end{equation}
For example, $A/z+nz-2A/t=-Av^2+O(sv^3)$, and the quadratic contribution of $h(x)-A$ is
$-a(y-\ii x_0v)^2/2$. This proves~\eqref{xpc:r2:eq:Q}, with the $D$ in~\eqref{xpc:r2:eq:D}. In particular $\Re Q$ is negative definite on the real $(y,v)$ plane.

\begin{lemma}[Uniform finite Taylor estimate]\label{xpc:r2:lem:remainder}
Fix integers $J\geq0$ and a fixed $K>0$. There are polynomials
$P_0,\ldots,P_J$ in $(y,v)$, with $P_0=1$, an integer $b_J$, and positive constants $C_J,t_J$ such that on~\eqref{xpc:r2:eq:box}, for $0<t<t_J$,
\begin{multline}\label{xpc:r2:eq:uniform}
 \left|\frac{T_k(e^{-z})e^{nz}}{g(x_0)e^{2A/t}}
       -e^{Q(y,v)}\sum_{m=0}^J P_m(y,v)s^m\right|\\
 \leq C_J s^{J+1}(1+|y|+|v|)^{b_J}
                 e^{-a y^2/4-Dv^2/4}.
\end{multline}
Furthermore $P_m(-y,-v)=(-1)^mP_m(y,v)$.
\end{lemma}
\begin{proof}
Use~\eqref{xpc:r2:eq:EM} at depth $p$ such that $4p-2\geq J+1$.
Substitute~\eqref{xpc:r2:eq:scaling} into its logarithm, subtract
$2A/s^2+\log g(x_0)+Q$, and denote the analytic part that remains by $R(s;y,v)$. The singular $s^{-2}$ and $s^{-1}$ terms cancel, so $R$ is analytic in $s$ at zero and $R(0;y,v)=0$. Each of its Taylor coefficients is a polynomial in $y,v$; the first has degree at most three.

Write $W=1+|y|+|v|$. For real $0\leq\sigma\leq s$, with $sW$ sufficiently small, the denominator $1-\ii\sigma v$ stays bounded away from zero and
$(x_0+\sigma y)(1-\ii\sigma v)$ stays in the same analytic neighborhood. Taylor's formula and the bounded derivatives of $h$, $g$, and the finitely many $e_{2r-1}$ show that every fixed $\sigma$ derivative of $R$ is bounded by a fixed polynomial in $W$. They also give
\[
 |R(\sigma;y,v)|\leq C\sigma(W^3+W).
\]
To see the Gaussian domination explicitly, $W^3\leq C_1W(1+y^2+v^2)$, and $sW\to0$ uniformly on~\eqref{xpc:r2:eq:box}. Hence, after reducing $t_J$,
\[
 \Re R(\sigma;y,v)\leq \frac a4 y^2+\frac D4v^2+C_2.
\]
The derivatives of $e^R$ are $e^R$ times polynomials in fixed derivatives of $R$. Taylor's integral remainder for $e^R$ at degree $J$, multiplied by $e^Q$, therefore gives~\eqref{xpc:r2:eq:uniform}. The Euler--Maclaurin error is $O(s^{4p-2})$ uniformly and is absorbed by the same majorant, using $|e^u-1|\leq |u|e^{|u|}$.

Finally the substitution is invariant under
$(s,y,v)\mapsto(-s,-y,-v)$; the remaining Euler--Maclaurin powers have the same invariance. The coefficient of $s^m$ thus has the claimed parity.
\end{proof}

The Gaussian estimate has no hidden logarithmic loss. Uniformly in the lattice shift and $0<s\leq1$,
\begin{equation}\label{xpc:r2:eq:latticemajorant}
 s\sum_y\int_{\R}W^{b_J}e^{-a y^2/4-Dv^2/4}\,dv\leq C'_J.
\end{equation}
For example, split the real line into unit intervals, use the maximum of the polynomial times a slightly weaker Gaussian on each interval, and bound the number of lattice points there by $1+s^{-1}$. The resulting summable Gaussian series proves~\eqref{xpc:r2:eq:latticemajorant}. Thus the relative integrated remainder in~\eqref{xpc:r2:eq:uniform} is $O(s^{J+1})$.

\subsection{The lattice produces no algebraic oscillations}
Each polynomial-Gaussian term can be extended to the full $(y,v)$ plane with an error of arbitrarily high fixed order by the choice of $K$. Integrate in $v$ first. Gaussian integration turns
$P_m(y,v)e^{Q(y,v)}$ into a polynomial in $y$ times
\[
 \exp\left[-\frac12\left(a+\frac{a^2x_0^2}{D}\right)y^2\right]
 =\exp\left[-\frac{aA}{D}y^2\right].
\]
Poisson summation for a polynomial times this Gaussian replaces its shifted mesh-$s$ sum by $s^{-1}$ times its integral with exponentially small error. Indeed its Fourier transform is a polynomial times
$\exp(-D\xi^2/(4aA))$. At nonzero frequencies $\xi=2\pi m/s$, its absolute value is bounded by a polynomial in $|m|/s$ times
\begin{equation}\label{xpc:r2:eq:poissonerror}
 \exp\left(-\frac{\pi^2 Dm^2}{aAt}\right),\qquad m\neq0.
\end{equation}
The lattice shift only contributes unimodular phases. This also proves uniformity in the fractional part of $x_0/t$. Poisson summation here can equivalently be proved by Fourier expansion of the periodized Schwartz function; its absolutely convergent Fourier series permits termwise evaluation.

The substitution $(y,v)\mapsto(-y,-v)$ and Lemma~\ref{xpc:r2:lem:remainder} make all odd-degree full-plane integrals vanish. Take $J=2M+1$ to obtain an integer-power expansion in $t$, with relative error $O(t^{M+1})$. The leading term, including $d\theta=t^{3/2}dv$ and lattice density $s^{-1}$, is
\begin{align*}
 \frac{g(x_0)t e^{2A/t}}{2\pi}\int_{\R^2}e^{Q(y,v)}\,dy\,dv
 &=\frac{g(x_0)t e^{2A/t}}{\sqrt{a(D+ax_0^2)}}\\
 &=\frac{g(x_0)t e^{2A/t}}{\sqrt{2aA}}
 =C n^{-1/2}e^{B\sqrt n}.
\end{align*}
This proves the existence and remainder claim of Theorem~\ref{xpc:r2:thm:coeff}. It remains to identify the coefficients.

\section{The finite coefficient rule}\label{xpc:r2:sec:rule}
Put $\ell(x)=\log g(x)$ and $H_m=h^{(m)}(x_0)$. Let $Y$ be a formal centered Gaussian with variance $1/a$, meaning that the linear functional $\E$ on polynomials is determined by
\[
 \E Y^{2m}=(2m-1)!!\,a^{-m},\qquad \E Y^{2m+1}=0,\qquad \E1=1.
\]
For every $j\geq0$ define
\begin{equation}\label{xpc:r2:eq:crule}
\begin{split}
 c_j=\coeff{s}{2j}\E\exp\Bigg(&
 \sum_{m\geq3}\frac{H_m}{m!}Y^m s^{m-2}
 +\sum_{m\geq1}\frac{\ell^{(m)}(x_0)}{m!}Y^m s^m\\
 &+\sum_{r\geq1}\sum_{m\geq0}
 \frac{e_{2r-1}^{(m)}(x_0)}{m!}Y^m s^{4r-2+m}
 \Bigg).
\end{split}
\end{equation}
This is a finite rule: discard every monomial of $s$-degree greater than $2j$ before exponentiating, and truncate the exponential at the same degree. Thus one needs $H_m$ only through $m=2j+2$, $\ell^{(m)}$ only through $2j$, and $4r-2+m\leq2j$ in the last sum. All derivatives in this rule are algebraic numbers in $\Q(\sqrt5)$; hence so is every $c_j$.

One way to check the algebra is to set $y=e^{-x}$ and differentiate using $-y\,d/dy$. For $m\geq2$,
\[
 H_m=2^m f^{(m-1)}(2x_0)-f^{(m-1)}(x_0),
\]
while
$\ell(x)=-3x+\tfrac12f(x)+\tfrac12f(2x)$.
At $x_0$ the arguments are $\phi^{-1}$ and $\phi^{-2}$, both in the same quadratic field. Some low-order checks are
\begin{align*}
 H_2&=4-3\phi,&H_3&=7-2\phi,&H_4&=21-24\phi,\\
 \ell'(x_0)&=-3\phi/2-2,&\ell''(x_0)&=\phi+5/2,\\
 e_1(x_0)&=1/12,&e_3(x_0)&=-(\phi+1)/120.
\end{align*}
Equation~\eqref{xpc:r2:eq:crule} is exactly the local real Laplace expansion: substitute $x=x_0+\sqrt{t}\,Y$ in~\eqref{xpc:r2:eq:EM}, divide by the leading Gaussian, and take its moments. The same localization and shifted-lattice estimates as in Section~\ref{xpc:r2:sec:two}, with $v=0$, show at every fixed order that
\begin{equation}\label{xpc:r2:eq:radial}
 F(e^{-t})=\lambda t^{-1/2}e^{A/t}
 \left(\sum_{j=0}^M c_jt^j+O_M(t^{M+1})\right).
\end{equation}
This supplies an elementary specialization of the prior radial theory; no originality is claimed for~\eqref{xpc:r2:eq:radial}. Applying~\eqref{xpc:r2:eq:crule} gives~\eqref{xpc:r2:eq:cs}. The package computes these values using exact rational pairs and checks them with a second finite polynomial calculation.

\begin{writenote}[A check of the rule, 7 October 2026]
The write evaluated the finite rule~\eqref{xpc:r2:eq:crule} with its own code
(mpmath at 50 digits, truncated power series in $s$ with polynomial
coefficients in $Y$, the Euler--Maclaurin terms through $r=2$): $c_0,\dots,c_4$
agree with~\eqref{xpc:r2:eq:cs} to better than $10^{-43}$, and the odd
$s$-degrees vanish. It also confirmed the displayed values of $H_2$, $H_3$,
$H_4$, $\ell'(x_0)$, $\ell''(x_0)$, $e_1(x_0)$, $e_3(x_0)$, $h(x_0)=A$ and
$\lambda$, and $a=3\phi-4=0.8541019662\ldots$, $D=0.4601936581\ldots>0$. This
is a numerical check of the stated algebra, not a second derivation.
\end{writenote}

\section{Identification with the coefficient saddle}\label{xpc:r2:sec:identify}
We justify the relation between~\eqref{xpc:r2:eq:crule} and the actual coefficients without assuming a sectorial version of~\eqref{xpc:r2:eq:radial}.

\begin{lemma}[Finite Gaussian contour identity]\label{xpc:r2:lem:contour}
For a polynomial $P$ and an affine line $\Gamma=b+e^{\ii\alpha}\R$, oriented in the increasing real parameter, with $|\alpha|<\pi/4$,
\[
 \int_\Gamma P(u)e^{-a u^2/2}\,du
 =\int_\R P(u)e^{-a u^2/2}\,du.
\]
\end{lemma}
\begin{proof}
First translate a truncated line and then rotate it within the sectors of Gaussian decay. The integrand is entire. Along the connecting segments at radius $R$, its modulus is bounded by a polynomial in $R$ times $\exp(-cR^2+O(R))$, for some $c>0$ if the rotation remains in a closed subsector of $|\alpha|<\pi/4$. Their integrals tend to zero. Cauchy's theorem and passage to the limit prove the identity. A fixed translation does not alter the quadratic decay.
\end{proof}

All coefficient calculations below are finite Taylor and Gaussian moment calculations already justified by Lemma~\ref{xpc:r2:lem:remainder}; the contour identity does not transfer an uncontrolled real-axis error. In~\eqref{xpc:r2:eq:scaling} use the principal square root and set $u=(x-x_0)/\sqrt z$. As real $y$ varies, its line is
\begin{equation}\label{xpc:r2:eq:uline}
 u=y\sqrt{1-\ii s v}-\frac{\ii x_0v}{\sqrt{1-\ii s v}},\qquad
 dk=z^{-1/2}\,du.
\end{equation}
For fixed real $v$ and small $s$, its direction lies inside a Gaussian decay sector. At each finite Taylor order the amplitude is a polynomial in $u$. Lemma~\ref{xpc:r2:lem:contour} therefore identifies the $k$-saddle moments with the radial moments~\eqref{xpc:r2:eq:crule}, giving the finite formal expression
\begin{equation}\label{xpc:r2:eq:formalradial}
 \lambda z^{-1/2}e^{A/z}\sum_j c_jz^j.
\end{equation}
Odd moments in $\sqrt z$ vanish.

This coefficientwise argument is compatible with the remaining $v$ integral. Before contour translation the absolute quadratic factor is
$e^{-a y^2/2-Dv^2/2}$. A translated-contour bound may introduce
$e^{a x_0^2v^2/2}$, but it is accompanied by $e^{-Av^2}$ from the outer saddle, leaving $e^{-Dv^2/2}$. Since $D>0$, every resulting finite polynomial moment is absolutely integrable. The actual remainders still come from~\eqref{xpc:r2:eq:uniform}, not from any assertion about~\eqref{xpc:r2:eq:formalradial} away from a finite expansion.

For clarity, the remaining one-dimensional moment algebra can be encoded by the modified Bessel saddle. The $j$th term corresponds to
\begin{equation}\label{xpc:r2:eq:Bessel}
 (A/n)^{j/2+1/4} I_{-j-1/2}(2\sqrt{An}).
\end{equation}
This notation records the formal dominant-exponential saddle of the contour integrand $z^{j-1/2}e^{A/z+nz}$; it is not an equality for a separately truncated Cauchy arc. The leading Gaussian normalizes $I_\nu(X)$ by $e^X/\sqrt{2\pi X}$. For a direct check of every coefficient, write its dominant series as $e^XX^{-1/2}\sum_{m\geq0}b_mX^{-m}$, $b_0=1$. Substitution into
\[
 X^2I_\nu''+XI_\nu'-(X^2+\nu^2)I_\nu=0
\]
gives
\[
 b_m=-\frac{4\nu^2-(2m-1)^2}{8m}\,b_{m-1}.
\]
The same differential identity for the formal saddle follows by differentiating the contour integrand and integrating a total derivative; Gaussian tails make its boundary term zero at each finite order. Thus the recurrence determines its moments uniquely. For $\nu=-j-1/2$ it terminates at $m=j$, since
\[
 \prod_{\ell=1}^m\bigl((2j+1)^2-(2\ell-1)^2\bigr)
 =4^m\frac{(j+m)!}{(j-m)!}\quad(0\leq m\leq j).
\]
Accordingly
\begin{equation}\label{xpc:r2:eq:Besselpoly}
 I_{-j-1/2}(X)\sim\frac{e^X}{\sqrt{2\pi X}}
 \sum_{m=0}^{j}\frac{(-1)^m(j+m)!}{m!(j-m)!(2X)^m}
\end{equation}
at algebraic order relative to $e^X$; the other exponential is negligible. These are also the standard coefficients in DLMF~\cite{dlmfBessel}.
Combining~\eqref{xpc:r2:eq:formalradial}--\eqref{xpc:r2:eq:Besselpoly}, the term indexed by $j,m$ has relative $n^{-(j+m)/2}$ coefficient
\[
 c_jA^{j/2}\frac{(-1)^m(j+m)!}{m!(j-m)!(4\sqrt A)^m}.
\]
Put $r=j+m$. The condition $0\leq m\leq j$ is exactly
$\lceil r/2\rceil\leq j\leq r$, proving~\eqref{xpc:r2:eq:transfer} and completing Theorem~\ref{xpc:r2:thm:coeff}.

\section{Inverse expansions and integer thresholds}\label{xpc:r2:sec:inverse}
Define the formal logarithmic coefficients
\begin{equation}\label{xpc:r2:eq:alpha}
 \ell_r=\coeff{z}{r}\log\left(1+\sum_{j\geq1}d_jz^j\right),
 \qquad \alpha_r=B^r\ell_r.
\end{equation}
Only finitely many terms enter at a fixed order. These $\ell_r$ are constants, distinct from the function $\ell(x)=\log g(x)$ in Section~\ref{xpc:r2:sec:rule}. For a large real threshold $Y$ write
\begin{equation}\label{xpc:r2:eq:Tw}
 T=\log\frac{Y}{CB},\qquad w=\log T.
\end{equation}
The variable $s=B\sqrt n$ used only in this section differs from the small parameter $s=\sqrt t$ in the saddle proof. The normalization is chosen so that Theorem~\ref{xpc:r2:thm:coeff} gives
\begin{equation}\label{xpc:r2:eq:loginverse}
 \log\frac{a(n)}{CB}
 =s-\log s+\sum_{r=1}^J\alpha_rs^{-r}+O_J(s^{-J-1}).
\end{equation}

\begin{theorem}[Continuous models and two ceiling bounds]\label{xpc:r2:thm:inverse}
For fixed $J\geq1$ define polynomials $P_j(w)$ successively by
\begin{align}
 S_0(u)&=1+uw,\notag\\
 P_j(w)&=\coeff{u}{j}\left\{\log S_{j-1}(u)
 -\sum_{r=1}^j\alpha_ru^r S_{j-1}(u)^{-r}\right\},\label{xpc:r2:eq:Prec}\\
 S_j(u)&=1+uw+\sum_{i=1}^jP_i(w)u^{i+1}.\notag
\end{align}
Set
\begin{equation}\label{xpc:r2:eq:X}
 X_J(Y)=B^{-2}\left(T+\log T+
             \sum_{j=1}^J\frac{P_j(\log T)}{T^j}\right)^2.
\end{equation}
Then $\deg P_j\leq j$, with
\[
 P_1(w)=w-\alpha_1,\qquad
 P_2(w)=-\tfrac12w^2+(1+\alpha_1)w-\alpha_1-\alpha_2.
\]
Along the exact values $Y=a(n)$,
\begin{equation}\label{xpc:r2:eq:Xerror}
 X_J(a(n))=n+O_J\left(\frac{(1+\log\log a(n))^{J+1}}
                              {(\log a(n))^J}\right).
\end{equation}
For $N(Y)=\min\{n\geq1:a(n)\geq Y\}$ there is a constant $K_J>0$ such that, for all sufficiently large $Y$,
\begin{equation}\label{xpc:r2:eq:ceil}
 \ceil{X_J(Y)-\epsilon_J(Y)}\leq N(Y)
 \leq\ceil{X_J(Y)+\epsilon_J(Y)},\qquad
 \epsilon_J(Y)=K_J\frac{(1+\log T)^{J+1}}{T^J}.
\end{equation}
Alternatively let $s_J(Y)$ be the large positive root of
\begin{equation}\label{xpc:r2:eq:implicit}
 s-\log s+\sum_{r=1}^J\alpha_rs^{-r}=T.
\end{equation}
It is eventually unique, and $Z_J(Y)=s_J(Y)^2/B^2$ satisfies
\begin{equation}\label{xpc:r2:eq:impliciterror}
 Z_J(a(n))=n+O_J(n^{-J/2}).
\end{equation}
There is also a two-ceiling bracket centered at $Z_J(Y)$ with radius $K'_JT^{-J}$ for some $K'_J>0$ and all sufficiently large $Y$.
\end{theorem}

\begin{proof}
Let $H_J(s)=s-\log s+\sum_{r=1}^J\alpha_rs^{-r}$.
For large $s$, $H_J'(s)=1+O(1/s)>1/2$. Since $H_J(s)\to\infty$, it has a unique inverse on a sufficiently large half-line. Substitute
$s=T S$ with $u=T^{-1}$ and $S=1+uw+\sum_{j=1}^JP_j(w)u^{j+1}$. The equation $H_J(s)=T$ is equivalent to
\[
 \sum_{j=1}^JP_j(w)u^j
 =\log S-\sum_{r=1}^J\alpha_ru^rS^{-r},
\]
up to the truncation error. The coefficient of $u^j$ on the right does not depend on $P_j$, which first enters $S$ at degree $j+1$. Hence~\eqref{xpc:r2:eq:Prec} cancels the residual successively. Formal power-series multiplication proves by induction that $\deg P_j\leq j$; more explicitly every term contributing at $u^j$ has total $w$-degree at most $j$. Evaluation at $w=\log T$ gives the residual bound
$O_J((1+w)^{J+1}T^{-J-1})$. These bounds also follow from convergent Taylor series of $\log(1+v)$ and $(1+v)^{-r}$ for $v=O(w/T)$. The mean value theorem then gives
\begin{equation}\label{xpc:r2:eq:serror}
 s_J(Y)=T+\log T+\sum_{j=1}^J P_j(\log T)T^{-j}
       +O_J((1+\log T)^{J+1}T^{-J-1}).
\end{equation}
The error in~\eqref{xpc:r2:eq:loginverse} changes $s_J(a(n))$ by
$O_J(s^{-J-1})$; squaring changes the index by $O_J(s^{-J})$. This proves~\eqref{xpc:r2:eq:impliciterror}. Squaring~\eqref{xpc:r2:eq:serror} costs one factor of $T$, proving~\eqref{xpc:r2:eq:Xerror}.

For the integer assertion, introduce the smooth log model
\[
 L_J(x)=\log C-\tfrac12\log x+B\sqrt x+
                      \sum_{r=1}^J\ell_rx^{-r/2}.
\]
Its derivative is $B/(2\sqrt x)+O(x^{-1})>0$ for large $x$ and
$\log a(n)=L_J(n)+O_J(n^{-(J+1)/2})$ for integers $n$.
Let $x_*=Z_J(Y)$; equation~\eqref{xpc:r2:eq:implicit} is exactly
$L_J(x_*)=\log Y$. In a fixed relative neighborhood of $x_*$, the derivative is bounded below by a positive constant times $x_*^{-1/2}$, while the integer error is bounded by a constant times $x_*^{-(J+1)/2}$. Choose
$\delta=Kx_*^{-J/2}$ with $K$ sufficiently large. The mean value theorem ensures that every integer in that neighborhood with $n<x_*-\delta$ has $a(n)<Y$, and that $n=\ceil{x_*+\delta}$ has $a(n)\geq Y$.

For completeness, these local comparisons imply the global threshold assertion. Section~\ref{xpc:r2:sec:shape} proves eventual strict increase. The finitely many earlier values are all below $Y$ when $Y$ is large. A nearby integer below $x_*-\delta$ bounds all earlier integers by eventual monotonicity. Therefore
\[
 \ceil{x_*-\delta}\leq N(Y)\leq\ceil{x_*+\delta}.
\]
There is no endpoint ambiguity: $n<\ceil{x_*-\delta}$ implies $n<x_*-\delta$, whether or not $x_*-\delta$ is integral. Since $x_*\asymp T^2$, this gives the radius $K'_JT^{-J}$ around $Z_J$. Equation~\eqref{xpc:r2:eq:serror} implies
$|x_*-X_J(Y)|=O_J((1+\log T)^{J+1}T^{-J})$; enlarging its constant gives~\eqref{xpc:r2:eq:ceil}.
\end{proof}

If the two ceilings in a bracket agree, their common value is the exact threshold. This is an asymptotic existence statement, not a certified numerical interval until its constants and range are made effective. In particular one may not infer $N(Y)=\ceil{X_J(Y)}$ for every large real $Y$: a threshold can be arbitrarily close to a jump. Even an $o(1)$ continuous error does not remove the need for rounding bounds. These inversion and ceiling mechanisms are standard asymptotic tools, also used in prior ProveIt reports~\cite{proveit}; no novelty is claimed for the general mechanism.

\begin{remark}[{[write]} The inverses and the transseries volume, 7~October 2026]\label{xpc:r2:rem:transseries}
Statement by statement, against the volume named in the Guide (its symbols
carry the subscript ``vol'').
\begin{enumerate}[label=(\alph*),leftmargin=*]
\item \emph{Instance.} For $Y$ beyond the early values, $N(Y)$ is the integer
staircase of \file{p0:def:three-inverses} with $A_n=a(n)$ and $n_1$ the onset
of Corollary~\ref{xpc:r2:cor:shape}, so \file{p0:thm:staircase}(1) applies.
In the b-file, the last $n$ with $a(n+1)\le a(n)$ is $n=42$; that is a finite
observation.
\item \emph{Formal instance.} The construction~\eqref{xpc:r2:eq:Prec} is
\file{p0:thm:flattening}(4) with $a_\vol=1$, $b_\vol=-1$, $q_r=\alpha_r$,
$L_\vol=T$ and $H_\vol=\log T=w$: the volume's ansatz
$x=L-bH+\sum_nP_n(H)L^{-n}$ is the bracket in~\eqref{xpc:r2:eq:X}, and its
Bell recurrence~\file{p0:eq:P-bell} gives $P_1=H-\alpha_1$ and
$P_2=-\tfrac12H^2+(1+\alpha_1)H-\alpha_1-\alpha_2$, the source's $P_1$, $P_2$
(checked by hand and symbolically); the leading coefficient $[H^n]P_n=b^{n+1}/n$
of \file{p0:eq:P-shape} matches $-\tfrac12$ at $n=2$. The remainder
in~\eqref{xpc:r2:eq:serror} is the source's own; its residual-over-slope step
is \file{p0:thm:backward-error}.
\item \emph{Instance, not used by the source.} $H_J(s)$ is exactly of
exponential--power type in $s$ with data $(1,1,1,-1,1,\sum_{r\le J}\alpha_r\tau^r)$
(\file{p0:def:model}), so its large root $s_J(Y)$ also has the Lambert-centred
expansion of \file{p0:thm:lambert-centered}(5) about $-W_{-1}(-e^{-T})$; the
source expands about $T$ instead, which is the flattened form of (b).
\item \emph{Analogues.} The bracket~\eqref{xpc:r2:eq:ceil} and the bracket
about $Z_J(Y)$ are, as stated and proved, analogues of
\file{p0:thm:staircase}(2).
\end{enumerate}
No novelty is claimed for any of these inversions; the source says the same
of its mechanism.
\end{remark}

\section{Eventual increase and strict log concavity}\label{xpc:r2:sec:shape}
\begin{corollary}\label{xpc:r2:cor:shape}
The sequence $a(n)$ is eventually strictly increasing and eventually strictly log-concave. In particular, eventually
$a(n)^2>a(n-1)a(n+1)$.
\end{corollary}
\begin{proof}
Retain the leading term and $d_1$ in~\eqref{xpc:r2:eq:main}; its relative remainder is $O(n^{-1})$. Taking the ratio at adjacent indices gives
\[
 \frac{a(n+1)}{a(n)}=1+\frac B{2\sqrt n}+O(n^{-1})>1.
\]
No smoothness of that remainder is needed: the ratio of the two remainder factors is $1+O(n^{-1})$.
For log-concavity use the expansion through $d_3$ and take a logarithm:
\[
 \log a(n)=\log C-\tfrac12\log n+B\sqrt n+
       \sum_{r=1}^3\ell_r n^{-r/2}+O(n^{-2}).
\]
The second difference of the remainder is still $O(n^{-2})$ by the triangle inequality. For the displayed smooth functions, Taylor's theorem on $[n-1,n+1]$ gives
\[
 \log a(n+1)-2\log a(n)+\log a(n-1)
 =-\frac B{4n^{3/2}}+O(n^{-2})<0.
\]
The main negative term dominates. This proof neither differentiates nor assumes differentiability of an asymptotic error attached to a discrete sequence.
\end{proof}
The argument proves no effective last exception. Finite searches are useful checks but cannot by themselves establish that their last observed failure is the global last failure.

\section{Computation and reproducibility}\label{xpc:r2:sec:repro}
The accompanying archive contains this TeX source, a compiled PDF, original exact code, generated tables, a manifest, and offline rebuild instructions. It does not include third-party papers, downloaded OEIS data, or private working reviews. Exact computations use integer and rational arithmetic in $\Q(\sqrt5)$, with no conversion to floating point. Input and consistency guards use explicit exceptions so they remain active under \texttt{python -O}. The fixed-order generator accepts a nonnegative integer order and truncates by the entire degree of the expansion parameter; it is not restricted to the four displayed coefficients.

The exact checks use two independently organized finite Gaussian calculations, and partition counts are obtained from~\eqref{xpc:r2:eq:F} by integer dynamic programming. A second small-range count routine checks the defining partition condition independently. The numerical diagnostics, when requested, use the separately disclosed optional \texttt{mpmath} dependency; they are not part of the exact proof. No network access is required by the rebuild.

For orientation, the relative errors of the coefficient approximations are defined as approximation divided by exact count, minus one. At $n=10000$, the approximation through $d_3$ has error approximately $-3.3063\cdot10^{-7}$; including $d_4$ changes it to approximately $-1.4600\cdot10^{-8}$. The three-term implicit logarithmic inverse at the exact $a(10000)$ returns an index about $0.0000271$ above $10000$. These are diagnostics at selected arguments, not certified remainder bounds, optimal truncation claims, or evidence about exact integer rounding at arbitrary thresholds.

The source and generated tables make the calculation auditable without making the proof depend on numerical agreement. Section~\ref{xpc:r2:sec:two} proves the asymptotic remainder, while Section~\ref{xpc:r2:sec:rule} and equation~\eqref{xpc:r2:eq:transfer} specify the constants.

\begin{writenote}[The diagnostics recomputed, 7 October 2026]
From~\eqref{xpc:r2:eq:cs} and~\eqref{xpc:r2:eq:ds}, and again from the
transfer formula~\eqref{xpc:r2:eq:transfer}, the write finds
$d_1=2.424917491063005\ldots$, $d_2=3.187616630712046\ldots$,
$d_3=9.439637377883243\ldots$, $d_4=32.379308690155\ldots$, whose twelve-place
truncations are the values printed in Theorem~\ref{xpc:r2:thm:coeff}. With
$a(10000)$ from the OEIS b-file, the relative errors through $d_3$ and $d_4$
are $-3.30626\ldots\cdot10^{-7}$ and $-1.45998\ldots\cdot10^{-8}$, and the
three-term implicit inverse returns $10000+2.709\ldots\cdot10^{-5}$, matching
the approximate values printed above. In the b-file, the last $n\le9999$ with
$a(n)^2\le a(n-1)a(n+1)$ is $n=445$ and the last with $a(n+1)\le a(n)$ is
$n=42$; finite observations, which prove no last exception.
\end{writenote}

\section{Source comparison and limitations}\label{xpc:r2:sec:sources}
The comparison below is deliberately bounded. It supports attribution and distinguishes neighboring models; it is not a priority certificate.
\begin{enumerate}[leftmargin=*]
\item The OEIS entry A118096~\cite{oeisA118096} gives~\eqref{xpc:r2:eq:F} and a leading equivalent attributed to V\'aclav Kot\v e\v sovec, dated 13 June 2025. The absence of a proof beside an OEIS formula does not make it an open conjecture or establish novelty for a proof.
\item McIntosh's general-product result~\cite{mcintosh}, in particular Theorem 2, and Zhou's Eulerian-series results~\cite{zhou}, Theorems 1--2, supply prior real-axis all-order machinery. The present series is an Eulerian-series specialization, since
\[
 F(q)=q^3\sum_{m\geq0}q^{3m}
          \frac{(q^{2m+3};q)_\infty}{(q^{m+1};q)_\infty}.
\]
The radial expansion and golden-ratio saddle are therefore prior-art territory. The complete-circle proof here is included to make the fixed-order coefficient argument self-contained, not to assert that such an argument is absent elsewhere.
\item Pittel~\cite{pittel} studies partitions with bounded largest-to-smallest ratio. Its publisher abstract identifies the condition $\max/\min\leq k$, $k>1$, and a parameter satisfying $w^k+w=1$. This is directly adjacent literature. The full theorem text was not accessible in the present comparison, so its enumerative precision and possible implications for this exact-boundary model remain unresolved here. No source clearance or exclusion of overlap is claimed.
\item The neighboring mock-theta models are distinct. Andrews, Dixit, and Yee~\cite{ady} identify $q\chi_1(q)$ with $\max<2\min$ partitions and associate $\chi_0(q)$ with $\max\leq2\min$ and a unique smallest part. If
\[
 H(q)=\sum_{k\geq1}\frac{q^k}{\prod_{j=k}^{2k}(1-q^j)},
\]
then direct termwise subtraction gives $F(q)=H(q)-q\chi_1(q)$. Indeed the strict class has summand $q^k/\prod_{j=k}^{2k-1}(1-q^j)$; subtracting from the weak class leaves $q^{3k}/\prod_{j=k}^{2k}(1-q^j)$. This does not identify $H$ with a classical product or mock theta function, and the unique-minimum condition must not be dropped.
\item The Bessel coefficient product in Section~\ref{xpc:r2:sec:identify} is classical~\cite{dlmfBessel}. Formal series reversion and derivative-based inverse control are standard; related inverse reports in ProveIt~\cite{proveit} are not a basis for claiming a new inversion method here.
\end{enumerate}
The appropriate description is an independently derived, explicitly computable refinement with an incomplete historical comparison. No claim is made to have solved a previously unsolved problem.

\section{Further questions}\label{xpc:r2:sec:questions}
Several questions remain outside the proved scope.
\begin{itemize}[leftmargin=*]
\item An accessible full reading of Pittel's paper is needed to complete the closest historical comparison, before any originality assessment.
\item Effective Euler--Maclaurin, localization, and Gaussian constants could turn the asymptotic inverse brackets into certified computable bounds and prove explicit last exceptions for increase or log-concavity.
\item For largest part $m$ times smallest part, $m\geq3$, the finite interval becomes $[k,mk]$. The same architecture suggests further coefficient rules, but the constants, angular estimates, and inverse conclusions require their own proofs.
\item The series $H(q)$ in Section~\ref{xpc:r2:sec:sources} may admit useful exact identities. The elementary mock-theta subtraction alone does not settle its analytic structure.
\item The present theorem fixes the order before $n\to\infty$. Growth of $c_j$, optimal truncation, and exponentially small corrections require analysis beyond the finite-order estimates proved here.
\end{itemize}
\begin{writenote}[7 October 2026]
Under Vladimir's standing rule of 4~October 2026, the write found no claim of
this Part that is unproved as stated and none that is wrong. The comparison
with Pittel's 2007 paper~\cite{pittel} remains open, as the first question
says: the write's attempt to read it at the publisher on 7~October 2026 was
refused (HTTP~403), so neither the source nor the write has read its theorems,
and nothing here should be read as excluding an overlap. By the source's
account (item~3 of Section~\ref{xpc:r2:sec:sources}), the abstract names a
parameter with $w^k+w=1$; at $k=2$ its root in $(0,1)$ is $\phi^{-1}=e^{-x_0}$,
this Part's saddle point, which makes the paper the closest comparison; that
coincidence is all the write can add. On the
second question, note only that in the b-file the last $n$ with
$a(n+1)\le a(n)$ is $n=42$ and the last with $a(n)^2\le a(n-1)a(n+1)$ is
recorded in the note at the end of Section~\ref{xpc:r2:sec:repro}; finite
searches prove no last exception.
\end{writenote}


\part{Deficit asymptotics for partitions with divisible extreme parts}\label{xpc:dv:part}
\setlength{\parindent}{\xpcdefparindent}\setlength{\parskip}{0pt plus 1pt}\setlist{itemsep=3pt,topsep=5pt}\allowdisplaybreaks[2]\setlength{\emergencystretch}{2em}
\partsource{Bundle Report~218, archive \file{Report218-reproducibility.zip}, delivered as \file{report218.tex} (22 pages). Delivered title block ``Deficit asymptotics for partitions with divisible extreme parts / Report 218 / 4 October 2026''. Printed as delivered, with \textbf{[write]} remarks and notes; Report~218's Section~$k$ is Section~$k+21$ here, its equation~($k$) is equation~($k+103$), and its appendix keeps the letter~A; statements move with their sections. Its labels carry \file{xpc:dv:}; its two unlabelled sections were labelled; its citation key \file{oeis} is \file{oeisA117086} and \file{nr} is \file{ngorhoades}. Its four generated tables are printed inline. This Part is self-contained.}
\begin{partabstract}
Let $a(n)$ count partitions of $n$ whose largest part is divisible by their smallest part, the sequence A117086, and put $D(n)=p(n)-a(n)$. We prove
\[
 D(n)\sim\frac{\pi}{2\sqrt{6n}}\,p(n)
\]
and give its expansion to every fixed order, including an exact coefficient algorithm and six explicit relative coefficients. A finite-polynomial reduction of the fixed-minimum Fourier series and an elementary bound on the whole coefficient circle control the residue obstruction. The minimum tail is estimated before coefficient extraction; a finite-shift argument based on the Rademacher formula then supplies all orders. We also prove fixed-minimum residue equidistribution, a comparison with a reciprocal-minimum weighted count, sharpened to an explicit stretched-exponential bound for $n\ge90$, finite-$n$ enclosures, and all-order Lambert-branch inverse brackets. Exact rational code regenerates the coefficients, counting checks, and a certified enclosure at $n=1000$. The elementary equivalence $a(n)\sim p(n)$ is background. No historical-priority claim, uniform growing-order expansion, or convergence claim is made.
\end{partabstract}

\section{The sequence and the main result}\label{xpc:dv:sec:result}
A partition is a finite nonincreasing list of positive integers. For a nonempty partition $\lambda$, write $|\lambda|$ for its sum and $\min\lambda$, $\max\lambda$ for its extreme parts. Let
\[
 a(n)=\#\{\lambda\vdash n:\min\lambda\mid\max\lambda\},\qquad
 D(n)=p(n)-a(n)\quad(n\ge1),
\]
where $p(n)$ is the unrestricted partition number and $p(0)=1$. For computational bookkeeping only, set $a(0)=0$. Thus $D(0)$ is not included in any generating series below.

The retrieved OEIS entry A117086, credited to Vladeta Jovovic on 17 April 2006, gives this definition and the generating function \cite{oeisA117086}
\begin{equation}\label{xpc:dv:eq:exact-a}
 \sum_{n\ge1}a(n)q^n
 =\sum_{L\ge0}\sum_{k\ge1}
   \frac{q^{(L+1)k}}{\prod_{i=k}^{Lk}(1-q^i)}.
\end{equation}
An empty product is one. The $L=0$ term counts the one-part partition $(k)$. For $L\ge1$, removing one smallest part $k$ and one largest part $Lk$ leaves arbitrary parts between those endpoints. When $L=1$, this counts at least two copies of $k$, so the singleton is counted exactly once overall.

Adjoining a part $1$ is an injection from partitions of $n-1$ to the partitions counted by $a(n)$. Consequently
\begin{equation}\label{xpc:dv:eq:sandwich}
 p(n-1)\le a(n)\le p(n)\qquad(n\ge1).
\end{equation}
In particular, $a$ is weakly increasing, since $a(n+1)\ge p(n)\ge a(n)$. The familiar partition asymptotic already gives $a(n)\sim p(n)$. The quantity requiring further analysis is the deficit $D(n)$.

Throughout put
\begin{equation}\label{xpc:dv:eq:scales}
 \A=\frac{\pi^2}{6},\qquad N=n-\frac1{24},\qquad
 t=\sqrt{\frac{\A}{N}},\qquad
 \B(n)=\frac{\exp(2\sqrt{\A N})}{4\sqrt3\,N}.
\end{equation}
Write $(q;q)_j=\prod_{h=1}^j(1-q^h)$, $(q;q)_0=1$, and
$P(q)=(q;q)_\infty^{-1}=\sum_{n\ge0}p(n)q^n$. For the fixed-order expansions and fixed-minimum residue results, every order, minimum, modulus, and truncation depth is fixed before $n\to\infty$. Section~\ref{xpc:dv:sec:quantitative} separately uses exact uniform inequalities at a specified growing auxiliary cutoff.

\begin{theorem}[All fixed orders]\label{xpc:dv:thm:main}
There are explicit polynomials $c_m\in\mathbb Q[\A^{-1}]$, with $c_0=1$, such that for every integer $R\ge0$,
\begin{equation}\label{xpc:dv:eq:main}
 D(n)=\frac{\B(n)t}{2}
 \left(\sum_{m=0}^R c_m t^m+O_R(t^{R+1})\right).
\end{equation}
They are generated by \eqref{xpc:dv:eq:c-generator} below. In particular,
\[
 \frac{D(n)}{p(n)}\sim\frac{\pi}{2\sqrt{6n}}.
\]
\end{theorem}
The other principal conclusions are Theorem~\ref{xpc:dv:thm:residues}, concerning fixed-minimum residue classes; Corollary~\ref{xpc:dv:cor:weight} and Theorem~\ref{xpc:dv:thm:quantitative}, comparing $D$ with a rational weighted count; Theorem~\ref{xpc:dv:thm:enclosure}, an effective finite-$n$ inequality; and Theorem~\ref{xpc:dv:thm:inverse}, an all-order asymptotic inverse bracket. All proofs are given below, using the classical eta transformation and Rademacher partition formula as stated explicitly. These standard input identities are distinguished from the elementary estimates derived here.

\begin{remark}[{[write]} The OEIS entry, quoted, 7~October 2026]\label{xpc:dv:rem:oeis}
OEIS A117086, read on 7~October 2026 in the internal format, is at revision
\#19 (3~September 2017), author Vladeta Jovovic (17~April 2006), offset~1. Its
name is ``Number of partitions of n such that the largest part is a multiple
of the smallest part.''; its formula line is the generating
function~\eqref{xpc:dv:eq:exact-a}; a comment of Emeric Deutsch (21~April 2006)
gives the conjugate reading (the number of parts is a multiple of the
multiplicity of the largest part); the example ``a(7)=12'' lists the three
excluded partitions [5,2], [4,3], [3,2,2]; it cross-references A118096
(Part~II) and carries a b-file by Alois P.~Heinz for $1\le n\le1000$. It
states no asymptotic formula and no conjecture. The bibliography
entry~\cite{oeisA117086} prints the name as a title in its own words (``is
divisible by'' for ``is a multiple of''). The revision number was visible to
the write; like the source, it did not inspect the revision history. All 1000
b-file terms agree with the shipped table
\file{data/218-divisible-generated-exact_values.csv}, whose partition numbers
the write also recomputed, and brute-force counts agree for $n\le40$. Nothing
was submitted to the OEIS.
\end{remark}

\section{Prior frameworks and source boundaries}\label{xpc:dv:sec:prior}
The coefficient and residue methods used here have substantial antecedents. Bridges, Franke, and Garnowski \cite{bfg} study the coefficients of $(\zeta q;q)_\infty^{-1}$ and residue biases for the number of parts. Partition conjugation gives the same unrestricted distribution for the largest part. Their Theorem 3.2 resolves much finer angular asymptotics than the whole-circle upper bound used in this report. Their Theorem 1.3 concerns the distinguished root $\zeta_b=\exp(2\pi i/b)$ for $b\ge5$; it must not be read as the same formula for every primitive $b$th root. The angular classification excludes two crossover angles without proving they are irrational multiples of $\pi$. Our elementary bound applies to every nontrivial unit-modulus twist, so it does not need to resolve those crossover cases. A fixed minimum is an additional restriction, handled here by an exact polynomial reduction.

Ngo and Rhoades \cite{ngorhoades} prove arbitrary-order asymptotics for smallest-part parity and expected smallest-part size in Theorems 1.1 and 1.4. They credit their polynomial coefficient-transfer proposition, Proposition 1.8, to Wright. Cesana, Craig, and Males \cite{ccm}, Proposition 2.3 and Theorem 3.1, give general coefficient-transfer and equidistribution frameworks. Those transfer theorems require control away from the positive radius as well as a major-arc expansion. A relative little-$o$ estimate on a twist alone does not establish agreement to every algebraic order. Here we prove an exponential Fourier bound and control the entire remaining minimum tail before extracting coefficients.

The weighted sum-of-tails identity in Section~\ref{xpc:dv:sec:weighted} is a specialization of established arbitrary-weight machinery. Andrews and Pedro Freitas \cite{af} give the general Abel-lemma framework. We verified their publication metadata but did not inspect the original subscription-only full text. Gupta \cite{gupta} reproduces their Theorem 4.1 as equation (1.9) in arXiv version 1 and equation (1.7) in the published article. We give a direct absolute-convergence proof of the specialization used here; that proof does not close the original-text access gap.

The retrieved A117086 entry displays no asymptotic formula, but its live revision history was not accessible. A bounded literature screen did not locate the exact deficit theorem or reciprocal-minimum specialization presented here. Non-discovery does not establish novelty. This report makes no global novelty claim, and no quantum-modularity claim for the weighted series. The general twisted-product, all-order transfer, and arbitrary-weight tail frameworks are prior work.

\section{Fixed-minimum Fourier reduction}\label{xpc:dv:sec:fourier}
For an integer $k\ge1$, $|q|<1$, and $|z|\le1$, define
\[
 H_k(z,q)=\sum_{\min\lambda=k}z^{\max\lambda}q^{|\lambda|},
 \qquad P_z(q)=\prod_{j\ge1}(1-zq^j)^{-1}.
\]
The products and sums converge absolutely in this domain.

\begin{lemma}[Exact polynomial reduction]\label{xpc:dv:lem:polynomial}
One has
\begin{align}
 H_k(z,q)&=q^k(q;q)_{k-1}P_z(q)+R_k(z,q),\label{xpc:dv:eq:reduction}\\
 R_k(z,q)&=(zq)^k-q^k\sum_{M=0}^{k-1}(zq)^M
                       \prod_{j=M+1}^{k-1}(1-q^j).\label{xpc:dv:eq:remainder}
\end{align}
The correction $R_k$ is a polynomial of $q$-degree at most $k(k+1)/2$, and $R_k(1,q)=0$.
\end{lemma}
\begin{proof}
For maximum $M>k$, removal of one endpoint of each size leaves the contribution
$z^Mq^{k+M}/\prod_{j=k}^{M}(1-q^j)$. For $M=k$, the same expression counts constant partitions of multiplicity at least two; adding $(zq)^k$ restores exactly the singleton. Thus
\[
 H_k(z,q)=q^k(q;q)_{k-1}
             \sum_{M\ge k}\frac{(zq)^M}{(q;q)_M}+(zq)^k.
\]
Euler's identity $\sum_{M\ge0}u^M/(q;q)_M=(u;q)_\infty^{-1}$ proves \eqref{xpc:dv:eq:reduction} after removing $M<k$. For completeness, the series on the left satisfies $(1-u)F(u)=F(qu)$ and $F(0)=1$; iterating this equation proves the product identity for $|u|<1$. Each summand removed has $q$-degree
\[
 k+M+\sum_{j=M+1}^{k-1}j
 =\frac{k(k+1)}2-\frac{M(M-1)}2.
\]
Finally $\sum_{M=0}^{k-1}q^M/(q;q)_M=1/(q;q)_{k-1}$, by finite telescoping, so $R_k(1,q)=0$.
\end{proof}
For example $R_1(z,q)=q(z-1)$: finite corrections cannot be discarded at every index.

Let $F_k(n)$ count partitions of $n$ with minimum exactly $k$. The lemma gives
\begin{equation}\label{xpc:dv:eq:Fk}
 \sum_{n\ge1}F_k(n)q^n=q^k(q;q)_{k-1}P(q).
\end{equation}
For $m\ge2$, $0\le b<m$, put
\[
 C_{k,m,b}(n)=\#\{\lambda\vdash n:\min\lambda=k,
                                 \ \max\lambda\equiv b\pmod m\}.
\]
Writing $\zeta_m=\exp(2\pi i/m)$, the elementary root filter is
\begin{equation}\label{xpc:dv:eq:filter}
 C_{k,m,b}(n)=\frac1m\sum_{h=0}^{m-1}\zeta_m^{-bh}
                         [q^n]H_k(\zeta_m^h,q).
\end{equation}
This is an exact identity, including small indices. Fixing the minimum under conjugation fixes the multiplicity of the largest part, rather than the minimum of the conjugate partition.

\section{An exponential bound on the whole circle}\label{xpc:dv:sec:gap}
\begin{lemma}[First harmonic gap]\label{xpc:dv:lem:gap}
If $0<r<1$, $|q|=r$, and $|\zeta|=1$, then
\begin{equation}\label{xpc:dv:eq:gap}
 |P_\zeta(q)|\le P(r)
  \exp\left(-\frac{r^2(1-\operatorname{Re}\zeta)}{1-r^2}\right).
\end{equation}
\end{lemma}
\begin{proof}
Write $q=r\e^{i\theta}$ and $\zeta=\e^{i\phi}$. Absolute convergence of the logarithmic series gives
\[
 \log P(r)-\log|P_\zeta(q)|
 =\sum_{h\ge1}\frac1h\sum_{j\ge1}r^{jh}
                    \bigl(1-\cos(h\phi+jh\theta)\bigr).
\]
All summands are nonnegative. Retaining $h=1$ gives
$r/(1-r)-\operatorname{Re}(\zeta q/(1-q))$. As $q$ traverses $|q|=r$, the point $q/(1-q)$ traverses a circle with center $r^2/(1-r^2)$ and radius $r/(1-r^2)$. Hence
\[
 \operatorname{Re}\frac{\zeta q}{1-q}
 \le\frac{r^2\cos\phi+r}{1-r^2}.
\]
Subtracting and exponentiating proves the claim.
\end{proof}
This bound is uniform in the argument of $q$. Combining it with Lemma~\ref{xpc:dv:lem:polynomial} and Cauchy's coefficient inequality gives, for $n>k(k+1)/2$,
\begin{equation}\label{xpc:dv:eq:fourier-cauchy}
 \bigl|[q^n]H_k(\zeta,q)\bigr|
 \le r^{-n}P(r)\,r^k\prod_{j<k}(1+r^j)
       \exp\left(-\frac{r^2(1-\operatorname{Re}\zeta)}{1-r^2}\right).
\end{equation}
The bound is valid at every $0<r<1$ and will also be used for finite-$n$ certification.

The classical eta transformation on the positive radius is
\begin{equation}\label{xpc:dv:eq:eta}
 P(\e^{-t})=\sqrt{\frac{t}{2\pi}}
        \exp\left(\frac{\A}{t}-\frac{t}{24}\right)
        P(\e^{-4\pi^2/t}).
\end{equation}
It follows from the modular identity for the Dedekind eta function \cite[\S23.18, eqs.\ (23.18.5)--(23.18.7)]{dlmf}; it is one of the standard partition identities used here. Since $P(x)=1+O(x)$ as $x\to0$, our choice of $t$ yields
\begin{equation}\label{xpc:dv:eq:cauchy-cost}
 \frac{\e^{nt}P(\e^{-t})}{\B(n)}
 =\frac{4\sqrt3\,\A}{\sqrt{2\pi}}\,t^{-3/2}
       \left(1+O(\e^{-4\pi^2/t})\right).
\end{equation}
Also $r^2/(1-r^2)=1/(2t)+O(1)$ at $r=\e^{-t}$. Thus any fixed nontrivial root of unity gives an exponential saving, even after multiplication or division by any fixed power of $t$.

\begin{theorem}[Fixed-minimum residue equidistribution]\label{xpc:dv:thm:residues}
For fixed $k\ge1$, $m\ge2$, and $0\le b<m$,
\begin{align}
 F_k(n)&=(k-1)!\,\B(n)t^{k-1}(1+O_k(t)),\label{xpc:dv:eq:Fk-leading}\\
 \frac{C_{k,m,b}(n)}{F_k(n)}
 &=\frac1m+O_{k,m}\!\left(\e^{-\gamma_m/t}\right),
 \qquad \gamma_m=\frac{1-\cos(2\pi/m)}4>0.\label{xpc:dv:eq:residues}
\end{align}
\end{theorem}
\begin{proof}
The finite-shift transfer in Lemma~\ref{xpc:dv:lem:transfer} below applied to
$q^k(q;q)_{k-1}$ proves \eqref{xpc:dv:eq:Fk-leading}, because its radial expansion begins with $(k-1)!z^{k-1}$. For each nonzero $h$ in \eqref{xpc:dv:eq:filter}, the real gap is at least $1-\cos(2\pi/m)$. Equations \eqref{xpc:dv:eq:fourier-cauchy} and \eqref{xpc:dv:eq:cauchy-cost}, divided by \eqref{xpc:dv:eq:Fk-leading}, leave a fixed power of $t$ times $\exp(-(1-\cos(2\pi/m))/(2t)+O(1))$. Decreasing the exponent constant to $\gamma_m$ absorbs that power. Summing the fixed number of twists proves the result.
\end{proof}
There is no assertion uniform in a growing minimum $k$ or modulus $m$.

\section{The minimum tail and a finite enclosure}\label{xpc:dv:sec:tail}
For $K\ge2$, define the finite polynomial
\begin{equation}\label{xpc:dv:eq:SK}
 S_K(q)=\sum_{k=2}^{K}\frac{k-1}{k}\,q^k(q;q)_{k-1}.
\end{equation}
The root filter with $m=k$, $b=0$ gives the exact decomposition
\begin{equation}\label{xpc:dv:eq:decomposition}
 \sum_{n\ge1}D(n)q^n=P(q)S_K(q)+E_K(q)+T_K(q),
 \qquad E_K(q)=-\sum_{k=2}^{K}\frac1k
                         \sum_{h=1}^{k-1}H_k(\zeta_k^h,q),
\end{equation}
where $T_K$ counts the actual deficit partitions with minimum at least $K+1$. In particular $T_K$ has nonnegative integer coefficients. The polynomial corrections remain included in $E_K$ until its coefficients are bounded.

For $|q|=r$ positivity gives
\begin{equation}\label{xpc:dv:eq:tailcircle}
 |T_K(q)|\le T_K(r)\le(q;q)_K\big|_{q=r}\,P(r)
 =\prod_{j=1}^K(1-r^j)P(r).
\end{equation}
The majorant counts all partitions with no parts $1,\ldots,K$, including an irrelevant empty partition. At $r=\e^{-t}$, Cauchy's inequality and \eqref{xpc:dv:eq:cauchy-cost} therefore prove
\begin{equation}\label{xpc:dv:eq:tailcost}
 [q^n]T_K(q)=O_K\bigl(\B(n)t^{K-3/2}\bigr).
\end{equation}
This explicitly retains the entire Cauchy loss $t^{-3/2}$. A radial Taylor expansion alone would not justify the coefficient estimate.

For $n>K(K+1)/2$, the polynomial corrections in $E_K$ vanish. The fixed finite set of nontrivial twists and \eqref{xpc:dv:eq:fourier-cauchy} make $[q^n]E_K$ exponentially small relative to $\B(n)$. Taking $K=R+4$ for any fixed $R\ge0$ in \eqref{xpc:dv:eq:decomposition} gives
\begin{equation}\label{xpc:dv:eq:finite-reduction}
 D(n)=[q^n]P(q)S_{R+4}(q)+O_R\bigl(\B(n)t^{R+2}\bigr),
\end{equation}
since $K-3/2=R+5/2>R+2$.

There is also the stronger coefficientwise bound
\begin{equation}\label{xpc:dv:eq:positive-tail}
 0\le[q^n]T_K(q)\le u_K(n):=[q^n]P(q)(q;q)_K
 =O_K\bigl(\B(n)t^K\bigr).
\end{equation}
The last equality follows from Lemma~\ref{xpc:dv:lem:transfer}; its leading term is $K!\B(n)t^K$. The weaker whole-circle argument already suffices for \eqref{xpc:dv:eq:finite-reduction}.

\begin{theorem}[Explicit finite-$n$ enclosure]\label{xpc:dv:thm:enclosure}
Let $K\ge2$, $n>K(K+1)/2$, $0<r<1$, and put
\begin{align}
 s_K(n)&=[q^n]P(q)S_K(q),\label{xpc:dv:eq:s}\\
 \mathcal E_K(n,r)&=r^{-n}P(r)\sum_{k=2}^{K}
   \frac{r^k\prod_{j<k}(1+r^j)}{k}
   \sum_{h=1}^{k-1}\exp\left(-\frac{r^2(1-\cos(2\pi h/k))}{1-r^2}\right).
 \label{xpc:dv:eq:Efinite}
\end{align}
Then
\begin{equation}\label{xpc:dv:eq:enclosure}
 s_K(n)-\mathcal E_K(n,r)\le D(n)
 \le s_K(n)+u_K(n)+\mathcal E_K(n,r).
\end{equation}
The quantities $s_K(n)$ and $u_K(n)$ are finite rational combinations of shifted partition numbers.
\end{theorem}
\begin{proof}
Bound $[q^n]E_K$ in absolute value by \eqref{xpc:dv:eq:fourier-cauchy}, and use the coefficientwise bounds $0\le[q^n]T_K\le u_K(n)$ in \eqref{xpc:dv:eq:decomposition}.
\end{proof}
Increasing $K$ need not improve this elementary enclosure: the least Fourier gap becomes smaller. Section~\ref{xpc:dv:sec:certificate} gives an entirely rational evaluation, with no floating exponentials or cosines.

\section{Finite shifts and exact correction coefficients}\label{xpc:dv:sec:transfer}
\subsection{The classical partition input}
The Rademacher series, in the normalization of Johansson \cite[eqs.\ (1.4)--(1.8)]{johansson}, is
\begin{equation}\label{xpc:dv:eq:rademacher}
 p(n)=\frac1{2\sqrt3\,N}\sum_{a\ge1}
       \frac{A_a(n)}{\sqrt a}\,
       U\left(\frac{2\sqrt{\A N}}a\right),
 \qquad U(x)=\cosh x-\frac{\sinh x}{x}.
\end{equation}
The usual Rademacher sum $A_a(n)$ is a sum of unit-modulus phases over reduced residue classes modulo $a$; hence $A_1(n)=1$ and $|A_a(n)|\le a$. No sharper estimate is needed. Since
\[
 U(x)=\frac1x\int_0^x u\sinh u\,\dd u
 \le\frac{x^2\e^x}{3}\qquad(x\ge0),
\]
the $a\ge2$ part is $O(\exp(\sqrt{\A N}))$: after division by $N$ its absolute sum is bounded by a constant times $\exp(\sqrt{\A N})\sum_{a\ge2}a^{-3/2}$. The growing exponential in the $a=1$ term gives
\begin{equation}\label{xpc:dv:eq:partition-smooth}
 p(n)=f(N)+O(\e^{\sqrt{\A N}}),\qquad
 f(N)=\frac{\e^{2\sqrt{\A N}}}{4\sqrt3\,N}
         \left(1-\frac1{2\sqrt{\A N}}\right).
\end{equation}
The decreasing exponential is absorbed in the same error. This estimate holds at every one of finitely many fixed integer shifts as $n\to\infty$.

\subsection{A terminating derivative formula}
For $\ell\ge0$, define
\begin{equation}\label{xpc:dv:eq:Q}
 Q_\ell(t)=\sum_{h=0}^{\ell+1}
 \frac{(-1)^h(\ell+h+1)!}{(\ell+1-h)!\,h!}
       \left(\frac{t}{4\A}\right)^h.
\end{equation}
Direct differentiation gives the exact identity
\begin{equation}\label{xpc:dv:eq:derivative}
 f^{(\ell)}(N)=\B(n)t^\ell Q_\ell(t).
\end{equation}
Indeed, $Q_0(t)=1-t/(2\A)$, $\dd t/\dd N=-t^3/(2\A)$, and
\begin{equation}\label{xpc:dv:eq:Qrec}
 Q_{\ell+1}(t)=\left(1-\frac{(\ell+2)t}{2\A}\right)Q_\ell(t)
                      -\frac{t^2}{2\A}Q'_\ell(t).
\end{equation}
The factorial expression \eqref{xpc:dv:eq:Q} satisfies this recurrence, proving \eqref{xpc:dv:eq:derivative} by induction.

\begin{lemma}[Finite-polynomial transfer]\label{xpc:dv:lem:transfer}
For a fixed polynomial $H(q)$ and a fixed integer $L\ge0$,
\begin{equation}\label{xpc:dv:eq:transfer}
 [q^n]P(q)H(q)
 =\B(n)\sum_{\ell=0}^{L}[z^\ell]H(\e^{-z})\,
                      t^\ell Q_\ell(t)
  +O_{H,L}\bigl(\B(n)t^{L+1}\bigr).
\end{equation}
\end{lemma}
\begin{proof}
Write $H(q)=\sum_{j=0}^{J}h_jq^j$. The coefficient is $\sum_jh_jp(n-j)$. Replace each partition number by $f(N-j)$ using \eqref{xpc:dv:eq:partition-smooth}, with an exponentially small total error. Taylor-expand this explicit smooth function at $N$ to degree $L$. By \eqref{xpc:dv:eq:derivative}, its $(L+1)$st derivative throughout each bounded interval $[N-J,N]$ is $O_{H,L}(\B(n)t^{L+1})$. The Taylor remainder has the stated size. Finally,
\[
 \sum_{j=0}^Jh_j\frac{(-j)^\ell}{\ell!}=[z^\ell]H(\e^{-z}).
\]
Only the explicit function $f$ is differentiated; no unknown error term on the integer sequence $p(n)$ is differentiated.
\end{proof}

\subsection{The formal coefficient algorithm}
Define a formal power series
\begin{equation}\label{xpc:dv:eq:G}
 G(z)=\sum_{k\ge2}\frac{k-1}{k}\,\e^{-kz}
                         \prod_{j=1}^{k-1}(1-\e^{-jz}),
 \qquad g_\ell=[z^\ell]G(z).
\end{equation}
The $k$th term has valuation $k-1$, so any specified coefficient involves finitely many rational operations. Finite telescoping using
$q^k(q;q)_{k-1}=(q;q)_{k-1}-(q;q)_k$ proves the equivalent formal identity
\begin{equation}\label{xpc:dv:eq:Gtail}
 G(z)=\sum_{j\ge1}\frac{\prod_{h=1}^{j}(1-\e^{-hz})}{j(j+1)}.
\end{equation}
No analytic convergence near $z=0$ is asserted. Define
\begin{align}
 \sum_{m\ge0}c_mt^m
 &=\frac2t\sum_{\ell\ge1}g_\ell t^\ell Q_\ell(t),\label{xpc:dv:eq:c-generator}\\
 c_m&=2\!\sum_{\substack{1\le\ell\le m+1\\ h=m+1-\ell\le\ell+1}}
 g_\ell\frac{(-1)^h(\ell+h+1)!}
              {(\ell+1-h)!\,h!\,(4\A)^h}.\label{xpc:dv:eq:c-finite}
\end{align}
The displayed sum always has $h\ge0$. Each coefficient belongs to $\mathbb Q[\A^{-1}]$ and is effectively computable.

\begin{proof}[Proof of Theorem~\ref{xpc:dv:thm:main}]
Use \eqref{xpc:dv:eq:finite-reduction} with fixed $R$, then Lemma~\ref{xpc:dv:lem:transfer} with $H=S_{R+4}$ and $L=R+1$. The coefficients of $S_{R+4}(\e^{-z})$ through degree $R+1$ equal $g_\ell$, because every omitted minimum has larger valuation. Terms of $t$-degree above $R+1$ in the finite sum contribute $O_R(\B(n)t^{R+2})$. Equations \eqref{xpc:dv:eq:c-generator}--\eqref{xpc:dv:eq:c-finite} now give \eqref{xpc:dv:eq:main}. The first radial coefficient is $g_1=1/2$, so $c_0=1$. Since $p(n)=\B(n)(1-t/(2\A))+O(\e^{\sqrt{\A N}})$ and $t\sim\pi/\sqrt{6n}$, the leading equivalent follows.
\end{proof}

The first eight radial coefficients are
\[
 (g_1,\ldots,g_8)=\left(
 \frac12,\frac1{12},\frac1{12},\frac{89}{720},
 \frac{61}{240},\frac{20287}{30240},\frac{21841}{10080},
 \frac{10023809}{1209600}\right).
\]
The first six relative coefficients in \eqref{xpc:dv:eq:main} are
\begin{align*}
 c_0&=1,\\
 c_1&=\frac16-\frac3{2\A},\\
 c_2&=\frac16-\frac1{2\A}+\frac3{4\A^2},\\
 c_3&=\frac{89}{360}-\frac5{6\A}+\frac5{8\A^2},\\
 c_4&=\frac{61}{120}-\frac{89}{48\A}+\frac{15}{8\A^2}-\frac5{16\A^3},\\
 c_5&=\frac{20287}{15120}-\frac{427}{80\A}
                         +\frac{623}{96\A^2}-\frac{35}{16\A^3}.
\end{align*}
Relative to the exact partition number, divide the formal series by
$1-t/(2\A)$ to obtain
\begin{align}
 \frac{D(n)}{p(n)}&=\frac t2
  \left(1+\rho_1t+\rho_2t^2+\rho_3t^3+O(t^4)\right),\label{xpc:dv:eq:relative-p}\\
 \rho_1&=\frac16-\frac1\A,\notag\\
 \rho_2&=\frac16-\frac5{12\A}+\frac1{4\A^2},\notag\\
 \rho_3&=\frac{89}{360}-\frac3{4\A}
                         +\frac5{12\A^2}+\frac1{8\A^3}.\notag
\end{align}
The shift $N=n-1/24$ remains part of these formulas; replacing $t$ by $\pi/\sqrt{6n}$ without re-expanding changes higher coefficients.

\begin{writenote}[Part~I's counterpart, 7 October 2026]
This section is, for the same partition function, Part~I's
Section~\ref{xpc:sl:sec:transfer}: \eqref{xpc:dv:eq:partition-smooth} is
Lemma~\ref{xpc:sl:lem:gap} (both from Johansson's normalization of the
Rademacher series, with the same bound on $U$), \eqref{xpc:dv:eq:Q},
\eqref{xpc:dv:eq:derivative} and~\eqref{xpc:dv:eq:Qrec} are
Lemma~\ref{xpc:sl:lem:Q} (the same polynomials $Q_\ell$), and
Lemma~\ref{xpc:dv:lem:transfer} is Proposition~\ref{xpc:sl:prop:transfer};
the eta identity~\eqref{xpc:dv:eq:eta} is~\eqref{xpc:sl:eq:eta}. It is kept
because the later sections of this Part cite these statements by their own
labels; the proofs are the same argument.
\end{writenote}

\section{The reciprocal-minimum weighted count}\label{xpc:dv:sec:weighted}
Define
\[
 W(n)=\sum_{\lambda\vdash n}\left(1-\frac1{\min\lambda}\right)
 \qquad(n\ge1).
\]
This rational number is a weighted count, rather than an integer count of a subset.
\begin{corollary}\label{xpc:dv:cor:weight}
For every fixed real $L>0$,
\begin{equation}\label{xpc:dv:eq:superalgebraic}
 D(n)-W(n)=O_L\bigl(p(n)n^{-L}\bigr).
\end{equation}
\end{corollary}
\begin{proof}
Choose a fixed integer $K\ge2L+2$. For minima $k\le K$, the difference between the deficit count and the weight $(k-1)/k$ times the fixed-minimum count is a finite sum of nontrivial Fourier coefficients. These are exponentially small relative to $p(n)$ by Section~\ref{xpc:dv:sec:gap}. Each remaining nonnegative tail, for $D$ or $W$, is at most $u_K(n)$, because its weights lie in $[0,1]$. Equation \eqref{xpc:dv:eq:positive-tail} bounds their sum by $O_K(p(n)n^{-K/2})$. This implies \eqref{xpc:dv:eq:superalgebraic}.
\end{proof}
The minimum and its associated modulus range over all integers in the complete sum. This fixed-cutoff proof alone gives no specified decay rate beyond all algebraic orders. Section~\ref{xpc:dv:sec:quantitative} strengthens it using exact uniform inequalities at a growing auxiliary cutoff.

For completeness, the exact weighted generating function is
\begin{equation}\label{xpc:dv:eq:Hweight}
 \sum_{n\ge1}W(n)q^n=P(q)H(q),\qquad
 H(q)=\sum_{k\ge2}\frac{k-1}{k}q^k(q;q)_{k-1}.
\end{equation}
For $|q|<1$, ordinary telescoping gives
\begin{equation}\label{xpc:dv:eq:Htail}
 H(q)=\sum_{j\ge1}\frac{(q;q)_j-(q;q)_\infty}{j(j+1)}
 =\sum_{j\ge1}\frac{(q;q)_j}{j(j+1)}-(q;q)_\infty.
\end{equation}
The last product is indispensable in the analytic identity, even though it has no contribution to the formal radial series \eqref{xpc:dv:eq:Gtail}. Nor is $P(q)H(q)$ the exact deficit generating function: the Fourier terms still distinguish the two.

Here is a direct derivation of the connection to the arbitrary-weight tail framework mentioned in Section~\ref{xpc:dv:sec:prior}. Define
\[
 g(x)=\sum_{j\ge1}\frac{x^j}{j(j+1)}
     =1+(x^{-1}-1)\log(1-x),\quad |x|<1,
\]
with the removable value $g(0)=0$ and logarithm given by its power series. Euler's identity gives
\[
 (q;q)_j-(q;q)_\infty
 =(q;q)_\infty\sum_{v\ge1}\frac{q^{v(j+1)}}{(q;q)_v}.
\]
For $|q|\le r<1$, finite products are bounded above by
$\prod_{h\ge1}(1+r^h)$ and their reciprocals by $P(r)$. The double series of absolute values is bounded by a constant times
\[
 \sum_{j\ge1}\frac1{j(j+1)}\sum_{v\ge1}r^{v(j+1)}<\infty.
\]
Interchanging sums is therefore justified and yields
\begin{equation}\label{xpc:dv:eq:abel}
 H(q)=(q;q)_\infty\sum_{v\ge1}\frac{q^v g(q^v)}{(q;q)_v}.
\end{equation}
Since $q^vg(q^v)=q^v+(1-q^v)\log(1-q^v)$, this also proves
\begin{equation}\label{xpc:dv:eq:log-identity}
 P(q)H(q)=P(q)-1+
           \sum_{v\ge1}\frac{\log(1-q^v)}{(q;q)_{v-1}}.
\end{equation}
These are exact analytic identities for the weighted statistic. They are not substitutes for the coefficient error estimates that prove Corollary~\ref{xpc:dv:cor:weight}.

\section{A quantitative reciprocal-minimum approximation}\label{xpc:dv:sec:quantitative}
The exact estimates proved earlier also permit a specified growing auxiliary cutoff. This yields a stronger global comparison than Corollary~\ref{xpc:dv:cor:weight}, without using a coefficient-transfer theorem uniform in that cutoff.

\begin{theorem}[Explicit stretched-exponential comparison]\label{xpc:dv:thm:quantitative}
Put
\[
 c_* =\left(\frac{\pi\sqrt6}{9}\right)^{1/3},\qquad
 C_0=\frac{32\sqrt3\,\A}{\sqrt{2\pi}},\qquad
 \Lambda=\log(N/\A).
\]
For every integer $n\ge90$,
\begin{equation}\label{xpc:dv:eq:explicit-stretched}
 \frac{|D(n)-W(n)|}{p(n)}
 \le C_0\exp\left(
 -c_*N^{1/6}\Lambda^{2/3}
 +3c_*\log2\,N^{1/6}\Lambda^{-1/3}
 +\frac54\Lambda\right).
\end{equation}
Consequently, with an absolute implied constant,
\begin{equation}\label{xpc:dv:eq:stretched}
 \frac{|D(n)-W(n)|}{p(n)}
 \le\exp\left(-c_*n^{1/6}(\log n)^{2/3}
               +O\bigl(n^{1/6}(\log n)^{-1/3}\bigr)\right).
\end{equation}
For each fixed $\varepsilon>0$, the right side can therefore be replaced, for all sufficiently large $n$, by
$\exp(-(c_*-\varepsilon)n^{1/6}(\log n)^{2/3})$.
\end{theorem}
The explicit bound may exceed one at small indices; it remains valid. It is an upper bound, not an asymptotic formula for the actual error. In particular it does not assert a rate $\exp(-c\sqrt n)$ for any fixed $c>0$.

\subsection{An exact bound uniform in the cutoff}
For any integer $K\ge2$, $n>K(K+1)/2$, and $0<r<1$, define
\[
 \Gamma_K(r)=\frac{r^2(1-\cos(2\pi/K))}{1-r^2}.
\]
Then the following bound is simultaneous in all its displayed parameters:
\begin{equation}\label{xpc:dv:eq:uniform-cutoff}
 |D(n)-W(n)|\le r^{-n}P(r)
 \left(2^K\e^{-\Gamma_K(r)}+\prod_{j=1}^{K}(1-r^j)\right).
\end{equation}
To prove it, let $C_k(n)=C_{k,k,0}(n)$. The minimum-$k$ contribution to $D-W$ is $F_k(n)/k-C_k(n)$. For $2\le k\le K$, the root filter and polynomial reduction give the negative sum of the nontrivial twists, divided by $k$. Every polynomial coefficient has vanished under the stated bound on $n$. All nontrivial $k$th roots have
$1-\operatorname{Re}\zeta\ge1-\cos(2\pi/K)$, because their angular distance from $1$ is at least $2\pi/k\ge2\pi/K$. Hence \eqref{xpc:dv:eq:fourier-cauchy} bounds the total small-minimum error by
\[
 r^{-n}P(r)\e^{-\Gamma_K(r)}
 \sum_{k=2}^{K}\frac{k-1}{k}r^k\prod_{j<k}(1+r^j)
 <r^{-n}P(r)2^K\e^{-\Gamma_K(r)}.
\]
For a partition with minimum $k>K$, its signed weight in $D-W$ is
$1/k-\mathbf1_{k\mid\max\lambda}$, whose absolute value is at most one. Thus the absolute tail is bounded by the restricted-partition coefficient $u_K(n)$ and then by its positive-radius Cauchy majorant. This proves \eqref{xpc:dv:eq:uniform-cutoff}, with only one tail majorant. At $r=\e^{-t}$, it gives
\begin{equation}\label{xpc:dv:eq:uniform-two-term}
 |D(n)-W(n)|\le \e^{nt}P(\e^{-t})
 \left(
 2^K\exp\left(-\frac{1-\cos(2\pi/K)}{\e^{2t}-1}\right)
 +t^KK!\right),
\end{equation}
since $1-\e^{-jt}\le jt$. No parameter-dependent asymptotic error occurs in this inequality.

\subsection{An explicit normalized Cauchy cost}
When $t\le\e^{-2}$, the exact eta identity \eqref{xpc:dv:eq:eta} and the Rademacher series give
\begin{equation}\label{xpc:dv:eq:explicit-cost}
 \frac{\e^{nt}P(\e^{-t})}{p(n)}
 \le C_1t^{-3/2},\qquad C_1=\frac{16\sqrt3\,\A}{\sqrt{2\pi}}.
\end{equation}
Here are explicit bounds establishing the constant and its onset. Set $w=2\A/t$. The tail estimate for \eqref{xpc:dv:eq:rademacher}, using
$\sum_{a\ge2}a^{-3/2}\le\int_1^\infty x^{-3/2}\dd x=2$, gives
\begin{equation}\label{xpc:dv:eq:p-lower-explicit}
 \frac{p(n)}{\B(n)}\ge1-\frac1w-\frac{4w^2}{3}\e^{-w/2}.
\end{equation}
The decreasing exponential in the $a=1$ summand is positive and has been discarded. At $t\le\e^{-2}$, elementary rational bounds imply
\[
 w\ge2\A\e^2>
 \frac{(25/8)^2(163/60)^2}{3}
 =\frac{664225}{27648}>24.
\]
Since $w^2\e^{-w/2}$ decreases for $w\ge4$,
\[
 \frac{p(n)}{\B(n)}\ge1-\frac1{24}-768\e^{-12}
 >1-\frac1{24}-768\left(\frac38\right)^{12}>\frac12.
\]
For $0<s\le1/10$,
\[
 \log P(s)\le\sum_{j\ge1}\frac{s^j}{1-s^j}
 \le\frac{s}{(1-s)^2}\le\frac{10}{81}<\frac12<\log2.
\]
Taking $s=\e^{-4\pi^2/t}<1/10$ gives $P(s)<2$. The last strict bound follows already from $4\pi^2/t>36>3$ and $\e^3>(8/3)^3>10$. Combining these estimates with the exact identity
\[
 \frac{\e^{nt}P(\e^{-t})}{\B(n)}
 =\frac{4\sqrt3\,\A}{\sqrt{2\pi}}\,t^{-3/2}
          P(\e^{-4\pi^2/t})
\]
proves \eqref{xpc:dv:eq:explicit-cost}.

The requirement $t\le\e^{-2}$ holds for every $n\ge90$, because
\begin{equation}\label{xpc:dv:eq:ninety}
 \A\e^4+\frac1{24}
 <\frac{(22/7)^2(87/32)^4}{6}+\frac1{24}
 =\frac{6935272345}{77070336}<90.
\end{equation}
For clarity, the elementary constants used here admit short exact certificates. Machin's identity
$\pi=16\arctan(1/5)-4\arctan(1/239)$ and alternating arctangent series give
$25/8<\pi<22/7$, hence $9<\pi^2<10$. Also
\[
 \frac{163}{60}=\sum_{j=0}^{5}\frac1{j!}<\e
 <\sum_{j=0}^{6}\frac1{j!}
       +\frac{1/7!}{1-1/8}<\frac{87}{32}.
\]
The upper tail bound follows because every subsequent ratio is at most $1/8$. These comparisons, including \eqref{xpc:dv:eq:ninety}, are regenerated with rational arithmetic in the accompanying code. The inequality $\log2>1/2$ follows directly from $\int_1^2\dd u/u>1/2$.

\subsection{Balancing a growing integer cutoff}
Put
\[
 L=\log(1/t)\ge2,\quad
 S=t^{-1/3}L^{2/3},\quad
 \kappa=\left(\frac{3\pi^2}{2}\right)^{1/3},\quad
 \beta_0=\left(\frac{4\pi^2}{9}\right)^{1/3}=\frac{2\kappa}{3},
\]
and choose
\begin{equation}\label{xpc:dv:eq:cutoff-choice}
 x=\kappa(tL)^{-1/3},\qquad K=\floor{x}.
\end{equation}
One has $x>3$ and $tx<1$. Indeed $x^3=(3\pi^2/2)\e^L/L$ increases for $L\ge2$, whereas $(tx)^3=(3\pi^2/2)\e^{-2L}/L$ decreases. At $L=2$, the first is greater than $48>27$, and the second is less than $(30/4)(3/8)^4<1$. Thus $K\ge3$, $K<1/t$, and
\begin{equation}\label{xpc:dv:eq:degree-growing}
 \frac{K(K+1)}2<\frac{t^{-2}+t^{-1}}2<t^{-2}=\frac N\A<n.
\end{equation}
All polynomial corrections therefore vanish even for this growing cutoff. This is a direct check of the exact inequality's domain, rather than a uniform extension of Lemma~\ref{xpc:dv:lem:transfer}.

For the twist term, $K\le x$ implies
$1-\cos(2\pi/K)\ge1-\cos(2\pi/x)$. Use
\[
 1-\cos u\ge\frac{u^2}2-\frac{u^4}{24},\qquad
 \frac1{\e^{2t}-1}\ge\frac1{2t}-\frac12.
\]
The first follows, for example, from the alternating cosine series for the present $u=2\pi/x<2\pi/3$; the second follows from $\coth t\ge1/t$, since $t\cosh t-\sinh t$ has nonnegative derivative and vanishes at zero. Both right sides are nonnegative here: $u^2<4\pi^2/9<12$ and $t<1$. Multiplying and dropping the positive final term gives
\begin{equation}\label{xpc:dv:eq:gap-growing}
 \Gamma_K(\e^{-t})\ge\frac{\pi^2}{tx^2}-\delta(t)
 =\beta_0S-\delta(t),\qquad
 \delta(t)=\frac{\pi^2}{x^2}+\frac{\pi^4}{3tx^4}.
\end{equation}
Hence
\begin{equation}\label{xpc:dv:eq:twist-growing}
 2^K\e^{-\Gamma_K(\e^{-t})}
 \le\exp(-\beta_0S+x\log2+\delta(t)).
\end{equation}
There is no omitted loss from rounding down: that rounding improves the lower bound on the angular gap.

To bound $\delta$, write
\[
 \frac{\delta(t)}L
 =\frac{\pi^2}{\kappa^2}\e^{-2L/3}L^{-1/3}
 +\frac{\pi^4}{3\kappa^4}\e^{-L/3}L^{1/3}.
\]
Both summands decrease for $L\ge2$; their logarithmic derivatives are
$-2/3-1/(3L)$ and $-1/3+1/(3L)$. At $L=2$, their cubes are respectively
\begin{align*}
 \frac{2\pi^2}{9}\e^{-4}
 &<\frac{20}{9}\left(\frac38\right)^4
 <\left(\frac9{25}\right)^3,\\
 \frac{32\pi^4}{2187}\e^{-2}
 &<\frac{3200}{2187}\left(\frac38\right)^2
 <\left(\frac35\right)^3.
\end{align*}
Thus $\delta(t)/L<9/25+3/5=24/25<1$, and $\delta(t)\le L$.

For the minimum tail set
\[
 H=-\log(tx)=\frac23L+\frac13\log L-\log\kappa>0.
\]
Using $K!\le K^K$, $K\le x$, and $K\ge x-1$ gives
\[
 \log(t^KK!)\le K\log(tK)\le-KH\le-(x-1)H.
\]
Since $L\ge2>3\pi^2/16=\kappa^3/8$, one has
$H\ge(2/3)L-\log2$; also $H=L-\log x\le L$. It follows that
\begin{equation}\label{xpc:dv:eq:tail-growing}
 \log(t^KK!)\le-\beta_0S+x\log2+L.
\end{equation}
This accounts explicitly for the possible floor-rounding loss, at most $H\le L$.

Equations \eqref{xpc:dv:eq:twist-growing}--\eqref{xpc:dv:eq:tail-growing} and the normalized Cauchy cost give, for every $n\ge90$,
\begin{equation}\label{xpc:dv:eq:stretched-t}
 \frac{|D(n)-W(n)|}{p(n)}
 \le2C_1\exp\left(-\beta_0S+\kappa\log2\,\frac SL+\frac52L\right).
\end{equation}
Now $L=\Lambda/2$ and
\[
 \beta_0S=c_*N^{1/6}\Lambda^{2/3},\qquad
 \kappa\log2\,\frac SL=3c_*\log2\,N^{1/6}\Lambda^{-1/3}.
\]
Indeed $c_*^3=\beta_0^3\A^{-1/2}/4=\pi^2/(9\sqrt\A)=\pi\sqrt6/9$.
Since $C_0=2C_1$, this proves \eqref{xpc:dv:eq:explicit-stretched}. Finally
\[
 N^{1/6}\Lambda^{2/3}
 =n^{1/6}(\log n)^{2/3}
       +O\bigl(n^{1/6}(\log n)^{-1/3}\bigr).
\]
The other terms in \eqref{xpc:dv:eq:explicit-stretched}, including $O(\log n)$ and $\log C_0$, fit in this error scale. This proves \eqref{xpc:dv:eq:stretched} and the eventual $\varepsilon$ statement, completing the proof of Theorem~\ref{xpc:dv:thm:quantitative}.

For a chosen $\varepsilon>0$ and a particular integer $n\ge90$, comparing the logarithm of the right side of \eqref{xpc:dv:eq:explicit-stretched} with
$-(c_*-\varepsilon)n^{1/6}(\log n)^{2/3}$ is a computable pointwise sufficient test. The first passing integer is not asserted to certify every larger integer or an optimal onset for the true error. Eventual validity follows from \eqref{xpc:dv:eq:stretched}.

\subsection{Optimality only for this two-term majorant}
The balance above optimizes the leading constant of the specific majorant
\[
 M(t,K)=2^K\e^{-\Gamma_K(\e^{-t})}+t^KK!.
\]
If $K\sim a(tL)^{-1/3}$ for a fixed $a>0$, elementary small-argument expansions and the logarithmic Stirling formula yield
\[
 \frac{\log(2^K\e^{-\Gamma_K})}{S}\longrightarrow-\frac{\pi^2}{a^2},
 \qquad
 \frac{\log(t^KK!)}S\longrightarrow-\frac{2a}{3}.
\]
For the second limit, use
$\log(tK)=-(2/3)L-(1/3)\log L+\log a+o(1)$ and
$\log K!=K\log K-K+O(\log K)$. Therefore
\[
 \frac{-\log M(t,K)}S\longrightarrow
       \min\left\{\frac{\pi^2}{a^2},\frac{2a}{3}\right\},
\]
whose unique maximum is $\beta_0$ at $a=\kappa$.

This also bounds arbitrary admissible integer-cutoff sequences with $K\ge2$. Write
$a(t)=K(tL)^{1/3}$ and $R(t)=-\log M(t,K)/S$. Since each term of $M$ is positive,
\[
 R(t)\le a(t),\qquad R(t)\le\frac{\pi^2}{a(t)^2}.
\]
The first uses $\log K!\ge0$; the second uses
$1-\cos u\le u^2/2$ and $\e^{2t}-1\ge2t$ and discards $K\log2\ge0$.
If $\limsup R>\beta_0$, a subsequence with $R\ge\beta_0+\eta$ confines $a(t)$ to a compact subset of $(0,\infty)$. A convergent further subsequence contradicts the preceding fixed-$a$ limit. Thus no cutoff sequence improves this leading constant within $M$; the choice \eqref{xpc:dv:eq:cutoff-choice} attains it. The logarithm of the Cauchy cost is $O(L)=o(S)$ and does not alter the conclusion.

This does not optimize the actual error, the sharper product tail in \eqref{xpc:dv:eq:uniform-cutoff}, a more detailed Fourier sum, or another decomposition. No absence of cancellation or actual-error asymptotic is proved, and subleading cutoff terms have not been optimized.

\section{All-order inverses and integer thresholds}\label{xpc:dv:sec:inverse}
\subsection{The common formal reversion}
For $X=D$ or $X=a$, use the weak threshold convention
\[
 \nu_X(y)=\min\{n\ge1:X(n)\ge y\}.
\]
Set $w=2\sqrt{\A(n-1/24)}$. Theorem~\ref{xpc:dv:thm:main} gives, at every fixed order,
\begin{equation}\label{xpc:dv:eq:forward-w}
 X(n)=C\e^w w^{-\beta}
        \left(1+\sum_{j=1}^{R}d_jw^{-j}+O_R(w^{-R-1})\right).
\end{equation}
For the deficit,
\begin{equation}\label{xpc:dv:eq:Ddata}
 C_D=\frac{\A^2}{\sqrt3},\qquad \beta_D=3,
 \qquad d_j^{(D)}=(2\A)^jc_j.
\end{equation}
In particular,
\[
 d_1^{(D)}=\frac\A3-3,\qquad
 d_2^{(D)}=\frac{2\A^2}3-2\A+3,\qquad
 d_3^{(D)}=\frac{\A(89\A^2-300\A+225)}{45}.
\]
For $a=p-D$,
\begin{equation}\label{xpc:dv:eq:adata}
 C_a=\frac\A{\sqrt3},\quad\beta_a=2,\quad
 d_1^{(a)}=-(\A+1),\quad
 d_j^{(a)}=-\A d_{j-1}^{(D)}\quad(j\ge2).
\end{equation}
Indeed $p/\B=1-1/w$ to all algebraic orders, whereas
$D/\B=(\A/w)(1+\sum_{j\ge1}d_j^{(D)}w^{-j})$.

For the relevant $C,\beta$, define
\begin{equation}\label{xpc:dv:eq:lambert}
 w_0=-\beta W_{-1}\!\left(-\frac{(C/y)^{1/\beta}}\beta\right).
\end{equation}
For $y>C(\e/\beta)^\beta$, this is the unique solution $w_0>\beta$ of
$C\e^{w_0}w_0^{-\beta}=y$. The lower real Lambert branch $W_{-1}$ is essential; the principal branch describes the other side of the minimum.

With $u=1/w_0$, seek a formal series $\Delta(u)=\sum_{j\ge1}v_ju^j$ satisfying
\begin{equation}\label{xpc:dv:eq:inverse-formal}
 \Delta-\beta\log(1+u\Delta)
 +\log\left(1+\sum_{j\ge1}d_j
                  \left(\frac{u}{1+u\Delta}\right)^j\right)=0.
\end{equation}
At degree $u^m$, the unknown $v_m$ occurs with coefficient one only in $\Delta$; elsewhere its occurrence gains at least one additional $u$. Therefore the solution is unique and triangularly computable. The first coefficients are
\begin{align}
 v_1&=-d_1,\notag\\
 v_2&=-\beta d_1+\frac{d_1^2}{2}-d_2,\label{xpc:dv:eq:v}\\
 v_3&=-\beta^2d_1+\left(\frac\beta2-1\right)d_1^2
       -\beta d_2-\frac{d_1^3}{3}+d_1d_2-d_3.\notag
\end{align}
For an integer $M\ge1$, retain the square in full and define
\begin{equation}\label{xpc:dv:eq:xM}
 x_M(y)=\frac1{24}+\frac1{4\A}
          \left(w_0+\sum_{j=1}^{M}\frac{v_j}{w_0^j}\right)^2.
\end{equation}

\begin{theorem}[Asymptotic integer inverse bracket]\label{xpc:dv:thm:inverse}
For $X=D$ or $X=a$ and every fixed $M\ge1$, there are constants $B_M>0$ and $y_M$ such that
\begin{equation}\label{xpc:dv:eq:inverse-bracket}
 \ceil{x_M(y)-B_Mw_0^{-M}}
 \le\nu_X(y)\le
 \ceil{x_M(y)+B_Mw_0^{-M}}\qquad(y\ge y_M).
\end{equation}
The constants and onset are not made effective here.
\end{theorem}
\begin{proof}
Use the truncated smooth forward function
\[
 F_M(x)=C\e^w w^{-\beta}
         \left(1+\sum_{j=1}^{M}d_jw^{-j}\right),
 \qquad w=2\sqrt{\A(x-1/24)}.
\]
It is positive and increasing for sufficiently large $x$. Its logarithmic derivative is
\[
 \frac{\dd}{\dd x}\log F_M(x)
 =\frac{2\A}{w}\left(1-\frac\beta w+O(w^{-2})\right).
\]
Substituting the truncated $\Delta$ into the logarithmic equation gives residual $O(w_0^{-M-1})$. The derivative of that equation with respect to $w$ is $1+O(w_0^{-1})$. Thus the error in the $w$ root is $O(w_0^{-M-1})$, and $\dd x/\dd w=w/(2\A)$ turns this into an $O(w_0^{-M})$ error in the smooth root.

At nearby integers the relative forward error in \eqref{xpc:dv:eq:forward-w} is $O(w_0^{-M-1})$. Moving the smooth root by $h=Bw_0^{-M}$ changes its logarithmic value by at least a fixed positive constant times $Bw_0^{-M-1}$. Choosing $B$ large enough dominates that error. If $x_*$ is the exact large root of $F_M(x)=y$, then
\[
 n_-:=\ceil{x_*-h}-1,\qquad n_+:=\ceil{x_*+h}
 \quad\hbox{satisfy}\quad X(n_-)<y\le X(n_+).
\]
The same conclusion remains valid if either test integer is farther from the root, since the smooth function is increasing and the displacement is still at most $1+h$.

The sequence $a$ is globally weakly increasing by \eqref{xpc:dv:eq:sandwich}. For $D$, use \eqref{xpc:dv:eq:main} through relative order one, including an $O(t^2)$ remainder. The leading factor $\B(n)t/2$ has ratio $1+t+O(t^2)$ at consecutive indices; the ratio of its remaining factors is $1+O(t^2)$, by comparing two bounded remainders directly. Consequently
\[
 \frac{D(n+1)}{D(n)}=1+t+O(t^2)>1
\]
eventually. No error term is differentiated. Also $D$ is eventually positive and unbounded. The finite initial segment is below all sufficiently large $y$, so the two test integers bracket the first threshold for either $X$. Absorb the smooth-root error into $B_M$ to obtain \eqref{xpc:dv:eq:inverse-bracket}.
\end{proof}
The interval before taking ceilings shrinks to zero. Its endpoint ceilings are eventually equal or adjacent, but targets arbitrarily near an integer boundary prevent an unconditional prescription $\nu_X(y)=\ceil{x_M(y)}$. The theorem is not a certified finite-$y$ numerical error bound. Deficit monotonicity is only eventual: exact counting gives $D(17)=44$ and $D(18)=43$.

\subsection{Explicit inverse corrections}
For the deficit the first two coefficients are
\[
 v_1^{(D)}=3-\frac\A3,\qquad
 v_2^{(D)}=\frac{21}{2}-\frac{11\A^2}{18}.
\]
Thus a smooth inverse center has expansion
\begin{equation}\label{xpc:dv:eq:inverse-D}
 x_D(y)=\frac1{24}+\frac{w_0^2}{4\A}
       +\frac3{2\A}-\frac16
       +\frac1{w_0}\left(\frac{21}{4\A}-\frac{11\A}{36}\right)
       +O(w_0^{-2}),
\end{equation}
with the deficit core from \eqref{xpc:dv:eq:Ddata}. For $a$,
\[
 v_1^{(a)}=\A+1,\qquad v_2^{(a)}=\frac{5\A^2}{6}+\frac52,
\]
and its own core from \eqref{xpc:dv:eq:adata} gives
\begin{equation}\label{xpc:dv:eq:inverse-a}
 x_a(y)=\frac1{24}+\frac{w_0^2}{4\A}
       +\frac12+\frac1{2\A}
       +\frac1{w_0}\left(\frac{5\A}{12}+\frac5{4\A}\right)
       +O(w_0^{-2}).
\end{equation}
These are expansions of smooth centers in the sense of \eqref{xpc:dv:eq:xM}, not integer thresholds.

\begin{proposition}[Exact two-index shortcut]\label{xpc:dv:prop:shortcut}
Let $\nu_p(y)=\min\{n\ge0:p(n)\ge y\}$. For every real $y>1$,
\begin{equation}\label{xpc:dv:eq:shortcut}
 \nu_p(y)\le\nu_a(y)\le\nu_p(y)+1.
\end{equation}
\end{proposition}
\begin{proof}
Set $j=\nu_p(y)$. For $n<j$, $a(n)\le p(n)<y$, while
$a(j+1)\ge p(j)\ge y$. Since $y>1=p(0)=p(1)$, no zero-index convention interferes.
\end{proof}
Hence evaluating $a(\nu_p(y))$ selects the correct one of the two candidates once the partition threshold is known.

The smooth inverses for $p$ and $a$ have the same $C=\A/\sqrt3$, $\beta=2$, and the same Lambert core. For $p$, $d_1=-1$ and $d_j=0$ for $j\ge2$, so $v_1=1$, $v_2=5/2$. Subtracting the corresponding centers gives
\begin{equation}\label{xpc:dv:eq:half-index}
 x_a(y)-x_p(y)=\frac12+\frac{5\A}{12w_0}+O(w_0^{-2}).
\end{equation}
The half-index shift concerns smooth centers. The exact discrete statement is \eqref{xpc:dv:eq:shortcut}.

\begin{remark}[{[write]} The inverses and the transseries volume, 7~October 2026]\label{xpc:dv:rem:transseries}
Statement by statement, against the volume named in the Guide (its symbols
carry the subscript ``vol'').
\begin{enumerate}[label=(\alph*),leftmargin=*]
\item \emph{Instances, and strict increase of $a$.} For $X=D$ and $y$ beyond
the initial segment, $\nu_D(y)$ is the integer staircase of
\file{p0:def:three-inverses}, so \file{p0:thm:staircase}(1) applies (in the
shipped table $D$ increases strictly for $18\le n<1000$, a finite
observation). For $X=a$ the write proves that $a$ is strictly increasing for
every $n\ge0$, so $\nu_a$ is the staircase with $n_1=0$. \emph{Proof.} By
\eqref{xpc:dv:eq:sandwich}, $a(n+1)\ge p(n)=a(n)+D(n)$, so it suffices that
$D(n)\ge1$. For odd $n\ge5$ the partition $(n-2,2)$ is counted by $D(n)$
($n-2$ is odd); for even $n\ge8$ so is $(n-5,3,2)$ ($n-5\ge3$ is odd and the
minimum is~$2$). Hence $a(n+1)>a(n)$ for $n=5$ and all $n\ge7$, and the
values $a(0),\dots,a(8)=0,1,2,3,5,6,11,12,20$ settle $n\le7$. (This
strengthens the weak monotonicity stated after~\eqref{xpc:dv:eq:sandwich};
note $D(6)=0$.)
\item \emph{Instance.} The Lambert root~\eqref{xpc:dv:eq:lambert} is
\file{p0:thm:lambert-core}(3), case $b_\vol<0$, with $X_\vol=w_0$, $a_\vol=1$,
$b_\vol=-\beta$, $L_\vol=\log(y/C)$; the condition $y>C(e/\beta)^\beta$ is the
fold condition.
\item \emph{Formal instance.} The equation~\eqref{xpc:dv:eq:inverse-formal} is
the formal-exactness identity of \file{p0:thm:lambert-centered}(4) with
$a_\vol=1$, $b_\vol=-\beta$, $Q_\vol(\tau)=\log(1+\sum_jd_j\tau^j)$; it is
Part~I's~\eqref{xpc:sl:eq:reversion} with $d=\beta$ and $d_j=(2\mathsf A)^jc_j$,
so $v_m$ is Part~I's $u_m$, and \eqref{xpc:dv:eq:v} is
\eqref{xpc:sl:eq:u1}--\eqref{xpc:sl:eq:u3} in these variables (checked
symbolically). The smooth root in the proof of Theorem~\ref{xpc:dv:thm:inverse}
is \file{p0:thm:lambert-centered}(5) in $w$ (the truncated $F_M$ is exactly of
exponential--power type in $w$), followed by $x=\tfrac1{24}+w^2/(4\mathsf A)$;
its residual-over-slope step is \file{p0:thm:backward-error}.
\item \emph{Analogue.} Theorem~\ref{xpc:dv:thm:inverse} is, as stated and
proved, an analogue of \file{p0:thm:staircase}(2).
\item \emph{Not shown to be an instance.} Proposition~\ref{xpc:dv:prop:shortcut}
compares the staircases of $p$ and $a$ directly through the injection
behind~\eqref{xpc:dv:eq:sandwich}; it involves no smooth model, and the write
knows no theorem of the volume of which it is an instance.
\end{enumerate}
No novelty is claimed for any of these inversions.
\end{remark}

\section{A rational certificate at a finite index}\label{xpc:dv:sec:certificate}
Theorem~\ref{xpc:dv:thm:enclosure} becomes a rational certificate for rational $r$ and $K\in\{2,3,4\}$. For these $K$, the nonzero gaps
$1-\cos(2\pi h/k)$ are rational:
\[
 k=2:\ (2),\qquad k=3:\ (3/2,3/2),\qquad
 k=4:\ (1,2,1).
\]
More generally, rational lower bounds for the gaps would suffice. Here no such approximation is needed.

Choose integers $J\ge1$ and $b\ge1$, set $S=10^b$, and define integer upward roundings
\[
 V_0=S,\qquad V_j=\ceil{\frac{V_{j-1}}{1-r^j}}\quad(1\le j\le J).
\]
Then $V_J/S\ge\prod_{j=1}^{J}(1-r^j)^{-1}$. The omitted tail satisfies
\begin{align*}
 \log\prod_{j>J}(1-r^j)^{-1}
 &=\sum_{j>J}\sum_{h\ge1}\frac{r^{jh}}h\\
 &\le\sum_{j>J}\frac{r^j}{1-r^j}
 \le\delta_J:=\frac{r^{J+1}}{(1-r)(1-r^{J+1})}.
\end{align*}
If $\delta_J<1$, then $\e^{\delta_J}\le(1-\delta_J)^{-1}$, whence
\begin{equation}\label{xpc:dv:eq:Pcertificate}
 P(r)\le\overline P(r):=
 \frac1S\ceil{\frac{V_J}{1-\delta_J}}.
\end{equation}
For rational $x\ge0$ and an integer $H\ge0$,
\begin{equation}\label{xpc:dv:eq:exp-certificate}
 \e^{-x}\le\left(\sum_{h=0}^{H}\frac{x^h}{h!}\right)^{-1}.
\end{equation}
Every operation in \eqref{xpc:dv:eq:Pcertificate} and \eqref{xpc:dv:eq:exp-certificate} is exact rational arithmetic or an integer ceiling. Substituting these upper bounds for the positive factors in \eqref{xpc:dv:eq:Efinite} yields a rational upper bound $\overline{\mathcal E}$. Set $E=\ceil{\overline{\mathcal E}}$. Because $D(n)$ is an integer, the certified enclosure is
\begin{equation}\label{xpc:dv:eq:integer-certificate}
 \ceil{s_K(n)-E}\le D(n)\le\floor{s_K(n)+u_K(n)+E}.
\end{equation}
This proves the certificate procedure, independently of any numerical exponential evaluation.

For $n=1000$, $K=3$, $r=24/25$, the supplied program uses $b=100$, chooses the first $J$ with $\delta_J<10^{-30}$, and uses $H=140$. It obtains $J=1771$ and the following exact values:
\begin{align*}
 s_3(1000)&=\frac{1420175722807416251109237414763}{3},\\
 u_3(1000)&=7689741909994010189877860653,\\
 E&=666467645896325327332069792.
\end{align*}
Therefore
\begin{equation}\label{xpc:dv:eq:worked-bound}
 \begin{split}
 472725439956575758375747068463&\le D(1000)\\
 &\le481748117158362419220289068699.
 \end{split}
\end{equation}
Independent exact counting gives
\[
 D(1000)=479221051806021018167049012329.
\]
The width reflects the elementary Fourier bound; it is not an estimate of the unknown asymptotic remainder. The certification concerns \eqref{xpc:dv:eq:worked-bound} only, not effective constants in Theorem~\ref{xpc:dv:thm:inverse}.


\begin{corollary}[Certified threshold decision]\label{xpc:dv:cor:decision}
Let $y>1$, $m=\nu_p(y)$, and choose $K\ge2$ with $m>K(K+1)/2$. Suppose $E$ is any rigorously evaluated upper bound for $\mathcal E_K(m,r)$. Then
\begin{align*}
 p(m)-s_K(m)-u_K(m)-E\ge y&\quad\Longrightarrow\quad\nu_a(y)=m,\\
 p(m)-s_K(m)+E<y&\quad\Longrightarrow\quad\nu_a(y)=m+1.
\end{align*}
If neither comparison settles the question, evaluate $a(m)$ exactly and test $a(m)\ge y$. For a rational target and rational enclosures, every comparison can be performed with exact rational arithmetic.
\end{corollary}
\begin{proof}
Subtract the two sides of \eqref{xpc:dv:eq:enclosure} from $p(m)$ to obtain a lower and an upper bound for $a(m)$. Apply Proposition~\ref{xpc:dv:prop:shortcut}. Equality belongs to the first case because the threshold convention is weak; the second case requires a strict inequality.
\end{proof}
This is an effective decision procedure once $m$ is known. It does not require effective constants in the asymptotic inverse theorem.

\section{Reproducible exact computations}\label{xpc:dv:sec:code}
The accompanying archive contains the article source, PDF, standard-library Python implementation, deterministic generated tables, and a manifest. Running the reproduction script computes all outputs anew. Saved tables or receipts are never used as arithmetic inputs. Integer and rational checks use explicit exceptions, so optimization with \texttt{python -O} cannot disable them. The program validates mathematical inputs before computation; in particular booleans are not accepted as integers.

\subsection{Two counting constructions}
The first exact algorithm implements \eqref{xpc:dv:eq:exact-a}. For each minimum $k$, it incrementally builds the bounded-part product $\prod_{j=k}^{M}(1-q^j)^{-1}$ by the usual unbounded coin-change recurrence. At maxima $M$ divisible by $k$, it adds the shifted coefficients for the two removed endpoints. It also adds the singleton terms separately. The unrestricted $p(n)$ table uses an independent standard partition recurrence.

The second algorithm uses conjugation. If $\mu$ is conjugate to $\lambda$, then
\[
 \max\lambda=\#\text{parts of }\mu,\qquad
 \min\lambda=\text{multiplicity of the largest part of }\mu.
\]
Before adding a possible largest part $M$, a two-dimensional table records partitions into parts $<M$ by total size $s$ and length $\ell$. Adding exactly $k$ copies of $M$ contributes size $kM$ and length $k+\ell$; it is admissible exactly when $\ell\equiv0\pmod k$. Summing this condition counts $a(n)$ independently of the minimum--maximum product algorithm. Direct enumeration at small sizes gives a third check. These finite checks validate implementations and examples; the theorems rest on the analytic proofs above.

\subsection{Independent coefficient checks}
The radial coefficients are computed from the truncated exponential-product series in \eqref{xpc:dv:eq:Gtail}. An independent route first constructs the finite $q$-polynomial $S_{L+1}(q)=\sum_jh_jq^j$ and uses the exact moment formula
\[
 g_\ell=\frac1{\ell!}\sum_jh_j(-j)^\ell\qquad(0\le\ell\le L).
\]
The terminating derivative polynomials are checked both by the factorial expression \eqref{xpc:dv:eq:Q} and the differential recurrence \eqref{xpc:dv:eq:Qrec}. The algorithm \eqref{xpc:dv:eq:c-finite} then generates correction coefficients as rational polynomials in $\A^{-1}$. Triangular formal reversion is checked by substituting its coefficients back into \eqref{xpc:dv:eq:inverse-formal} with exact rational arithmetic. These residual tests supplement the symbolic proof of the general triangular construction; a finite collection of substitutions alone is not a proof of a symbolic identity for every parameter.

The Fourier identity is checked in the exact group algebra $\mathbb Q[z]/(z^m-1)$, so no approximate complex roots occur. The rational certificate uses only the integer/rational operations proved in Section~\ref{xpc:dv:sec:certificate}. Separate rational endpoint checks certify the elementary $\pi$ and $\e$ bounds, the $n\ge90$ onset, and the cutoff inequalities in Section~\ref{xpc:dv:sec:quantitative}; the analytic monotonicity proof extends the endpoint checks throughout their stated ranges. The archive documents check ranges and commands, and includes a deterministic receipt that records their outcomes.

\subsection{Scope of reproducibility}
The default run computes $p,a,D$ through $1000$, with the more expensive independent conjugate comparison through $220$. A larger exact count can be requested explicitly. The archive rebuild command regenerates the data, compiles the PDF until consecutive passes have identical bytes, and creates a new archive in a fresh output directory without overwriting its input bundle. A normal run and an optimized run should produce identical semantic receipts and tables. The build uses a fixed timestamp and strips variable PDF metadata to make the packaged result predictable.

No floating-point diagnostic is used to certify a theorem, a remainder constant, a finite-$y$ inverse, or the rational enclosure. The source files contain no private working notes, review reports, or copied source papers. The bibliography supplies public primary-source links instead.

\section{Limitations and conclusions}\label{xpc:dv:sec:limitations}
The leading equivalent for $a(n)$ follows immediately from the sandwich \eqref{xpc:dv:eq:sandwich}. The refinement is that the missing partitions have relative mass $\pi/(2\sqrt{6n})$ to first order, with an exact algorithm for every fixed correction order. The residue obstruction becomes exponentially small at each fixed minimum, and the remaining minimum tail can be made smaller than any prescribed algebraic order by fixing a sufficiently deep truncation.

The quantifiers are essential. The series are asymptotic formal expansions, with no convergence assertion. The fixed-order expansion is not asserted uniform in increasing minimum, modulus, or order. The separate stretched-exponential comparison uses a growing auxiliary cutoff only through exact uniform inequalities. It is not an actual-error asymptotic and does not prove an $\exp(-c\sqrt n)$ rate. There is no global monotonicity assertion for $D$. The smooth inverse bracket has ineffective constants and onset, and does not license blind rounding at all targets. The exact two-index shortcut for $a$ and the rational finite-$n$ certificate are separate effective statements.

The source comparison establishes close prior frameworks and clearly identifies two access limitations. It does not establish first priority for the specialization, coefficients, or combination of results. The outcome is a fully specified application with proofs, exact coefficient generation, and reproducible arithmetic checks.

\begin{samepage}
Several further questions are natural:
\begin{itemize}
\item Make the inverse error constants and their onset explicit, beyond the separate finite-$n$ certification method.
\item Sharpen local twist bounds and fixed-minimum rates using the finer established twisted-product asymptotics, with explicit treatment of their angular scope.
\item Study other minimum weights, first comparing each proposed specialization with existing arbitrary-weight sum-of-tails identities and checking its tail control.
\end{itemize}
\end{samepage}
\begin{writenote}[7 October 2026]
Under Vladimir's standing rule of 4~October 2026, the write found no claim of
this Part that is unproved as stated and none that is wrong, so nothing was
moved to these questions or refuted. Remark~\ref{xpc:dv:rem:transseries}(a)
adds a proved fact (strict increase of $a$ for every $n$), not an answer to
any of the questions above.
\end{writenote}

% [write] Part III's appendix keeps its letter.
\setcounter{section}{0}\renewcommand{\thesection}{\Alph{section}}\renewcommand{\theHsection}{xpcdvapp.\Alph{section}}
\section{Tables regenerated by the companion program}\label{xpc:dv:app:tables}
The tables in this appendix are included from files written by the exact reproduction program during the build. The printed expressions in the main text are also tested by that program.

\subsection{Exact sample counts}
\begin{center}\small
% [write] inlined from the delivered generated/exact_counts.tex (shipped as data/)
% Generated by code/reproduce.py; do not edit.
\begin{tabular}{r r r r}
\hline
$n$ & $p(n)$ & $a(n)$ & $D(n)$ \\
\hline
10 & 42 & 37 & 5 \\
20 & 627 & 556 & 71 \\
50 & 204226 & 187251 & 16975 \\
100 & 190569292 & 179039597 & 11529695 \\
200 & 3972999029388 & 3800030996209 & 172968033179 \\
\hline
\end{tabular}

\par\medskip

% Generated by code/reproduce.py; do not edit.
\begin{tabular}{l r}
\hline
Quantity & Exact value \\
\hline
$p(500)$ & 2300165032574323995027 \\
$a(500)$ & 2235865358882256076332 \\
$D(500)$ & 64299673692067918695 \\
$p(1000)$ & 24061467864032622473692149727991 \\
$a(1000)$ & 23582246812226601455525100715662 \\
$D(1000)$ & 479221051806021018167049012329 \\
\hline
\end{tabular}
\end{center}

\subsection{Radial coefficients}
\begin{center}\small
\renewcommand{\arraystretch}{1.85}
% [write] inlined from the delivered generated/radial_coefficients.tex (shipped as data/)
% Generated by code/reproduce.py; do not edit.
\begin{tabular}{r l r l}
\hline
$\ell$ & $g_\ell$ & $\ell$ & $g_\ell$ \\
\hline
1 & $\frac{1}{2}$ & 7 & $\frac{21841}{10080}$ \\
2 & $\frac{1}{12}$ & 8 & $\frac{10023809}{1209600}$ \\
3 & $\frac{1}{12}$ & 9 & $\frac{26596837}{725760}$ \\
4 & $\frac{89}{720}$ & 10 & $\frac{44079183131}{239500800}$ \\
5 & $\frac{61}{240}$ & 11 & $\frac{82627721881}{79833600}$ \\
6 & $\frac{20287}{30240}$ & 12 & $\frac{8426813263570739}{1307674368000}$ \\
\hline
\end{tabular}
\end{center}

\subsection{Finite enclosure and check receipt}
\begin{center}\small
% [write] inlined from the delivered generated/finite_enclosure.tex (shipped as data/)
% Generated by code/reproduce.py; do not edit.
\begin{tabular}{l r}
\hline
Quantity & Exact value \\
\hline
Lower bound & 472725439956575758375747068463 \\
$D(1000)$ & 479221051806021018167049012329 \\
Upper bound & 481748117158362419220289068699 \\
Product factors & 1771 \\
Exponential degree $H$ & 140 \\
\hline
\end{tabular}
\end{center}
\begin{center}\small
% [write] inlined from the delivered generated/verification_summary.tex (shipped as data/)
% Generated by code/reproduce.py; do not edit.
\begin{tabular}{l l}
\hline
Exact check & Range or count \\
\hline
Euler recurrence versus partition knapsack & $0\le n\le 1000$ \\
Original extrema versus conjugate statistic & $0\le n\le220$ \\
Direct partition enumeration & $1\le n\le35$ \\
Fourier reduction in the integral group ring & 49 cases, $0\le n\le70$ \\
Radial product versus finite moments & degree 20 \\
Derivative recurrence versus factorial formula & $Q_0,\ldots,Q_{20}$ \\
Independent forward-coefficient generators & $c_0,\ldots,c_{12}$ \\
Rational inverse specializations, both statistics & 40 cases, degree 12 \\
Exact threshold breakpoint tests & 864 \\
Certified inverse endpoint tests & 8 \\
Uniform error bound, $n\ge90$ & 19 rational endpoint certificates \\
Uniform certificate invalid-input probes & 18 \\
Rejected invalid-input probes & 22 \\
\hline
\end{tabular}
\end{center}


\begin{writenote}[The merged bibliography, 7 October 2026]
The three delivered bibliographies are merged, and every detail of the
delivered entries is kept; an entry cited by more than one Part is printed
once, with each Part's own note on what it inspected. Bracketed Part numbers
mark entries cited by one Part only. Renamed keys: Part~II's \file{oeis}
and \file{dlmf} are \file{oeisA118096} and \file{dlmfBessel}; Part~III's
\file{oeis} is \file{oeisA117086}, and its \file{nr} is Part~I's
\file{ngorhoades}.
\end{writenote}
\begingroup\raggedright
\begin{thebibliography}{99}
\bibitem{oeis350879} OEIS Foundation Inc., \emph{The On-Line Encyclopedia of Integer Sequences}, A350879, triangle of partitions with $k$ times greatest part equal to length; includes the 17 October 2024 equivalent attributed to V. Kot\v{e}\v{s}ovec. \url{https://oeis.org/A350879}. Checked 4 October 2026. [Part~I]
\bibitem{oeis237753} OEIS Foundation Inc., A237753, partitions with twice the greatest part equal to the number of parts. \url{https://oeis.org/A237753}. Source-listed definition and terms checked 4 October 2026. [Part~I]
\bibitem{oeis237751} OEIS Foundation Inc., A237751, partitions with twice the greatest part less than the number of parts. \url{https://oeis.org/A237751}. Source-listed definition and terms checked 4 October 2026. [Part~I]
\bibitem{oeisA118096} OEIS Foundation Inc., \emph{A118096, Number of partitions of $n$ in which the greatest part is twice the smallest part}, record consulted 4 October 2026. \url{https://oeis.org/A118096}. [Part~II]
\bibitem{oeisA117086} OEIS Foundation, \emph{A117086: Number of partitions of $n$ such that the largest part is divisible by the smallest part}. Entry credited to V. Jovovic, 17 April 2006. Retrieved 4 October 2026; live revision history not inspected. \url{https://oeis.org/A117086}. [Part~III]
\bibitem{szekeres} G. Szekeres, \emph{Asymptotic distribution of partitions by number and size of parts}, in \emph{Number Theory}, Colloquia Mathematica Societatis J\'{a}nos Bolyai 51, proceedings of Budapest 1987, North-Holland, 1990, pp. 527--538. Digitized original volume: \url{https://www.mek.oszk.hu/24100/24107/pdf/24107_1.pdf}. [Part~I]
\bibitem{jiangwang} T. Jiang and K. Wang, \emph{A generalized Hardy--Ramanujan formula for the number of restricted integer partitions}, Journal of Number Theory \textbf{201} (2019), 322--353. \href{https://doi.org/10.1016/j.jnt.2019.02.006}{doi:10.1016/j.jnt.2019.02.006}. \url{https://arxiv.org/html/1805.06108}. [Part~I]
\bibitem{mpp} S. Melczer, G. Panova, and R. Pemantle, \emph{Counting partitions inside a rectangle}, SIAM Journal on Discrete Mathematics \textbf{34} (2020), 2388--2410. \href{https://doi.org/10.1137/20M1315828}{doi:10.1137/20M1315828}. \url{https://arxiv.org/html/1805.08375}. [Part~I]
\bibitem{richmond} L. B. Richmond, \emph{A George Szekeres Formula for Restricted Partitions}, arXiv:1803.08548 (2018). \url{https://arxiv.org/abs/1803.08548}. [Part~I]
\bibitem{hwang} H.-K. Hwang, \emph{Distribution of integer partitions with large number of summands}, Acta Arithmetica \textbf{78} (1997), 351--365. \url{https://matwbn.icm.edu.pl/ksiazki/aa/aa78/aa7844.pdf}. [Part~I]
\bibitem{ngorhoades} H. T. Ngo and R. C. Rhoades, \emph{Integer partitions, probabilities and quantum modular forms}, Research in the Mathematical Sciences \textbf{4}, article 17 (2017). \href{https://doi.org/10.1186/s40687-017-0102-4}{doi:10.1186/s40687-017-0102-4}. [Parts~I and~III; Part~III's key \file{nr}]
\bibitem{johansson} F. Johansson, \emph{Efficient implementation of the Hardy--Ramanujan--Rademacher formula}, LMS Journal of Computation and Mathematics \textbf{15} (2012), 341--359. \href{https://doi.org/10.1112/S1461157012001088}{doi:10.1112/S1461157012001088}. \url{https://arxiv.org/html/1205.5991}. [Parts~I and~III]
\bibitem{rademacher} H. Rademacher, \emph{A convergent series for the partition function $p(n)$}, Proceedings of the National Academy of Sciences \textbf{23} (1937), 78--84. \href{https://doi.org/10.1073/pnas.23.2.78}{doi:10.1073/pnas.23.2.78}. Bibliographic citation to the original; the normalization used here is read from Johansson~\cite{johansson}. [Part~I]
\bibitem{dlmf} NIST Digital Library of Mathematical Functions, Section 23.18, modular transformations, equations 23.18.5--23.18.7. \url{https://dlmf.nist.gov/23.18#E5}. Checked 4 October 2026. (Part~III cites the same section at \url{https://dlmf.nist.gov/23.18}.) [Parts~I and~III]
\bibitem{dlmfBessel} NIST Digital Library of Mathematical Functions, Sections 10.40.1 and 10.17.1, modified Bessel asymptotics and coefficient product. \url{https://dlmf.nist.gov/10.40#E1}; \url{https://dlmf.nist.gov/10.17#E1}. [Part~II]
\bibitem{dyson} F. J. Dyson, \emph{A new symmetry of partitions}, Journal of Combinatorial Theory \textbf{7} (1969), 56--61. Primary paper: \url{https://www.mi.uni-koeln.de/NumberTheory/teaching/Seminare/Partitionen/pages/vortragS12.pdf}. [Part~I]
\bibitem{berkovichgarvan} A. Berkovich and F. G. Garvan, \emph{Some observations on Dyson's new symmetries of partitions}, author's version, especially equations (6.1) and (6.2). \url{https://qseries.org/fgarvan/papers/dyson.pdf}. [Part~I]
\bibitem{mcintosh} R. J. McIntosh, \emph{Asymptotic transformations of $q$-series}, Canadian Journal of Mathematics \textbf{50} (1998), 412--425. \href{https://doi.org/10.4153/CJM-1998-022-x}{doi:10.4153/CJM-1998-022-x}. [Part~II]
\bibitem{zhou} N. H. Zhou, \emph{Asymptotic formulae for Eulerian series}, arXiv:1709.08550v1 (2017). \url{https://arxiv.org/abs/1709.08550}. [Part~II]
\bibitem{pittel} B. Pittel, \emph{Limit shape of a random integer partition with a bounded max-to-min ratio of parts sizes}, Journal of Combinatorial Theory, Series A \textbf{114} (2007), 1238--1253. \href{https://doi.org/10.1016/j.jcta.2007.01.006}{doi:10.1016/j.jcta.2007.01.006}. Publisher abstract and metadata inspected; full theorem text not inspected in this comparison. [Part~II]
\bibitem{ady} G. E. Andrews, A. Dixit, and A. J. Yee, \emph{Partitions associated with the Ramanujan/Watson mock theta functions $\omega(q)$, $\nu(q)$ and $\phi(q)$}, Research in Number Theory \textbf{1} (2015), article 19. \href{https://doi.org/10.1007/s40993-015-0020-8}{doi:10.1007/s40993-015-0020-8}; \url{https://arxiv.org/abs/1502.06860}. [Part~II]
\bibitem{proveit} V. Reshetnikov, \emph{ProveIt}, repository of mathematical reports. \url{https://github.com/VladimirReshetnikov/ProveIt}. Cited for context on previously used inverse methods, not as a comprehensive duplication check. [Part~II]
\bibitem{bfg}
W. Bridges, J. Franke, and T. Garnowski,
\emph{Asymptotics for the twisted eta-product and applications to sign changes in partitions}, Research in the Mathematical Sciences \textbf{9} (2022), article 61.
\url{https://doi.org/10.1007/s40687-022-00355-x}. [Part~III]
\bibitem{ccm}
G. Cesana, W. Craig, and J. Males,
\emph{Asymptotic equidistribution for partition statistics and topological invariants}, Journal of Number Theory \textbf{256} (2024), 373--396.
\url{https://doi.org/10.1016/j.jnt.2023.10.007}.
Primary preprint version 3: \url{https://arxiv.org/html/2111.13766v3}. [Part~III]
\bibitem{af}
G. E. Andrews and P. Freitas,
\emph{Extension of Abel's Lemma with q-Series Implications}, Ramanujan Journal \textbf{10} (2005), 137--152.
\url{https://doi.org/10.1007/s11139-005-4844-z}.
Original full text not inspected; the used general identity was checked through Gupta's reproduction. [Part~III]
\bibitem{gupta}
R. Gupta,
\emph{On sum-of-tails identities}, Journal of Combinatorial Theory, Series A \textbf{184} (2021), 105521.
\url{https://doi.org/10.1016/j.jcta.2021.105521}.
Andrews--Freitas identity: equation (1.7) in publication, equation (1.9) in primary preprint version 1,
\url{https://arxiv.org/html/2002.00447v1}. [Part~III]
\end{thebibliography}
\endgroup
\end{document}
